\chapter{灵界}

内域道争的硝烟，已然散去七年。

七年光阴，足以让许多事物焕然一新。

云霞山半山腰的云隐别院，早已不再是当初那座仅供数人清修的僻静院落。随着叶牧谦声望日隆，以及问道途逐渐成为内域显学，这里在事实上已然取代了昔日天枢城的地位，成为了内域诸宗门联盟共治的行政中枢。

山道被拓宽，以青石板铺就，沿途楼阁依山而建，错落有致。越来越多的修士被吸引而来，或为聆听大道，或为寻求机遇，或单纯向往此地的清静与秩序。他们在云霞山脚及周边平原聚居、贸易、交流，一座崭新的城镇——云霞镇——已初具规模，屋舍俨然，人流如织，虽不及昔日七大城宏伟，却充满蓬勃朝气。

巡天司的总部便设在镇中心，由白言冰、吴东浩统筹，陈默、吴刚等一批转修问道途或理念相合的修士担任骨干。七年来，巡天司法规严明，处事公允，有效调解纷争，清剿残余邪祟，宵小之辈闻风丧胆。在相对平和的环境下，市井之间颇见安居乐业之象。

问道途的传播也比预想中更为顺利。尤其是在资源相对匮乏的低阶修士和散修群体中，这条不依赖外界灵气浓度、强调内求己身的道路，展现出独特的吸引力。

随着陈千羽和脱离宗门政务的王丹两位丹道大师（至少在内域看来，确实称得上一句“大师”）加入转修问道途的相关研究，事情很快就出现了转机：玄脉境界及以上修士，在特定丹药的作用下，可以暂时保留灵枢和玄脉，开辟气海，之后再将早已空置的灵枢和玄脉逐步崩解吸收。如此过程要比先前安全得多了，但主动崩解灵枢与玄脉的痛苦却似乎完全不能抑制。

也有的人试图偷懒，试图长期保留灵枢和玄脉，结果就是他辛辛苦苦温养出来的元气全都沿着空置的灵枢与玄脉泄了个干净，最终也只能被迫崩解灵枢与玄脉。

然而，尽管转修过程，尤其是崩解灵枢与玄脉的过程依旧伴随痛苦，但成功先例的增多和巡天司营造的稳定环境，让越来越多的人愿意尝试。云霞镇及周边，专修或转修问道途的修士数量稳步增长，虽未超越根深蒂固的灵枢径，却已成为内域不可忽视的一股显学力量。

云隐别院内，叶牧谦的三位亲传弟子——白言冰、陈千羽、沈瑶，历经七载磨砺，皆已踏入通明大成之境，真元凝练，各有所长。白言冰剑意愈发澄明刚正，陈千羽对生机平衡的掌控出神入化，沈瑶的“千幻流光”之术更是诡变莫测，三人已成为内域新生代的翘楚，也是巡天司震慑四方的底气之一。

而叶牧谦本人，在享受了数月悠闲日子、作为联盟的精神领袖参知政事，并顺手解决了徒弟家父母养老带来的扩建问题后，近来又将大部分心神投入了对炼器、阵法的研究，以及对下一境界“凝一”的推演之中。

至于为什么叫凝一……

通明之境，见微知著，洞察万物微妙之“异”。然而，若不能从这无穷的“异”中，找到那贯穿始终的“同”，又如何能触摸到那缥缈的“道”的根本？他的道路既名“问道”，那么真正的目标，就应当是那“近乎于道”的层次吧？只是这“凝一”关乎根本道途的整合与升华，远比之前任何一境的突破都要艰难晦涩。

他时而独坐静室，对着一炉香烟冥思；时而翻阅从各方搜集来的古老残卷；时而唤来弟子，于庭院中坐而论道。然而，那一层关键的明悟，却始终如同隔着一层沾湿的窗纸，朦朦胧胧，看得见后面的光，却触不到实质。

\vspace{2em}

时值深冬，云霞山银装素裹。

这一日，天色灰蒙蒙的，似有雪意。

骤然间，东北方向的天空，毫无征兆地扭曲了一下。

那感觉极其怪异，仿佛一块平整的琉璃被无形的力量从内部顶起，凸出、旋转，又骤然平复。空间的褶皱一闪而逝，若非一直关注那片空域，几乎会以为是风雪来临前的光影错觉。然而，紧接着，三道璀璨夺目的华光便悍然撕裂了灰暗的云层，如同从九霄之外坠落的流星，拖着长长的、灼热的光尾，裹挟着令人心悸的呼啸声，朝着大地狠狠砸落！

就在它们即将与地面撞击，引发山崩地裂的刹那，光速骤然衰减，变得轻柔如羽。三道身影从消散的光晕中踏步而出，轻飘飘地落地，点尘不惊。

为首者是一名身着锦蓝云纹长袍的男子，约莫三十许相貌，面容俊朗，但眉宇间带着一股毫不掩饰的倨傲，负手而立，眼神扫过四周荒山雪岭，如同君王审视一片刚刚征服，却贫瘠得让他毫无兴趣的疆土。他便是徐峰，修为已达法相境通灵期。

其左侧是一名身材高瘦、眼神锐利如鹰的青年，穿着暗红色劲装，外罩一件华贵裘氅，嘴角噙着一丝玩味又轻蔑的笑，正是周锐，法相境外显期。

右侧则是一位女子，柳眉凤目，容貌娇艳，身段婀娜，着一袭鹅黄色宫装长裙，外披雪白狐裘，显得雍容华贵。她名唤柳媚，亦是法相境外显期。此刻她正微微蹙眉，以袖掩鼻，仿佛空气中有什么难以忍受的气味。

“啧，这破地方，灵气稀薄得像掺了水的劣酒，还带着一股子土腥味。”周锐率先开口，语气满是嫌弃，“陆执事怕不是立功心切想疯了，这种鸟不拉屎的下界，能掘出什么好苗子？平白浪费我等一个月的光阴，真是晦气！”

徐峰冷哼一声：“既来之，则安之。陆洲新官上任三把火，抓了前任渎职怠政的小辫子，然后派我等下来走个过场，给他的考绩添点彩头罢了。随便找两个看得过去的，带回去交差便是。至于这下界之人……”他眼中闪过一丝漠然，“与蝼蚁何异？不必在意。”

柳媚轻轻弹了弹狐裘上并不存在的灰尘，柔声道：“徐师兄说的是。只是这环境着实恼人，还是尽快寻个落脚处，应付完这差事为好。前方似有烟火气，不如去看看？”

三人神识略一扫荡，便锁定了三十里外那座新兴的、人气旺盛的云霞镇。

“走。”徐峰话音未落，三人已化作三道流光，朝着云霞镇方向疾掠而去，姿态从容，却带着一种居高临下的随意，仿佛不是去往一个陌生的城镇，而是踏入自家后花园。

他们确是来自那传说中、资源丰饶大道昌盛的上界——灵界，隶属于灵界中州皇朝麾下，专司巡查下界、接引有潜力的修士前往灵界的衙门“界门司”。

上一任负责这片下域接引事务的执事，因懒政怠政，将本该每十年一次的正规资质考察与接引流程，次次都敷衍了事，甚至常年懒于派人下界，只胡乱上报“下界无资质优秀者可接引”，用“下界补贴”的灵石中饱私囊。此事近年终究败露，原执事被革职查办。新任执事陆洲为彰显手腕、稳固权位，便火速派遣了徐峰这支小队下来，强制要求他们必须在此地下界至少滞留满一月，并且必须带回至少一名“资质上佳者”以充业绩，堵住悠悠众口，证明他陆洲雷厉风行，绝非前任可比。

徐峰三人心中对此行极为不满。在他们看来，这等灵气贫瘠、道统落后的下界根本就是蛮荒之地，此地的修士与凡人与未开化的野兽无异，岂配与他们这些灵界修士相提并论？他们打定主意，随便逛逛，寻个勉强能入眼的带走了事，其余时间便当是来这穷乡僻壤“体验生活”了。

\vspace{2em}

云霞镇，醉仙居。

这是镇上最好的酒楼，三层木质结构，虽无雕梁画栋，但也拾掇得窗明几净。菜肴以本地山野时鲜和低阶灵兽肉为主，价格虽稍贵但对这些菜肴来说相当实惠，颇受往来修士与富裕凡人的欢迎。

正值午时，大堂里坐了近七成的客人，交谈声、碗筷碰撞声、伙计殷勤的吆喝声混杂在一起，透着热闹的暖意。掌柜是个圆脸的中年人，穿着绛色绸衫，笑容可掬地在柜台后拨弄着算盘。

就在这时，门帘被一股无形的力量悄然拂开，没有带进一丝冷风。

徐峰三人缓步走了进来。

刹那间，仿佛有一盆冰水迎头浇下，大堂里喧嚣的声音不由自主地低落、消失。所有食客，无论是粗豪的体修还是谨慎的法修，都不由自主地停下了动作，目光被这三个突兀出现的身影牢牢吸住。那并非因为他们的容貌衣着多么绝世——虽然确实华贵不凡——而是他们身上自然散发出的那种气质，一种与整个云霞镇、乃至与整个内域都格格不入的疏离与居高临下。他们站在这里，不像客人，倒像是误入羊圈的三头猛虎，尽管收敛了爪牙，却依旧让所有羊羔感到源自生命本能的战栗与不安。

掌柜到底是见多识广，强行压下心头莫名的不适，脸上堆起加倍热情的笑容，小跑着迎上前：“哎哟，三位贵客光临，小店蓬荜生辉！快，快里面请，楼上有雅座！”他眼光毒辣，看出这三人绝非寻常，伺候得格外小心。

徐峰眼皮都未抬一下，径直走向临窗一张空着的、显然是为别人预留了的最好位置。掌柜亲自擦拭本已纤尘不染的桌面，奉上香茗，递上制作精美的菜牌。

“把你们这儿最拿手的，最好的，都端上来。”周锐懒洋洋地靠在椅背上，手指随意点了点菜牌，“酒要最好的。”

“是，是，马上就来，保准让三位仙师满意！”掌柜连声应诺，亲自跑去后厨吩咐。

菜肴很快如流水般呈上。清蒸的寒潭银鱼灵气氤氲，红烧的赤睛鹿腩色泽诱人，炭烤的云雀外焦里嫩，更有几样掌柜珍藏的、来自远方的稀罕山珍。酒是埋藏了三十年的“火烧云”，开坛香气四溢。

徐峰三人动筷品尝，姿态优雅，却吃得很快。他们不曾交谈，偶尔周锐会对某道菜撇撇嘴，柳媚则会用清茶漱了漱口。他们的吃相并不粗鲁，但那种对食物的态度，并非品尝珍馐，更像是在完成一项任务，或者检验这些“下界之物”是否够格入他们的口。

不到半个时辰，满桌佳肴美酒便被扫荡一空。徐峰拿起一方洁白的丝帕，轻轻擦了擦嘴角，站起身，便欲离去。

掌柜一直亲自在旁伺候，见状连忙又堆起笑脸，搓着手上前，微微躬身：“三位仙师，可还满意？承惠，一共是八块下品灵石，或者四钱银子，您看……”他话里带着十二分的小心，甚至不敢直接说要结账，只敢暗示。

话音未落。

“聒噪。”徐峰眉头一皱，仿佛听到了一只苍蝇在耳边嗡嗡作响。他甚至没有正眼看掌柜，只是随手一挥，如同拂去一片落叶。

“啪！！”

一声极其清脆响亮的耳光声，震得整个大堂嗡嗡作响。掌柜那圆润的脸上，瞬间出现了一个清晰的、红肿的掌印。他连惨叫都未能发出，整个人就像被无形的巨锤狠狠砸中，离地飞起，横着撞翻了两张厚重的实木方桌，杯盘碗盏稀里哗啦破碎一地，汤汁酒液溅得到处都是。他滚落在地，口鼻之中鲜血汩汩涌出，身体抽搐了两下，便一动不动，昏死过去。

死寂。

醉仙居内，落针可闻。所有食客、伙计都僵在原地，脸上血色瞬间褪尽，惊恐地瞪大眼睛，看着那三个缓缓向门口走去的身影，又看看地上生死不知的掌柜，连呼吸都仿佛停滞了。

周锐甚至懒得再看那狼藉一眼，嘴角那丝玩味又轻蔑的笑意加深了些，仿佛只是做了一件微不足道、且理所当然的小事。

“下界的贱民，也配向我要钱？”

他低声嗤笑，声音不高，却清晰地钻入每个人的耳中，带着刺骨的寒意。

三人就这样扬长而去，穿过死寂的厅堂，掀开门帘，步入外面灰蒙蒙的冬日街道，留下满屋的狼藉、刺鼻的血腥味，以及深入骨髓的冰冷恐惧。

终于有人反应过来：“快，快去报巡天司！赶紧把掌柜的搬到安平医舍去！”

修士们手忙脚乱地把掌柜搬到椅子上，两名身强力壮的体修抬起椅子就跑。陈千羽诊治一番，发现掌柜全身上下多处骨折，但好在暂无性命之虞。剩下的人此时也没有吃饭的心情了，于是各个带着沉重的心情散了。

\vspace{2em}

而这，仅仅只是开始。

在接下来的几天里，这三个天上来的恶客，将云霞镇当作了他们肆意取乐的猎场。他们看上的东西，无论是店铺里陈列的、带有微弱灵气的药材矿石，还是摊贩手中精巧的凡俗玩物，直接伸手便拿，如同取用自家之物。若有谁敢露出半分不满或迟疑，轻则被一道气劲打得筋断骨折，吐血倒飞；重则如同醉仙居掌柜一般，生死难料。

他们的恶行迅速从财物蔓延到人身。镇东头新来的戏班子里，一个嗓音清亮、容貌秀美的花旦被徐峰看中，当夜便从暂居的后台失踪。次日清晨，有人在她常去吊嗓子的溪边林地里，发现了她衣衫破碎、布满淤青与伤痕的尸体，双目圆睁，凝固着无边的恐惧与绝望。

巡天司也确实并非没有反应，最初接到零散的报案，执勤的巡天卫还试图按章程前去询问调查，结果连对方的身都近不了，便被无形的威压震慑得神魂欲裂，吐血败退。

两人前来阻止的人果然蝼蚁一般，随即更加大胆，开始当街调戏良家妇女。一名卖绣品的少女被周锐捏住下巴，吓得花容失色。其兄长是镇上一名低阶灵枢径修士，怒而上前理论，被周锐像拎小鸡般掐住脖子，狞笑着拎出镇外。不久，周锐独自返回，指尖犹带一丝未擦净的血迹，而那青年的身影，再未出现。

柳媚始终冷眼旁观。她虽不屑于徐峰、周锐这般急色与暴虐的行径，觉得有失身份，但也从未出言劝阻。在她根深蒂固的观念里，下界之人，无论修士还是凡人，皆与蝼蚁草芥无异。蝼蚁被踩死，草芥被折断，算得了什么大事？她只是每日用着镇上最好的客栈最好的房间，用着强征来的、带有清香的花露，然后站在窗边，冷漠地俯瞰着这座因他们三人而逐渐被恐慌笼罩的小镇

短短两三日之间，云霞镇已是人心惶惶。很快，六人失踪，生不见人死不见尸；三位良家女子受辱后惨死；数十名凡人百姓或低阶修士因“冒犯”而被打伤，店铺被毁、货物被抢之事更是不计其数。昔日祥和的小镇，如今笼罩在一片愁云惨雾之中，白昼行人稀少，入夜更是家家闭户，灯火黯淡。

很快，巡天司副使吴刚与陈默两人知道了这件事。

两人皆是经历过道争血火、真心信奉并维护眼下这片安宁的修士，闻听竟有如此暴虐狂徒在云霞镇肆意践踏法度、残害百姓，顿时怒发冲冠。他们即刻点齐一队最为精干的巡天卫，根据线报，在镇外通往一处风景山谷的路上，截住了正闲逛着、寻找新“乐子”的徐峰与周锐。

“站住！尔等何方狂徒，竟敢在云霞镇境内屡施暴行，伤人害命！立刻束手就擒，随我回巡天司受审！”

吴刚身材魁梧，声如洪钟，手中制式长刀已然出鞘，转修问道途后凝练的元气在体内鼓荡，蓄势待发。他身后，陈默面色沉凝如铁，同样暗运元气，十几名巡天卫结成简单的阵势，刀剑出鞘，寒光闪闪，怒视着前方两人。

徐峰与周锐停下脚步，缓缓转过身，脸上没有丝毫意外或紧张，反而像是看到了什么极其有趣的事物。

周锐上下打量着吴刚一行人，尤其是他们身上制式的、在他看来粗陋不堪的甲胄和武器，嗤笑出声：“巡天司？呵……下界蝼蚁们自己弄出来过家家的玩意儿？也配拦我的路？”

陈默强压怒火，声音因为极致的愤怒而显得无比冰冷：“无论你们来自何处，有何背景，既踏入内域，行于此地，便需守此地的规矩！伤人害命，天理难容！今日若不给出交代，休想离开！”

“规矩？天理？”

徐峰仿佛听到了世间最荒谬的笑话。他略微偏头，似乎第一次正眼看了看吴刚和陈默。

随即，他那张俊朗的脸上，轻蔑之色如同潮水般蔓延开来，再无丝毫掩饰。

“连法相境的门槛都未摸到的废物，体内那点可怜的真气驳杂不纯，也敢在我面前谈规矩，论天理？”

他摇了摇头，仿佛失去了最后一点耐心。

“也罢。今日便叫你们这些坐井观天的下界蛙虫，亲眼见识一下，何为真正的天高地厚，何为我等与尔等之间不可逾越的鸿沟！”

他甚至没有做出任何动手的架势，只是心念微动，将自身法相境通灵期的威压，如同解开了某种封印，毫无保留地、狂暴地向四周倾泻而出！

“轰——！！”

无声的轰鸣在在场每一个巡天司修士的神魂深处炸响！刹那间，以徐峰为中心，方圆百丈内的天地灵气猛地紊乱、暴动，仿佛平静的海面骤然掀起滔天巨浪。一股浩瀚如海、沉重如山岳的恐怖气息凭空降临，那气息中不仅蕴含着纯粹的力量压迫，更仿佛带着一丝勾连天地灵机本源的玄奥意味，直接作用于生灵的神魂本源。

吴刚与陈默首当其冲，两人只觉得周身奔流的元气瞬间凝滞，像是被冻在了气海之中，呼吸为之断绝，心脏仿佛被一只无形大手狠狠攥住。更为可怕的是神魂层面的冲击，如同万钧重锤狠狠砸在识海，眼前猛地一黑，耳中嗡鸣不休，喉头一甜，鲜血已不受控制地从嘴角溢了出来，身体晃了两晃，几乎站立不稳。

他们身后那些修为稍逊的巡天卫更是惨不忍睹，纷纷脸色煞白如纸，闷哼声中，修为最弱的几人直接双眼翻白，一声不吭地软倒在地，昏迷过去。其余人也东倒西歪，勉力支撑着兵器，才没有当场跪倒。其中有一位巡天卫出身青石城，他只觉得如此威压比当年他经历的沈天穹威压更加强横！

仅仅是一个照面，对方甚至未曾抬手，未曾念咒，仅仅凭借境界的绝对差距带来的威压，便让他们这支巡天司的精锐小队溃不成军，人人带伤！

吴刚与陈默心中掀起了惊涛骇浪，那不仅仅是震惊，更是一种近乎绝望的明悟。

他们终于真切地感受到了对方那深不可测的实力，那是一种超越了内域当前认知层次的力量。眼前之敌，绝非他们所能抗衡，怪不得那些巡天卫全都受了伤。

甚至可能倾尽如今内域现存的所有灵枢径修士之力，都难以与之匹敌！对方看他们的眼神，与看地上挣扎的虫豸，确实毫无分别！

“撤！快撤！！”陈默性格更为冷静坚韧，强忍着神魂撕裂般的剧痛和全身骨骼欲碎的压迫感，猛地一咬舌尖，借着剧痛刺激清醒了一瞬，一把拉住受伤更重、几乎要被那威压按倒在地的吴刚，同时用尽力气掷出早已扣在手中的数枚烟雾符箓。

“噗噗噗！”

浓密的、混杂着刺鼻气味的灰白色烟雾瞬间爆开，笼罩了方圆数十丈，暂时隔断了视线与部分气机感应。

“走！回山！禀告叶先生！”陈默嘶声低吼，与勉强提气的吴刚一起，带着残存的、连滚爬的巡天卫，用尽平生最快的速度，向着云霞山方向亡命遁逃，狼狈不堪。

烟雾缓缓散去，徐峰与周锐依旧站在原地，连衣角都未曾凌乱一分。徐峰看着那些人远遁的背影，嘴角撇了撇，吐出两个字：“无趣。”

周锐舔了舔嘴唇：“师兄，要不要追上去玩玩？”

“不必了，一群杂鱼，杀了都嫌脏手。”徐峰转身，“走吧，听说这镇子东头新来了个戏班子，去看看有没有新鲜货色。”他的语气平淡，却带着一种绝对的掌控感。仿佛云霞山上的存在，也仅仅是这片“蝼蚁之地”里，一只稍微强壮些、或许能让他多费点手脚的大号蝼蚁罢了。

\vspace{2em}

片刻后，别院的正厅内，气氛凝重得如同冻结。叶牧谦端坐主位，旁边是刚刚停止对弈的沈天穹（退休）。下方，吴刚与陈默脸色惨白，身上带着威压冲击造成的内伤，气息不稳，正急促地讲述着镇上的惨状与他们遭遇的惨败。

白言冰、陈千羽、沈瑶被紧急召来，三人静立一旁，面色冷峻如冰，眼中却似有火焰在无声燃烧。

“法相境……还是至少外显期，甚至可能有通灵期……”沈天穹沉吟，他曾经站在这个高度，深知其可怕，“内域现存修士，即便是老夫全盛时期，对付一个外显期也需小心。至于通灵期……我内域灵枢径修士，若无特殊机缘或至宝，绝难正面抗衡。”

当然，更吸引叶牧谦本人的，是这几个人的出身。

灵界。

没错，这个词又一次撞进他的视野中来。

而且这次是真真正正的“灵界”，不是沈天穹等人提及的、虚无缥缈的、只出现在传说和古老记载中的“灵界”。

如果能把他们请上山来，一定要好好盘问一番。

当然，现在的目标是先把他们抓住。

叶牧谦闭目片刻，神识遥遥感知了一下云霞镇方向那几道毫不掩饰、充满恶意与傲慢的强大气息，缓缓睁开眼：“灵界来客……果然不善。沈府主所言不虚，灵枢径修士在此等敌人面前，确实先天受制。”

他的目光扫过自己三位已达通明境大成的弟子，语气转为一种沉着的分析，“不过，我问道途修士，真元源于自身气海，自给自足，对天地灵气的依赖远小于彼辈。彼倚仗外灵，威势虽猛，变化与持久却未必占优。或许……尚有一战之力。”

白言冰三人感受到那远方传来的、令人心悸的压迫感，非但没有畏惧，眼中反而燃起熊熊战意。七年平静修炼，他们的锋芒并未被磨去，只是更加内敛。如今恶客临门，践踏他们守护的一切，正是试剑之时！

“师尊，请下令！”白言冰抱拳，声音坚定。

叶牧谦看着三位战意昂扬的弟子，又看了看面露忧色却同样坚定的沈天穹、吴刚、陈默等人，心中开始计议。这三人行事毫无底线，残忍暴虐，且对下界充满歧视，谈判妥协绝无可能。

外界恶客，自然要以雷霆手段应对。

但是，他可是要抓活的。依然要有万全的计划。

不仅是避免三人逃脱，也要把破坏与恐慌限于一地之内。

“好。”叶牧谦说。随即与几人商议定下具体实施手段如此这般。

\vspace{2em}

云霞镇的清晨，依旧被一股无形的低气压笼罩着。

连续几日的阴霾天气并未好转，铅灰色的云层低低压在镇子上空，仿佛一块脏污的棉絮，吸饱了水汽，却迟迟不肯落下雪来。街道上的行人比往日稀少了许多，即便有事不得不外出，也是步履匆匆，目光警惕地扫视四周，尤其避开镇上那家最气派的“云来阁”客栈。空气中弥漫着一种恐惧与愤怒交织的沉默，连往日清晨最热闹的早点摊子，今日都显得有气无力，灶火不旺，蒸气稀薄。

在这片压抑之中，客栈里却悄然完成了一场无声的替换。原先那个战战兢兢、面有菜色的店小二不见了踪影，取而代之的，是一对相貌出众、气质迥异于常人的年轻夫妇。

男子约莫三十出头，身姿挺拔如松，穿着一身浆洗得有些发白的粗布短打，袖口挽到小臂，露出一截线条流畅、结实有力的手臂。他正拿着抹布，一丝不苟地擦拭着大堂里一张榆木桌子的边角。他的动作看似寻常，甚至有些过于认真，每一个擦拭的轨迹都平稳而连贯，透着一股寻常脚夫或伙计绝难具备的沉静与专注。最引人注目的是他的眉宇，双眉如墨，斜飞入鬓，眼神清亮湛然，即便做着最低微的活计，眉梢眼角也自带一股难以掩藏的凛然正气，仿佛是刻意收敛了锋芒的宝剑，剑鞘朴素，却压不住那股隐而不发的锐意。

柜台后，那布衣荆钗的女子正低头整理着一本陈旧的账簿，实则只是在随手翻动，心思全不在那上面。她乌黑的秀发仅用一根普通的木钗松松绾起，几缕碎发垂在雪白的颈侧，身上是再寻常不过的靛蓝色粗布衣裙，洗得微微发白，却异常洁净。她未施半点脂粉，可那张清丽绝俗的容颜，在清晨略显昏暗的光线下，仿佛自带柔光。尤其是她偶尔抬起眼睑，那双眸子清澈灵动得如同山间最纯净的泉眼，顾盼之间，流转着一种难以言喻的灵韵与风致。

巳时刚过，木质楼梯传来一阵略显沉重的脚步声，伴随着肆无忌惮的谈笑。

徐峰与周锐一前一后，大摇大摆地走下楼来。两人皆穿着灵界款式的锦缎法袍，色泽鲜艳，质地华贵，与这朴素甚至有些破旧的客栈大堂格格不入。徐峰脸上带着宿醉未消的慵懒与惯有的傲慢，周锐则揉着肚子，嘴里嘟囔着“这破地方的饭食寡淡无味，勉强果腹罢了”。

周锐的目光漫不经心地扫过大堂，先是掠过正在擦柱子的白言冰——一个普通的店伙计，换了就换了，引不起他丝毫兴趣。然而，当他的视线落在柜台后那道低眉顺眼的蓝色身影时，就像被磁石牢牢吸住，再也挪不开了。他脚步一顿，眼睛骤然亮起贪婪的光芒，用胳膊肘狠狠捅了一下身前的徐峰，压低了声音，嗓音里满是发现猎物的兴奋：“师兄，快看！昨日还没见过这小娘子，啧啧，这穷乡僻壤竟还藏着如此绝色货色！可比前几天那些庸脂俗粉强出百倍！”

徐峰正憋闷呢，闻言顺着周锐所指望去。只一眼，他脸上那漫不经心的慵懒瞬间被一种更为露骨的轻佻与贪婪取代，眼中闪过毫不掩饰的异彩。他完全无视了近在咫尺、仿佛背景板一般的白言冰，径直走到柜台前，伸出指节敲了敲硬木台面，发出“笃笃”的声响。他身体微微前倾，刻意将声音压低，带着一种自认为迷人、实则充满恶意的暧昧腔调：“小娘子，新来的？爷怎的以前从未见过你？生得这般水灵，待在这破店真是暴殄天物了。”

他嘴角勾起一抹邪笑，“不如跟爷去灵界，那才是真正的仙家福地，繁华盛景，保你见了之后，再也看不上这下界的穷酸模样。跟着爷，保你享不尽的荣华，用不完的资源，如何？”

陈千羽握着账簿的手指微微收紧，指节有些发白。她依着计划，缓缓抬起头，脸上恰到好处地浮现出惊慌与羞愤，清澈的眸子里蒙上一层水雾，身体向后瑟缩了一下，仿佛受惊的小鹿，声音细若蚊蚋，却带着不容侵犯的坚决：“客官……请自重。”

“自重？”徐峰仿佛听到了天大的笑话，嗤笑出声，眼中的欲望更加炽烈，“在这下界，爷看上你，是你八辈子修来的福分，别不识抬举！”说着，他竟直接伸手，食指弯曲，以一种极其轻佻的姿态，朝着陈千羽光滑细腻的下巴挑去。动作熟练，显然不是第一次做这等事。

“住手！”

一声愤怒至极的暴喝，像惊雷般在大堂炸响。声音未落，只听“砰”的一声闷响，一块湿漉漉的抹布被狠狠摔在旁边一张空桌上。

白言冰身影如电，一个箭步便从柱子旁蹿至柜台前，结结实实地挡在了陈千羽身前。他胸膛剧烈起伏，双目喷火般怒视着徐峰，因为激动而脸色涨红，脖颈上青筋微微凸起，将一个眼见妻子受辱、血气上涌的普通丈夫演绎得淋漓尽致。

他伸手指着徐峰，手指因为愤怒而微微颤抖：“光天化日，朗朗乾坤！你们……你们竟敢如此调戏良家妇女！还有没有王法了！”

这怒火，三分是听闻镇上惨剧的愤懑，三分是演技，四分是见到这家伙竟然真在自己面前调戏自己的爱人导致的血气上涌。三者交融，竟逼真得让徐峰一时都未起疑。

徐峰伸出的手僵在半空。这几日在云霞镇乃至整个内域，他所到之处，凡人修士无不战战兢兢，避之唯恐不及，何曾被人如此当面厉声呵斥？尤其呵斥他的，还是一个他眼中蝼蚁不如的下界凡人，理由竟是为了保护另一个蝼蚁。他先是错愕，随即一种被严重冒犯的暴怒如同火山般从心底喷发出来，瞬间淹没了那点因陈千羽容貌而起的旖旎心思。

“王法？哈哈哈哈哈！”徐峰气极反笑，笑声尖厉，眼神却冰冷阴鸷得如同毒蛇，死死钉在白言冰脸上，“在这鸟不拉屎的下界，爷说的话，就是王法！你是个什么东西，也敢在爷面前狂吠？找死！”

他本就因任务不顺而心火旺盛，眼前这个不知死活的凡人丈夫，恰好成了他宣泄怒火的最佳对象。他脑海中瞬间闪过之前处理那些胆敢反抗的下界修士的场景，嘴角咧开一个残忍的弧度，冷笑道：“好，好得很！不是要护着你婆娘吗？有骨气！爷就给你一个机会。”

他抬手指向镇外，“镇子东边那片荒林子，看见没？有种，就带着你的小娘子跟过来！爷今日心情好，亲自指点指点你们，什么叫真正的天高地厚，什么叫有些人，是你们永远也得罪不起的！”

他打算像捏死蚂蚁一样，将这不知天高地厚的夫妇引出镇外无人处，随手灭杀，那个让他心痒难耐的清丽女子，自然就成了他随意处置的战利品，到时候……

白言冰脸上适时地露出剧烈的“挣扎”与“恐惧”，他回头望了一眼“瑟瑟发抖”的陈千羽，又看了看面前气焰嚣张的徐峰，最终像是被逼到了绝境，猛地一咬牙，脸上涌起一股豁出去的悲壮：“去就去！这世道，还没个说理的地方了？我还怕你不成！”

说着，他一把拉住陈千羽微凉的手，用力握了握，传递着只有彼此才懂的信号，然后做出一副“夫妻同心、其利断金”的模样，拉着陈千羽，径直穿过大堂，头也不回地朝着客栈门外、镇东的方向大步走去，背影竟带着几分决绝的意味。

周锐全程抱着胳膊，斜倚在楼梯扶手旁，津津有味地看着这出闹剧。见白言冰夫妇不知死活地真敢应战，他不由得打了个长长的哈欠，脸上写满了无聊与轻蔑。他踱步到徐峰身边，拍了拍徐峰的肩膀，懒洋洋地道：“师兄，两个彻头彻尾的凡人罢了，蝼蚁一样的东西，也值得你亲自跑一趟？动动手指头就能碾死的货色。”他揉了揉肚子，“师弟我还没吃饱呢，昨儿那坛醉仙酿还剩半坛，我再去叫两个小菜，就在这儿，温着酒，静候师兄佳音便是。”

他压低声音嘿嘿笑道：“等你料理了那不知死活的小子，嘿嘿，师兄吃肉，可别忘了给师弟我也留口热汤啊。”

他压根没将这对凡人夫妇放在心上，认定徐峰此去不过是弹指间的事，自己乐得在温暖舒适的客栈里继续享受酒食，等着坐享其成，分享“战利品”。

徐峰鼻腔里发出一声冷哼，算是默认了周锐的安排。他连看都懒得再看这懒蛋一眼，身形一晃，便如一片没有重量的羽毛般飘出了客栈大门，然后不紧不慢地，遥遥缀在白言冰夫妇身后，朝着镇东荒林而去。姿态悠闲，宛如踏青，只是眼中闪烁的，是猫戏老鼠般的残忍快意。

镇东的荒林，是云霞镇百姓平日砍柴都很少深入的偏僻之地。时值寒冬，前几日刚落过一场大雪，林间的积雪大多未曾融化，只在某些向阳的坡地露出些许灰黑色的冻土。无数光秃秃的树木伸展着嶙峋扭曲的枝丫，刺向铅灰色的天空，像一丛丛狰狞的鬼爪。寒风穿过林木，发出低沉的呜咽，卷起地面上的雪沫，更添几分肃杀与荒凉。这里寂静得可怕，连一声鸟叫都听不见，仿佛所有的生灵都预感到了即将到来的血腥，早早逃遁了。

徐峰姿态从容地踏入林中一片较为开阔的空地，正好看见前方那对夫妇停下了脚步，转过身来，面对着他。两人脸上早已没有了之前在客栈里的惊慌与愤怒，只是平静地望着他，那平静之下，是一片深不见底的冰冷。

徐峰脸上残忍的笑意扩大了，他环顾了一下四周的凄清景象，用一种宣布死刑般的口吻悠然道：“选这么个偏僻又清净的地方做你们的葬身之地，倒是有点眼光。也省得脏了镇子的地。”他目光落在白言冰脸上，充满了怜悯与嘲弄，“小子，下辈子投胎，记得长点记性，没那个实力，就别学人家逞英雄，护什么女人。有些东西，不是你该碰，更不是你能护得住的。”

他自觉废话已经足够，对付这种下界蝼蚁，他甚至懒得动用神通法宝。心念一动，法相境通灵期的磅礴威压毫无保留地骤然释放！精神层面的压迫感如同无形的海啸，伴随着他刻意引动的一丝天地灵气躁动，朝着前方两人猛然拍去！

他仿佛已经看到了这对凡人夫妇被这远超他们理解的力量压得肝胆俱裂，瘫软在地，涕泪横流地跪地求饶的丑态——就像之前那个不知好歹的吴刚一样。他甚至准备好了几句更为恶毒的嘲弄，只等对方出丑。

然而，预料中的场面并未出现。

那海啸般的威压及体的瞬间，白言冰与陈千羽甚至连衣角都未曾拂动一下。两人依旧稳稳地站在原地，如同脚下生根，与这片荒凉冻土融为一体。他们脸上的平静没有被打破，反而，之前那仅存的一丝伪装出来的情绪也彻底消失，只剩下一种剔除了所有杂质的、纯粹的冰冷，看向徐峰的眼神，如同在看一个死人。

空地上寒风依旧，卷起的雪沫在三人之间打着旋儿，气氛诡异地安静了一瞬。

徐峰心头猛地一沉，那股不对劲的感觉瞬间放大！“咯噔”一下，如同冰冷的石块投入心湖。

怎么回事？这两个“凡人”怎么可能在他的威压下毫无反应？！难道是身上佩戴了什么罕见的护身符箓？不，不对，就算是符箓，也绝不可能如此轻松地抵御法相境通灵期的神魂威压！

还未等他从惊疑中理清头绪，异变，就在他心神微分的这一刹那，毫无征兆地陡然爆发！

只见那一直垂手而立、看似柔弱无助的布衣女子，忽然抬起了双手。她的手指纤细白皙，此刻却以肉眼难以捕捉的速度飞速变幻，结出了一个玄奥古朴、充满了勃勃生机又暗藏无尽肃杀之意的法印！樱唇轻启，吐出一个清晰而短促的音节：

“生！”

没有惊天动地的巨响，却有一股磅礴到难以想象、精纯到令人心悸的生机领域，以陈千羽的身体为中心，轰然向四面八方展开，瞬间便笼罩了方圆数十丈的每一寸空间！在这天寒地冻、万物凋零的严冬荒林之中，颠覆常理、令人瞠目结舌的一幕发生了：

皑皑白雪覆盖的地面之下，传来密集而轻微的“噼啪”碎裂声；那些因为深冬而干枯着的老树树皮骤然皲裂，冒出了新生的枝桠；甚至连三人周围的空气，都仿佛被无形的生命力所浸润。

随即，无数嫩绿欲滴、却生满了尖锐倒刺的荆棘藤蔓，凭空疯狂地滋长出来！ 它们从雪层下钻出，从枯树干上抽出，从虚无的空气中凝聚成形，速度快得超出了目光捕捉的极限！这些藤蔓如同拥有独立智慧的灵蛇，又像是训练有素的军队，刚一出现，便带着“嘶嘶”的破空之声，从四面八方、上下左右，朝着徐峰电射缠绕而去！

首先是脚踝，冰冷的、滑腻的、带着毛刺的触感瞬间包裹了徐峰的双足，并将他牢牢钉在原地。紧接着，大腿、腰身、手臂……仅仅呼吸之间，徐峰整个人就被这突如其来的绿色狂潮所淹没，无数坚韧无比的藤蔓层层叠叠地缠绕上来，不仅束缚他的行动，藤蔓上那些细密的尖刺更试图扎透他法袍自带的微光，一股奇异的封禁之力混合着精纯的灵机，顺着藤蔓渗透，不断干扰、阻滞着他体内灵力的自然运转。

“什么？！！”徐峰惊恐交加地发出一声怪叫，瞳孔骤然缩成了针尖大小！

这哪里是什么凡人手段？！这分明是造诣极高、甚至带有某种道韵的木属神通！而且这神通发动之迅捷，领域展开之磅礴，藤蔓操控之精妙，绝非普通下界修士所能拥有！

巨大的震惊之后是暴怒。

“蝼蚁安敢欺我！”

他怒吼一声，被戏耍的羞辱感和潜在的危机感让他彻底爆发。体内澎湃的法力如同决堤洪流，轰然运转！炽热而耀眼的赤红色灵光从他周身每一个毛孔迸发出来，仿佛一尊瞬间点燃的火炉！缠绕在他身上的藤蔓被这高温灼烧，立刻发出嗤嗤的刺耳声响，青烟冒起，最内层的藤蔓以肉眼可见的速度变黑、碳化、断裂。

徐峰心中一喜，正要鼓荡法力将这些烦人的东西彻底震成齑粉，脸上的表情却猛地僵住。

那些被烧断的藤蔓断口处，非但没有萎缩枯萎，反而以更快的速度蠕动、再生！断裂的绿色躯体仿佛拥有无穷无尽的生命力，几乎是刹那间就重新抽出新的嫩芽，并以疯狂的速度蔓延、生长、缠绕！更要命的是，周围雪地、枯树、空气中，更多的藤蔓还在源源不断地滋生出来，前赴后继，无穷无尽！他刚震碎、烧却一层，立刻就有两层、三层缠绕上来，而且新生的藤蔓似乎对火焰高温有了一定的适应性，变得愈发坚韧难缠。这根本不是普通的木系束缚法术，而是蕴含了某种“生生不息”法则真意的神通领域！

就在徐峰被这诡异难缠、越挣扎越紧的藤蔓弄得手忙脚乱，注意力被完全吸引的瞬间——

那一直沉默站立、仿佛只是旁观者的布衣男子，动了。

没有掐诀念咒，没有蓄势聚气，甚至没有带起一丝劲风。他只是平静地抬起右手，并拢食指与中指，捏成剑诀，向着前方被藤蔓包裹的徐峰，极其随意地、虚虚一划。

“铮——！”

一道清越冰冷的剑鸣，仿佛自九幽深处响起，又似从灵魂层面直接迸发，毫无阻碍地响彻了整个死寂的荒林，甚至压过了寒风呼啸与藤蔓摩擦的声响！

白言冰手中空空如也，并无剑器。但就在他指尖划动的轨迹上，一道凝练到极致、纯粹到令人心悸的寒冰剑意凭空生成，正是白言冰以真元凝成的剑意！这道剑意随着他手指的动作，如同被无形的力量骤然击碎的万载玄冰，骤然散开！

“咻咻咻咻——！！”

每一道剑意都细碎如晶莹的冰棱，边缘却锋利得足以切割空间；每一道剑意都闪烁着清冽刺骨的寒光，轨迹飘忽，毫无规律可循。它们汇聚在一起，如同被极地暴风卷起的、冰冷致命的暴风雪霰，呼啸着，旋转着，笼罩了徐峰周身前后左右、上下内外所有可能闪避的空间，铺天盖地地攒射而去！

修习至巅峰的“寒光散乱”！

其核心依然是防御，但是是那种专门寻找、渗透、瓦解的“进攻式”防御！

徐峰浑身的汗毛在这一刻全部倒竖起来，一股前所未有的致命寒意，从尾椎骨直冲天灵盖！

他从未见过如此古怪又如此凌厉的剑意！

它散乱，却仿佛每一道细碎寒光都长着眼睛，精准地指向他防御的薄弱处；它冰冷，那寒意并非仅仅作用于肉体，更带着一种冻结魂魄、侵蚀灵力的特性；最让他感到骇然的是，这剑意中蕴含的力量本质，似乎与他惯常依赖、操控的天地灵气截然不同。那是一种更加内敛、更加坚韧、仿佛源于生命自身的力量，对他以灵枢径法力构筑的防御层，竟然有着一种奇特的、如同沸汤泼雪般的穿透与消融效果！

“给我破开！！”徐峰发出野兽般的嘶吼，惊怒交加。他再也顾不得节省法力或保持风度，双手在身前舞出一片赤红色的残影，体内法力疯狂涌出，在身前刹那间布下了足足七重灼热耀眼、符文流转的火焰灵盾，每一重灵盾都足以轻易抵挡真符境及以下修士的全力一击。

然而，那暴风雪般袭来的散乱寒光，竟似拥有某种灵性，或者说，驱使它们的人，拥有着洞穿一切虚妄的眼睛。赤红的火焰灵盾光芒流转，灵力节点与防御强度在洞察之下并非铁板一块。

“嗤嗤嗤嗤——！”

令人牙酸的密集穿透声响起，那千百道寒光剑气总能匪夷所思地找到灵盾能量流转间那细微到几乎可以忽略不计的薄弱之处，或者灵力衔接的缝隙，然后无孔不入地钻透进去！

第一重灵盾，破！第二重，破！第三重……眨眼之间，七重火焰灵盾如同被戳破的泡沫，接连碎裂，化为漫天逸散的火星！

尽管大部分寒光剑气也在穿透过程中被灼热的火灵力消耗湮灭，但仍有数十道漏网之鱼，如同冰锥般狠狠攒射在徐峰身上！他身上的锦缎法袍发出不堪重负的撕裂声，留下了数十道深深的、边缘带着白霜的裂口！更有几道剑气穿透法袍的灵光防御，直接切割在他的皮肉之上！

“呃啊——！”徐峰闷哼一声，伤口处传来的不仅仅是切割的痛楚——这种痛楚他经历过多次，没被刀剑砍过的修士不算好修士——还有一种刺骨的、仿佛连灵魂都要冻僵的寒意。

更重要的是，那剑气中蕴含的奇异真元如同附骨之疽，钻入他的经脉，开始侵蚀、阻滞他自身法力的运转！剧痛与寒意让他瞬间清醒，一个让他如坠冰窟的念头无比清晰地浮现出来：

踢到铁板了！这对狗男女，根本不是什么任人宰割的凡人！他们是修为极高、功法诡异、配合默契到了极点的下界修士！而且，他们修炼的体系，似乎对灵枢径有着某种先天的克制！

恐惧，第一次真正攫住了这位来自灵界的法相境修士。

“你们……你们究竟是谁？！哪个宗门派来的？可知我乃灵界界门……”徐峰厉声喝问，试图用身份震慑对方，同时手忙脚乱地一拍腰间储物袋，祭出了一面巴掌大小、通体赤红、铭刻着火焰纹路的小盾法宝。小盾迎风便长，化作一面门板大小的火焰巨盾，挡在他身前，散发出灼热而稳定的灵压，试图稳住摇摇欲坠的阵脚。

白言冰与陈千羽甚至连回答的兴趣都欠奉。两人只是极其短暂地对视了一眼，目光交汇的刹那，十二年共同修行、三十年朝夕相处培养出的无间默契便已完成了所有的交流。

陈千羽双眸之中碧光大盛，手中印诀猛然一变！

“咯吱吱——！”

缠绕在徐峰身上的所有藤蔓，在这一刻如同得到了号令的蟒群，同时爆发出惊人的力量，猛然向内收紧！那些尖锐的倒刺深深嵌入徐峰的护体灵光，甚至刺破法袍，扎入皮肉！ 不仅如此，所有藤蔓表面瞬间分泌出一种淡绿色的、散发着奇异甜腥气息的汁液，强烈的麻痹与阻滞神魂的效果通过伤口和灵光缝隙，疯狂向徐峰体内渗透！徐峰只觉得身体一沉，四肢百骸传来酸软无力之感，连法力的催动都变得滞涩起来，那面赤红火盾的光芒也随之黯淡了一分！

就在徐峰被藤蔓骤然收紧和麻痹气息弄得身形微僵、心神分散的极好时机，白言冰的身影便如同融入寒风的幽灵，只留下一道淡淡的残影，瞬间切入了徐峰因操控火盾而微微敞开的中门！他与徐峰之间的距离，从数丈缩短到三尺，仿佛从未存在过空间间隔。

并指如剑，指尖吞吐着凝练到极致的寒芒，不带一丝烟火气，却快如惊雷，直点徐峰胸腹之间的膻中大穴！

指尖未至，那凝聚于一点的恐怖剑意已然先行触及徐峰的肌肤与神魂，带来一种针砭般的刺痛与深入骨髓的寒冷预感，让他神魂都为之颤栗！

徐峰亡魂大冒，几乎出于本能地狂催法力，那面赤红火盾光芒一闪，就要回防格挡这近在咫尺的致命一指。

然而，白言冰仿佛早已预料。他点出的手指在距离火盾还有寸许距离时，极其精妙地、以一个微小到近乎不可能的角度轻轻一转。指尖那凝实的寒冰剑意骤然由实转虚，化整为零，仿佛水滴融入沙地，竟凭空绕过了炽热的火焰盾面，从另一个刁钻的角度，再次点向徐峰毫无防护的右侧肋下！与此同时，他脚下踩着的、由陈千羽操控的数根主藤蔓猛地向下一拽！

上下交攻，虚实变幻！

徐峰本就因麻痹而反应慢了半拍，此刻更是顾此失彼。他慌忙想侧身避让肋下袭来的虚指，脚下却被藤蔓狠拽，身形顿时失去平衡，一个明显的趔趄，胸前空门大开！

就是这电光石火间露出的、稍纵即逝的破绽！

对于将通明洞察力运用到战斗本能中的白言冰而言，这破绽如同暗夜中的灯塔般耀眼。他并拢的双指如同穿花蝴蝶，连点而出，速度快得在空中拖曳出一片模糊的指影幻象。

“噗噗噗噗”一连串轻微却沉闷的声响，如同雨打芭蕉，精准无比地落在徐峰周身十几处关键大穴之上。

每一指落下，都有一股精纯凝练、冰冷坚韧的真元，如同最锋利的冰针，强行破开徐峰体表残存的灵光防御，刺入他的穴窍深处。这些真元并非简单封堵他的气血运行，更带着一种独特的截断与凝固特性，如同在奔流的江河中打下无数坚固的冰桩，瞬间扰乱并强行截断了他灵枢与外界天地灵气之间那精密的勾连通道！

“嗬……嗬……”徐峰双眼骤然凸出，布满血丝，脸上充满了极致的惊骇与难以置信。他只感觉周身奔腾如江河的澎湃灵力，在那一瞬间完全凝固了！就像是汹涌的岩浆突然被投入了绝对零度的深渊，瞬间冻结成坚硬的石块，再也调动不了分毫！一种从未体验过的、身为修士最为恐惧的无力感淹没了他：四肢沉重如灌铅，丹田空空如也，神魂与天地灵气的联系被硬生生掐断！

他堂堂法相境通灵期修士，竟然……竟然在这灵气稀薄、被他视为蛮荒之地的下界，被两个他眼中的“下界蝼蚁”，以这样一种干脆利落、近乎羞辱的方式，联手生擒了？！

“唔！”他想要怒吼，却只发出一声含糊的闷哼，眼前阵阵发黑，软软地向前栽倒。

陈千羽素手轻扬，早已准备好的几根最为坚韧、闪烁着金属光泽的深绿色藤蔓如灵蛇出洞，瞬间缠绕而上，将徐峰的双臂反剪到背后，捆得结结实实，连手指都无法动弹一下。

白言冰上前一步，俯视着瘫软在雪地泥泞中、面如死灰的徐峰，神色依旧平静。他并指连点，又在徐峰后颈、丹田等数处要害补下了几道更为复杂的真元禁制，确保其体内哪怕有一丝一毫的灵力异动都会被瞬间镇压，彻底断绝其任何反抗或自戕的可能。

做完这一切，他才轻轻舒了一口气，抬起头，与不远处的陈千羽目光相触。

陈千羽眉眼弯弯，散去神通，林中疯狂滋长的藤蔓迅速枯萎、化为精纯的木灵之气消散，只在徐峰身上留下了那几根特制的束缚藤索。荒林恢复了一片狼藉但宁静的景象，仿佛刚才那场短暂却凶险激烈的战斗从未发生。

二人相视一笑，无需言语，一切尽在眼中。默契配合，效果斐然。以通明境修为，凭借功法特性与精妙战术，成功擒拿一名法相境通灵期修士，这份战绩虽不知高低，但足以自傲。

白言冰弯腰，像提起一条死狗般，单手抓住徐峰后背的衣襟，将他拎了起来。入手沉重，但对于他此刻的体魄而言，轻若无物。

“走，”他对着陈千羽微微点头，声音平和，“带回去，交给师尊发落。”

两人身影一晃，化作一青一蓝两道不易察觉的流光，掠过荒林枯枝，悄无声息地朝着云霞镇的方向返回，眨眼间便消失在茫茫雪景与铅灰的天色之中。

\vspace{2em}

与此同时，云霞镇中心，一座临街的二层茶楼，成了风暴眼中一处奇异的宁静点。茶楼门面古朴，此刻一楼大堂里坐着三三两两的茶客低声交谈着，人人脸上都带着挥之不去的忧色。而在二楼一处临街的雅座，叶牧谦独自一人，凭窗而坐。

他面前的红木方桌上，只摆着一壶冒着袅袅热气的清茶，两碟最普通的盐水花生和炒瓜子。他手持白瓷茶杯，慢条斯理地浅啜一口，目光似乎落在窗外萧瑟的街道上，神态悠闲惬意，仿佛只是一位闲来无事、消磨时光的普通茶客。

然而，这只是表象。他强大的神识早已无声无息地延展出去，悄然而精准地覆盖了整个云霞镇的每一个角落。镇东荒林方向传来的细微灵力波动、生机领域的骤然展开、清冽剑意的炸裂，以及随后迅速平息的禁锢……所有这些，都如同水面漾开的涟漪，清晰地倒映在他浩瀚如海的神识之中。而云来阁客栈里，另一个目标周锐那百无聊赖、充满期待的等候姿态，也尽在他的掌握。

他坐在这里，既是为执行荒林行动的白言冰二人掠阵接应，以防不测，同时也是监控全局，预防任何计划外的变故。窗边的位置，进可俯瞰长街，退可瞬息远遁，是最佳的观测与策应点。

楼梯口传来一阵轻微的脚步声，这脚步声与茶楼里其他茶客的沉稳或匆忙不同，带着一种刻意放柔的韵律，随之飘来的，还有一丝馥郁甜腻、与清茶香气格格不入的脂粉香风。

叶牧谦眼帘都未抬一下，依旧专注于杯中澄黄的茶汤，仿佛那里面有着比来人更值得关注的天地。

来者正是柳媚。她因实在厌烦了徐峰与周锐那毫无品味、白日宣淫且手段暴虐的低级行径，觉得与那两人为伍有失自己身份，便独自出来，在这小镇上闲逛，名为散心，实则也是想看看这蛮荒下界有没有什么能稍稍入眼的乐子。逛了半晌，只觉得满街都是衣衫简朴、面带惊惶的土包子，空气中都弥漫着一股穷酸的土腥味。正感无趣，打算折返时，不经意间抬眼，目光扫过茶楼二楼的窗口。

只一眼，她的脚步便顿住了。

窗边，一位身着朴素青衫的男子独坐。侧脸望去，线条清晰分明，鼻梁挺直，下颌的弧度带着一种干净的利落感。他只是安静地坐着，手里捧着一杯茶，可周身却萦绕着一股沉静淡然的气质，仿佛有一层无形的屏障，将他与茶楼里略显压抑的喧嚣、与窗外灰色沉郁的街景，彻底隔绝开来，自成了一片静谧的小天地。在这满是土腥味的下界人群中，此人简直鹤立鸡群，那份独特的气场瞬间就抓住了柳媚的目光。

柳媚凤目微微一亮，心头那股无聊瞬间被一股新的兴趣取代。

她舔了舔红唇，心中暗忖：徐峰周锐那两个只知道用下半身思考的蠢货，眼里只有那些瑟瑟发抖的庸俗女子，手段粗鄙至极。我柳媚岂能与他们一般？这男子气质如此独特，皮相更是上佳……即便他只是个空有皮囊的下界凡人，能有这般容貌气质，采补了助我修行，或是单纯吸干了取乐，也绝不亏本。

她对自己的魅术与手段向来极有信心，以往在灵界，不知多少修为心性尚可的男修，最终都心甘情愿地拜倒在她的石榴裙下，成为她的垫脚石。

一个下界男子，还不是手到擒来？

主意既定，柳媚嘴角勾起一抹志在必得的浅笑。她款步上楼，木质楼梯发出有节奏的轻响。她径直走向叶牧谦所在的方桌，毫不客气地在那空着的长凳上坐下，与叶牧谦隔着桌子相对。她将手肘支在桌面上，托着香腮，嫣然一笑，声音刻意调整得娇柔婉转，仿佛掺了蜜糖：“这位公子，独自一人品茶，岂不寂寞？小女子路过此地，见公子气质不凡，心中仰慕，可否……叨扰一杯清茶？”

眼波流转，眸光潋滟，在叶牧谦脸上流连，暗含的挑逗之意几乎要满溢出来。

叶牧谦这才仿佛被声音惊动，缓缓抬起眼皮，看了她一眼。那眼神平静得如同古井深潭，既无寻常男子乍见绝色时的惊艳与痴迷，也无对唐突之人的厌恶与不耐，平淡得就像扫过桌角的花生碟子，或者窗外一块寻常的街砖。

仅仅一眼之后，他便重新垂下眼帘，目光落回自己的茶杯，端起来又抿了一口，仿佛眼前坐着的不是一位千娇百媚、主动搭讪的美人，而是一团空气。

柳媚脸上那精心勾勒的嫣然笑容，瞬间僵住了。一丝难以置信的错愕爬上她的眉梢。

她……竟然被无视了？

如此彻底、如此理所当然的无视！

尤其对方还是个她打心眼里瞧不起的下界之人！

一股被严重冒犯的羞恼顿时冲上心头，烧得她脸颊微热。她深吸一口气，压下怒意，决定不再玩什么欲擒故纵的把戏，直接动点真格的。看来这男子是个不解风情的木头，或者是在故作矜持？

她微微调整了一下坐姿，腰肢轻扭，曲线毕露。同时，她悄然运转功法，一股若有若无、甜腻得仿佛熟透的蜜桃快要腐烂时的香气，开始从她周身肌肤散发出来，丝丝缕缕地飘向对面的叶牧谦。

她的眼眸在这一刻变得更加水润明亮，瞳孔深处似乎有旖旎的流光辗转，红唇轻启，吐气如兰，声音比刚才又酥软了三分，仿佛带着无数细小的钩子，能一直挠到人心底最痒的地方：“公子……可是嫌弃小女子蒲柳之姿，入不得您的眼？”她身子又前倾了几分，领口若隐若现，“您看，这长夜漫漫，寒寂无聊……何不寻些……彼此都能快活的趣事呢？小女子……定然让公子您，终身难忘……”

言语露骨挑逗的同时，那源自灵界的高阶魅惑术已被她全力催动！此术竟能直接撩拨心神，放大内心潜藏的欲望与情感漏洞；寻常心志不坚的修士，哪怕修为略高于她，稍有不慎也会中招，意乱情迷。

然而，预想中对方眼神迷离、呼吸急促的场景并未出现。对面的青衫男子，依旧保持着端坐品茶的姿势，甚至连握着茶杯的手指都没有颤动一下。那股足以让寻常法相境修士都心神摇曳的魅惑之力，在触及他身周三尺之地时，便如同泥牛入海，悄无声息地消散湮灭，没能激起半分涟漪。

不仅如此，叶牧谦甚至几不可察地微微蹙了一下眉头，仿佛茶楼里突然飘进了一股令人不悦的、低劣的异味，让他品茶的心情受到了些许干扰。

柳媚心中猛地一震，先前的羞恼瞬间被一股冰冷的警觉取代！

不对！这男人绝非凡人！甚至可能不是普通的下界修士！

自己的魅惑术竟然对他毫无作用？这怎么可能？！

她终于后知后觉地意识到，自己可能犯了一个严重的错误，踢到了一块看不透的铁板。她立刻收敛了所有媚态，眼神变得锐利，下意识地就想调动体内灵力，先行戒备，甚至准备抽身撤退！

然而，就在她心神从魅惑转换到警惕、灵力将发未发的这一刹，异变突生！

她突然感到周身经脉中奔流的灵力猛地一滞！仿佛一瞬间从顺畅的江河变成了粘稠的胶泥，运转变得无比晦涩迟缓，任凭她如何催动，都难以如臂指使。与此同时，她耳边传来了那青衫男子平淡如水的声音，可这句话，却不是对她说的：

“好了瑶儿，戏看够了。捆紧点，带回去吧。”

柳媚这时才发现自己后背被贴上了一枚符箓，惊骇欲绝地扭头，只见自己身侧——原本空无一物的空气里，如同水波荡漾般，光影一阵扭曲，悄无声息地多出了一位身着鹅黄色衣裙的青年女子，看上去约莫二十五六岁，容貌精致绝伦，眉目如画，皮肤胜雪，其美丽竟丝毫不输于她。

此刻，那青年女子一双妙目圆睁，狠狠地瞪着她，眼神里充满了毫不掩饰的鄙夷、恼怒，以及一丝计划得逞的小小得意。

而直到此时，柳媚才骇然发现，以她自己为中心的方圆数尺之内，空气都在微微地、持续地扭曲着，泛起一层层肉眼难辨、却真实存在的水波状涟漪——她早已在毫无察觉的情况下，落入了对方精心布置的幻术与禁制交织的天罗地网之中！

正是奉师命潜伏在侧的沈瑶！她凭借千幻流光之术，将自身气息与周围环境完美融合，如同变色龙般隐藏在侧，耐心等待着柳媚心神松懈、全力施展魅术而疏于防范的最佳时机，悄然布下无形罗网，此刻骤然收网，一击即中！

“你……！你们……！”柳媚想挣扎，想怒斥，但灵力被封，身体也被无形的禁制力量束缚，让她连动弹一下都困难。

极度的恐惧淹没了她，她慌忙又看向叶牧谦，眼中挤出泪水，声音发颤，再没有了丝毫媚意，只剩下苦苦哀求：“前……前辈饶命！是晚辈有眼无珠，冒犯了前辈！我……我不知道您有道侣在此！我向您道侣道歉！向这位仙子道歉！求您高抬贵手，饶了我吧！”

她误将沈瑶当成了叶牧谦的道侣。

叶牧谦闻言，原本平淡的脸上不禁流露出一丝又好气又好笑的神情，摇了摇头：“现在说这些，却饶你不得。”

沈瑶同样被柳媚那句“道侣”和之前的所作所为气得脸颊微鼓，清脆的声音带着怒意：“你这妖女，在我面前就敢如此放肆，勾引我师尊，胆子倒是不小！”

素手一扬，一道流光溢彩、不知何种材质炼制的绳索状法宝从她袖中如灵蛇般激射而出，在空中划过一道绚丽的弧线，瞬间将柳媚从头到脚，捆了一圈又一圈，勒得紧紧的，更是将她体内残存挣扎的灵力彻底锁死。

柳媚被捆得结实实，心中恐惧更甚，眼珠子急转，看到沈瑶的脸庞，忽然像是抓住了一根救命稻草，哪怕嘴巴被勒得有些变形，还是急急开口，声音含糊却带着诱惑：“我说……仙子！你……你放了我，我……我教你！我教你几套顶尖秘术，保证……保证你能把你师尊服侍得舒舒服服、欢欢喜喜，让他眼里再也看不下别的女人！怎么样？你把我放了，我把秘术全告诉你！绝对是真传！”

在她看来，这女子既是这强大男子的徒弟兼道侣，必然极在意对方，以此诱惑，或许有一线生机。

叶牧谦听得简直是哭笑不得，沈瑶更是瞬间整张脸连同脖子根都红透了。她心念一动，那绳索立刻又延伸出一截，如同活物般，精准地缠绕而上，将柳媚还在试图喋喋不休的嘴巴也严严实实地封住了，终于只剩下“呜呜”的闷哼声。

世界清净了。

叶牧谦这才放下一直握在手中的茶杯，杯底与桌面轻轻磕碰，发出一声轻响。他站起身，对沈瑶点了点头，目光中带着赞许：“走吧，去和言冰他们会合。”

沈瑶脸上的红晕稍褪，应了一声，手指一勾，那流光绳索便拽着被捆成粽子、满眼绝望的柳媚，轻若无物地飘起，跟在叶牧谦身后。两人一俘下楼，叶牧谦随手在桌上放下几枚铜钱结账，便径直离开了茶楼，身影很快融入街道，朝着云来阁的方向而去。茶楼里其他客人，竟无人察觉二楼雅座方才发生了一场无声的、碾压般的擒拿，只当那对相貌出众的男女带着一个被布帛包裹的“行李”离开了。

\vspace{2em}

云来阁客栈。

大堂里，周锐已经独自解决了整整一桌酒菜。

他惬意地靠在椅背上，剔着牙，面前又新上了一壶号称是店家珍藏的“十年陈酿”，自斟自饮，好不快活。他脑子里反复勾勒着徐峰师兄此刻在荒林中的场景——必定是以猫戏老鼠的姿态，将那不知死活的布衣小子折磨得哭爹喊娘，跪地求饶，最后再像捏死虫子一样了结他。

然后嘛……那个清丽脱俗的小娘子，便成了惊惶无助的待宰羔羊，徐师兄享用过后，按照约定，自然会留给自己……想到那女子清澈灵动的眸子和那难掩的风韵，周锐小腹便是一阵火热，嘴角咧开，露出得意的笑容，仿佛美人已然在怀。

“徐师兄怎么去了这么久？连个传讯符都没发回来……”他等得有些心焦不耐，低声嘀咕了一句，但旋即又释然了，“啧，定是玩得兴起，忘了时间。两个手无缚鸡之力的凡人，还能翻了天不成？说不定徐师兄正在慢慢‘教导’那小娘子呢，嘿嘿……”

他完全没往坏处想，只觉得是徐峰一时兴起，拖延了时间，甚至脑补出更多画面，让他更加心痒难耐，只好又猛灌了一口酒，压一压躁动。

就在他端起酒杯，准备再饮一口时，客栈门口的光线，忽然暗了一下。

周锐随意地瞥了一眼，起初并未在意，以为是又有客人进出遮挡了光线。但下一秒，他眼角的余光捕捉到了走进来的人影，动作瞬间僵住。

他的目光呆滞地移动，从当先那位气质沉静的青衫男子，扫到他身旁鹅黄衣裙、容貌极美的青年女子，再看到她手中一道流光牵着的、被捆得严严实实、嘴里塞着东西、满脸惊恐灰败的……柳媚？！

周锐的脑袋“嗡”的一声，像是被重锤砸中。

他的视线继续向后挪，心脏几乎停止了跳动。青衫男子身后，跟着的正是那对“布衣夫妇”！可此刻，那男子手中提着的……那个像死狗一样瘫软着、人事不省、衣着华贵却沾满尘土雪沫的人……不是徐峰师兄，又是谁？！

四人进门，加上被提着的徐峰和被捆着的柳媚，六个人，十二道目光，在这一刻，齐刷刷地、冰冷地聚焦在了独自坐在大堂酒桌旁、一只手还握着酒杯的周锐身上。

时间仿佛凝固了。客栈大堂里原本零星的几个伙计和客人，早已吓得缩到了墙角，大气不敢出。

“啪嗒！”

周锐手中的白瓷酒杯脱手坠落，在桌面上滚了两圈，摔落在地，发出一声脆响，碎裂开来。冰凉的酒液泼洒在他珍贵的法袍下摆上，浸湿了一片，他却浑然不觉。

他只感觉一股寒意从脚底板直冲天灵盖，瞬间蔓延全身，四肢百骸一片冰冷，如坠万丈冰窟！ 徐峰师兄……法相境通灵期的徐峰师兄，竟然被生擒了？！看那样子，毫无反抗之力！柳媚也落得如此下场！那对“凡人”夫妇果然是伪装的！还有这突然出现的青衫男子和黄衣女子……他们身上没有刻意散发任何强大的气息，但正是这种深不可测的平淡，让周锐本能地感到了一种近乎窒息的心悸，尤其是那青衫男子，明明只是平静地站着，却仿佛是整片空间的中心，让人望之生畏。

逃！必须逃！这个念头刚一升起，就被绝望掐灭。对方四人能如此轻松地拿下徐峰和柳媚，实力绝对远超自己。他几乎是刹那间就在心里完成了实力对比：徐峰师兄通灵期都被擒，自己仅仅是外显期，对方至少有四人，个个气息沉稳，眼神锐利，一看就不好惹……这还怎么打？任何反抗都无疑是螳臂当车，自取灭亡！

灵界养尊处优、视下界众生如蝼蚁草芥的傲慢，在同伴被擒的残酷现实和绝对的实力差距面前，顷刻间土崩瓦解，碎得连渣都不剩。

生存的本能压倒了一切。

“我……我……”周锐脸色煞白如纸，冷汗如同瀑布般从额头、鬓角渗出，瞬间浸湿了内衫和后襟，身体控制不住地微微发抖。他非常、非常识时务地，猛地从椅子上滑下来，几乎是瘫跪在地，高高举起了双手，挤出一个扭曲的、比哭还要难看十倍的笑容，声音干涩嘶哑，充满了恐惧与哀求：“我……我投降！别……别动手！我认栽！我什么都配合！求……求各位前辈、仙子高抬贵手！” 他毫不犹豫地放弃了所有抵抗的念头，只求能保住性命。

叶牧谦神色依旧平淡，只是微微颔首，算是接受了他这毫无骨气的投降。白言冰上前一步，眼神锐利如鹰隼，先是在周锐身上迅速扫视一圈，确认他没有隐藏什么自爆或传讯的法器符箓，然后才出手，快速而专业地搜查了他的全身，卸下了腰间储物袋和几件零碎佩饰。

陈千羽则上前，素手一翻，指间现出数根细如牛毛、却闪烁着淡淡金芒的长针，她出手如电，精准无比地将金针刺入周锐周身数处关键大穴。金针入体，周锐只觉浑身猛地一麻，随即一种空荡荡的虚弱感传来，灵枢与周身经脉的灵力联系被彻底截断，连抬起一根手指都变得困难无比，闷哼一声，彻底瘫软下去，眼中只剩下绝望的顺从。

叶牧谦目光扫过地上瘫着的周锐，又看了一眼被提着的徐峰和被捆着的柳媚，声音平静无波，却为这场干净利落的抓捕行动画上了句点：

“完事，一并带回去。”

\vspace{2em}

很快，云隐别院前的广场，已被临时布置成公审之所。虽无高堂明镜，但青石铺就的地面，环绕的古松，以及后方云雾缭绕的云霞山主峰，自有一股庄严肃穆之气。得到消息的云霞镇百姓、受害者的亲属、以及闻讯赶来的众多修士，将广场外围挤得水泄不通。人人脸上交织着悲愤、期待与一丝忐忑。

辰时正，叶牧谦、沈天穹、吴东浩等人于临时搭建的审案台后落座。白言冰、陈千羽、沈瑶侍立两侧，神色肃然。随着一声清越的钟鸣，全场渐渐安静下来。

“带人犯！”沈天穹沉声开口，声音不大，却清晰地传遍全场。

在众人灼灼目光注视下，徐峰、周锐、柳媚三人被巡天卫押了上来。他们身上的禁制并未完全解除，灵力被封，步履略显虚浮。徐峰神情带着惯有的倨傲，柳媚满脸惊疑不定，周锐则彻底服了软，双目空洞。华丽的灵界服饰沾了尘土，显得有些狼狈。

沈天穹作为昔日天师府府主、如今内域联盟德高望重的前辈，主持此次公审。他目光如电，扫过三人，缓缓开口，声音带着沉痛与威严：“徐峰、周锐、柳媚，尔等自灵界而来，于云霞镇滞留七日。现有苦主多人控告，尔等犯下杀人、奸污、伤人、抢劫等累累罪行。人证物证俱在，尔等可认罪？”

徐峰冷哼一声，昂着头：“是又如何？不过几个下界蝼蚁，杀了便杀了，玩了便玩了，算什么大事？”他语气满不在乎，仿佛在说踩死了几只蚂蚁。

柳媚抿了抿唇，没有像徐峰那样嚣张，但也没有否认，只是微微侧过头，避开了那些充满恨意的目光。

周锐则只是沉默，似乎是吓傻了。

他们的态度，如同火上浇油，激起了更大的愤怒。人群中传来压抑的哭泣和怒骂声。

沈天穹强压怒火，示意证人上前。一位失去女儿的老妪颤巍巍地指着周锐，哭诉女儿如何被其调戏后掳走，次日发现惨死林间；被打成重伤、店铺被砸的掌柜，拄着拐杖，控诉徐峰吃饭住店不给钱还出手伤人；几位侥幸逃得性命、却目睹亲人被害的修士，也纷纷站出来指证……一桩桩，一件件，血泪交织，令人发指。

面对这些指控，徐峰起初还试图狡辩或嗤之以鼻，但在确凿的人证和从他们临时落脚点搜出的、带有云霞镇标记的财物等物证面前，终究无可抵赖。或许是觉得承不承认对这些“下界蝼蚁”而言并无区别，又或许是存着别的念头，徐峰最终对大部分恶行供认不讳，态度依然嚣张；周锐也认了罪，希望全数赔偿能饶他不死；柳媚则承认了自己旁观、未曾阻拦，以及同样有吃住不给钱、随手伤人的行为，但对徐峰周锐的奸杀暴行，她坚称自己并未参与。

罪证确凿，供认不讳。场间群情激愤，要求严惩的声浪一阵高过一阵。

\vspace{2em}

然而，在最终宣判之前，叶牧谦却抬手示意众人稍安勿躁。他看向台下三名囚犯，目光平静无波，开口问道：“在处置尔等之前，我有几个问题，关于灵界。”

此言一出，徐峰三人眼中骤然爆发出希冀的光芒！他们以为自己抓住了救命稻草——这下界强者果然对灵界心存好奇乃至敬畏！或许，他们可以利用信息作为筹码，换取一线生机！

“前辈请问！我等知无不言，言无不尽！”徐峰抢着回答，语气甚至带上了一丝谄媚，与方才的嚣张判若两人。周锐和柳媚也连忙点头，眼巴巴地看着叶牧谦。

叶牧谦不置可否，缓缓问道：“灵界确实存在？尔等来自何处，所为何来？”

“存在！千真万确！”徐峰忙不迭地回答，“我等来自灵界中州皇朝麾下的界门司，专司巡查下界、接引有潜力的后辈修士前往灵界修行！此次……此次便是奉命前来贵界，寻找资质上佳者接引！”

他刻意强调了“奉命”和“接引”，试图给自己披上一层官方使者的外衣。

殊不知，在听到“皇朝”的时候，叶牧谦眼前一亮。

竟然与他对设定集的记忆完全一致！

所以……

“灵界广袤，共分五域：东域、西疆、南领、北境，以及最核心富庶的中州；中州乃皇朝直接统御之地，其余四域名义上亦尊皇朝为主，但天高皇帝远，皇朝实际难以有效节制四域具体事务，对吧？”

叶牧谦这次是极其笃定地问道。

“对对对！您是来过灵界吗？怎么对这些事情了如指掌？”周锐抢着回应道，“皇朝难以节制外域，正是因为中州之外，四域之中，盘踞着无数‘宗门’与‘世家’等修士集团。按其底蕴、强者数量、掌控资源，大致分为一流、二流、三流。这些宗门世家往往并非单打独斗，而是进一步抱团，形成不同的利益联盟或阵营。其实力强弱，直接决定了他们在资源争夺、地盘冲突中的地位与话语权。”

叶牧谦努力遏制着自己的笑意。

这一通盘问，效果比他预期的还要好上数倍！

至少他设定集中勾勒过的内容，与现实竟然如出一辙！

他微微颔首，问出关键：“修行境界，法相之上，可还有路？”

柳媚接过话头，声音柔媚，努力展现自己的价值：“有！自然有！法相境在灵界亦非顶点。据晚辈所知，完整的大境界依次为：引灵枢、通玄脉、绘真符、现法相、植道种、转天轮、凝圣魄、入神墟、归源寂！不过……”她语气稍顿，“灵界能容纳的最高境界，便是天轮境。更高境界的大能，会前往所谓‘天界’谋进一步发展。”

叶牧谦再次确认道：“至于这个容纳……是不是因为你们灵枢径修士，一旦转移到灵气稀薄的地界，体内灵气就会逸散、无法支持高境界的消耗，从而导致跌落？”

几人更是惊诧起来。

这个人，不简单！八成是灵界某法相境大能和几位子弟下来探亲的！

只是……这样的大能，难道会问出这么基础的问题？

但眼下，表示认同才是对的。

于是几人疯狂点头。

沈天穹一直静静听着，此时忽然身体前倾，声音带着一丝不易察觉的颤抖，插言问道：“尔等可知，灵界可有一位名叫‘沈玉’的修士？”这是他数百年来深埋心底的执念。

徐峰三人面面相觑，努力回忆，最终皆茫然摇头。周锐道：“沈玉？未曾听闻。灵界修士亿万，名姓相似者众多……前辈，此人可有何特征？是何境界？若是无名之辈或年代久远，我等不知也是常理。”

他们三人最大的也不过二百岁，对更早年间的人物，自然毫无所知。

沈天穹眼中希冀的光芒黯淡下去，整个人仿佛瞬间苍老了些许，靠回椅背，默然不语，失望之情溢于言表。

叶牧谦又陆续问——与其说“问”，不如说是和自己的记忆相互印证——了一些关于灵界风土、资源、通行货币、各大势力知名强者等细节，徐峰三人为了活命，争先恐后地回答，甚至互相补充，生怕自己知道得少了，价值不够。

当然，叶牧谦也问了一嘴“天界”的事情，这几个人只知道所谓天界和灵界的通道是“天灵走廊”（甚至这个名称也和叶牧谦先前的设定集惊人的重合）。但具体天界如何，他们也没去过。

\vspace{2em}

问询持续了近一个时辰。三人自觉已将所知和盘托出，态度愈发“诚恳”，眼巴巴地望着叶牧谦等人，等待着“宽大处理”。徐峰甚至表示，若放他们回去，必在陆执事面前美言，为内域争取接引名额云云。

然而，叶牧谦问完后，只是平静地坐了回去，对沈天穹微微点头。

沈天穹深吸一口气，重新坐直身体，脸上恢复了冰冷与威严。他目光扫过满怀期待的徐峰三人，又看向台下无数双殷切的眼睛，朗声宣判，声音回荡在广场上空：

“人犯徐峰、周锐、柳媚，倚仗修为，自灵界而来，视我内域生灵如草芥，于云霞镇期间，杀人越货、奸淫掳掠、伤人毁物，罪恶滔天，证据确凿，尔等亦已供认不讳！”

“今依《问道戒律》之精神，循天理人情，判决如下：”

“徐峰，身犯故意杀人、奸淫、重伤、抢劫等多重罪孽，手段残忍，情节极其恶劣，罪无可赦！不杀不足以告慰亡魂，不足以平民愤！判处死刑，立即执行！”

“周锐，所犯罪行与徐峰类同，残忍暴虐，丧尽天良，亦属罪大恶极！判处死刑，立即执行！”

“柳媚，虽未直接参与杀人、奸淫，但同行共处，目睹暴行而不加劝阻，同流合污；且亦有伤人、抢劫之实。判处杖刑一百，削其修为，罚没随身财物以赔偿受害者损失，随后逐出内域，永不得再入！”

判决既出，徐峰三人如遭雷击，脸上的期待瞬间化为难以置信的惊骇与恐惧！

“不！你们不能杀我！我是灵界界门司使者！杀了我，皇朝绝不会放过你们！灵界大军顷刻即至，尔等下界弹丸之地，必将灰飞烟灭！”徐峰疯狂挣扎嘶吼，试图搬出灵界震慑众人。

周锐也面色惨白，尖声叫道：“放过我！我的家族在灵界颇有势力，你们敢动我，必遭灭门之祸！那些蝼蚁死了就死了，我们愿意赔偿！十倍！百倍赔偿！”

然而，他们的威胁在满腔悲愤的民众和意志坚定的审判者面前，显得苍白无力。台下响起震天的怒吼：“杀了他！为死去的乡亲报仇！”“血债血偿！”“灵界了不起吗？就可以随意杀人？！”

柳媚相对冷静，但脸色也极其难看，她咬牙道：“我已回答你们所有问题！你们岂能出尔反尔？一百杖，削修为……你们可知这意味着什么？”这对于习惯高高在上的灵界修士而言，是比死更难受的屈辱。

叶牧谦面无表情：“尔等罪行，自有律法公断。灵界又如何！”

徐峰周锐眼见求饶威胁无用，绝望之下，眼中凶光毕露，竟试图引爆体内残存的禁制，做最后的鱼死网破！然而，叶牧谦早已料到，轻轻一抬手，一股无形的力量瞬间加诸二人身上，将他们刚刚凝聚起的一丝紊乱灵气彻底掐灭，同时加固了禁制。两人闷哼一声，彻底瘫软，眼中只剩下无尽的恐惧与悔恨。

白言冰与陈默上前，将面如死灰的徐峰、周锐拖至广场一侧空旷处。没有繁复的仪式，白言冰并指如剑，一道清冽剑光闪过，徐峰头颅飞起；陈默则挥刀斩落周锐首级。两人元神早在禁制下被锁死，随着肉身死亡而彻底湮灭。

干脆利落，恶贯满盈者伏诛。

轮到柳媚。她被按倒在地，封住修为后，由巡天卫行刑。沉重的法棍落下，噼啪作响，每一棍都结结实实。柳媚起初还能咬牙硬撑，很快便忍不住发出凄厉的惨叫，臀背血肉模糊。一百杖毕，她已奄奄一息。陈千羽上前，以金针之术，配合特殊手法，将其苦修多年的法相境外显期修为根基彻底废去，只保留了一丝微弱的元气维系性命，从此与大道无缘。她随身的储物法器、华服首饰等一切有价值之物，皆被罚没，登记造册，准备用于赔偿受害者。

整个过程，广场上寂静无声，只有法棍着肉声与柳媚的惨呼。当行刑完毕，柳媚像破布一样被拖走时，不知是谁先喊了一声：“好！”

随即，“青天大老爷！”“巡天司万岁！”“叶仙师公义！”的欢呼声、掌声、痛哭声轰然爆发，如山呼海啸，响彻云霞山麓！积压了七日的恐惧与悲痛，在这一刻化为对正义得以伸张的由衷快意与宣泄。许多人相拥而泣，告慰逝去的亲人。

叶牧谦等人起身，看着台下激动的人群，心中并无多少喜悦，唯有沉重。灵界之门已开，恶客虽除，但更大的风波或许还在后头。然而，今日之举，至少昭示了一点：无论来自何方，践踏此界生灵者，必受严惩！

\vspace{2em}

当日，云隐别院。

疏散完聚集在广场上的百姓们之后，内域有头有脸的人物都回到了云隐别院。徐峰、周锐、柳媚三人留下的储物袋一字排开，东西被尽数倾倒在一张宽大的青玉案几上。灵石光华莹莹，符箓灵光隐现，法器形态各异，丹药瓷瓶剔透，还有不少灵界特有的华美衣物与零碎物件，堆成了几座小山。

“清点一下吧。”叶牧谦开口道，“凡俗金银折算之物，以及适合折现的灵石，悉数用于赔偿镇民损失，抚恤死者家属。务必足额，若有剩余，纳入巡天司公库，用于日后抚恤及公益之事。”

众人点头。吴东浩主动请缨，与陈默一同负责此事，他们熟悉庶务，也能确保公平。

“话说叶先生怎么知道那么多与灵界相关的事宜？”沈天穹最终还是没忍住，问道，“难道说叶先生早年间去过灵界游历？”

“这个，真没有。”叶牧谦道，“我和这片天地之间的联系……可能你们目前确实难以理解，不过之后，有机会，我还是会告诉你们的。”

于是这件事情就被搁置下来了。

叶牧谦的目光落在那些符箓上。灵界的符箓，无论是材质、纹路还是蕴含的灵力性质，都与内域乃至荒域所见大不相同，更为繁复精妙。“瑶儿，你心思灵巧，于幻术、变化之道颇有天赋，这些符箓便交予你研究，或能触类旁通，精进你的千幻流光之术。”

沈瑶眼睛一亮，仔细端详那些符箓，郑重点头：“是，师尊。弟子定当仔细参详。”

接着是那些法器。刀、剑、环、镜、印……形制各异，灵光内蕴，至少也是真阶水准，甚至有几件法阶法器，放在内域足以引起哄抢。叶牧谦随手拿起徐峰那面破损的赤红小盾，指尖流露出一丝真元探入，仔细感知其内部结构。

“炼器手法确有独到之处，材料也非下界常见。这些东西，我暂且留下，拆解研究一番，或许对提升我等的炼器水准有所助益。”他如今对炼器布阵兴趣浓厚，这些灵界法器正是极好的样本。

最后是那些丹药。瓷瓶上贴着标签，字迹是灵界文字，但通过灵气波动大致能判断用途。王丹作为丹道大家，早已按捺不住好奇，征得叶牧谦同意后，小心地取过几瓶，拔开塞子轻嗅，又倒出些许在玉盘上仔细观察。“药力精纯，君臣佐使的搭配思路与下界丹方颇有差异，有几味主药我竟从未见过……”她看向陈千羽，微笑道，“千羽，你精研生机平衡、医药之理，这些丹药便交由你与老身一同研究如何？或许能从中悟出新的疗伤、补益思路，甚至反推出一些灵界药草的特性。”

陈千羽欣然应允：“能与王阁主一同钻研，弟子求之不得。”

衣物等私人物品并无实际价值，但料子极佳，也蕴含微弱灵气，叶牧谦吩咐一并折价处理，所得同样并入赔偿款项。

就在清点接近尾声时，白言冰从一个锦囊的夹层里，摸出了十几枚非金非玉、入手温润的令牌。令牌约莫巴掌大小，造型古朴，正面镌刻着复杂的云纹与一个古老的“界”字，背面则是一些编号似的细小符文。

“师尊，这是何物？”白言冰将令牌递给叶牧谦。

叶牧谦接过，仔细感应。令牌本身材质特殊，内部结构精巧，似乎蕴含着某种稳定的空间道纹。他尝试着向其中注入一丝真元。

嗡——

令牌轻轻一震，正面云纹仿佛活了过来，缓缓流转，散发出柔和的白色光晕。紧接着，光晕在令牌前方尺许处的空中汇聚，竟缓缓拉开了一道仅容一人通过的、模糊扭曲的光之门扉！门扉那边，隐约传来截然不同的、更为浓郁精纯的灵气波动，以及一股浩瀚苍茫的隐约气息，与下界迥异。

“这是……空间通道？”沈天穹惊疑道。

叶牧谦维持着真元输入，仔细感知门扉的稳定性和对面的情况。片刻后，他撤去真元，光之门扉晃动几下，随即消散。“一次性的单向传送门，或者短期通行凭证。”他沉吟道，“根据令牌内道纹的提示，需在下界环境下，灌入足够强度的能量方能激发，目的地应是固定指向灵界某处……很可能就是那‘界门司’的接引点或附近。”

他逐一检查了十几枚令牌，结构基本一致，只是背面编号不同。“没有被动过手脚的痕迹，没有隐藏的追踪或反制符纹。”叶牧谦得出结论，眼中闪过一丝锐芒，“这可真真正正是好东西。恐怕是那界门司配发给接引使者的标准制式物品，用于完成任务后返回，或紧急情况下启用。徐峰三人下来时用过一枚，这些是备用的。”

拥有能自行打开通往灵界门户的令牌，意味着主动权！不必等待灵界不定期的、可能充满不确定性的“接引”，而是可以自主选择时机前往！

室内的气氛变得有些微妙。灵界，那个传说中的上界，环境优渥，大道昌盛，是无数下界修士向往之地。然而，刚刚发生的血腥事件，徐峰三人视人命如草芥的傲慢与残忍，却给这层向往蒙上了厚厚的阴影。

叶牧谦缓缓摩挲着一枚令牌，目光扫过众人：“经此一役，灵界……至少是这‘界门司’以及类似徐峰之流的做派，诸位也看到了。弱肉强食，视下界为刍狗，绝非善地。”

众人默然，沈天穹、吴东浩等人脸上都露出沉痛与愤慨之色。沈瑶更是握紧了拳头，她对这种恃强凌弱、践踏弱小的行径深恶痛绝。

“但是，”叶牧谦话锋一转，语气平静却坚定，“无论如何，灵界是真实存在的，是更广阔的天地。我们不能因噎废食，因为几个恶徒，就永远蜷缩在下界，对上方世界闭目塞听。至少，我们得亲眼上去看看，看看那所谓的上界，到底是个何等模样，其‘坏’，又到了何种程度？其‘好’，是否又有值得借鉴之处？只有亲眼所见，亲身所感，方能做出准确的判断，决定日后如何自处，是闭关锁域，还是有限交流，抑或其他。”

他看向自己的三位弟子：“我打算，用这些令牌，上去一探。人数不宜多，就我们师徒四人。内域还需诸位坐镇。”

白言冰、陈千羽、沈瑶相互对视，眼中并无惧色，反而燃起熊熊的探索之火。他们本就是因“人外有人，天外有天”的信念才走出荒域，来到内域，如今有通往更广阔世界的机会，岂会因前路可能艰险而退缩？更何况，师尊同行，心中自有底气。

“弟子愿往！”三人异口同声，深以为然。

“好。”叶牧谦点头，“此事需谨慎。我们离开后，内域不可无人主持。沈府主、吴盟主，还有陈默、吴刚，内域联盟及巡天司事务，就劳烦诸位多多费心了。我们会尽快返回，短则数月，长则数年。若遇变故，也会设法传递消息。”

接下来几日，云隐别院与巡天司开始了紧张而有序的事务交割与准备。叶牧谦将一些重要的阵法维护、丹药调配心得留给了王丹和沈天穹；白言冰与陈千羽将巡天司日常巡视、案件处理流程与陈默吴刚做了细致交接；沈瑶则整理了部分幻术与禁制的研究笔记。同时，四人也在默默调整状态，准备应对灵界可能截然不同的环境与规则。

\vspace{2em}

灵界，中州，界门司某处分衙。

光影晃动，形容狼狈、气息萎靡到了极点的柳媚，从一座闪烁不定的传送阵中跌跌撞撞地滚了出来，趴在地上剧烈咳嗽，背后血肉模糊的杖伤虽然经过初步处理不再流血，但剧痛与修为被废的虚弱感让她几乎昏厥。

值守的修士吓了一跳，待看清是柳媚，更是大惊失色，连忙上前搀扶，并迅速上报。

很快，柳媚被带到了此处分衙的主事——陆洲面前。陆洲看上去四十余岁模样，面容清癯，三缕长须，身着界门司执事官袍，气息渊深，已达道种境凝道期。他看着跪在下方面无人色、瑟瑟发抖的柳媚，眉头紧锁。

“怎么回事？徐峰和周锐呢？怎么就你一个人回来？还弄成这副模样！”陆洲声音不高，却带着不容置疑的威严。

柳媚早已打好了腹稿，闻言立刻抬起苍白的脸，眼中挤出泪水，声音虚弱而惊惶：“陆主事……属下……属下等人奉命下界巡查接引，不料……不料那下界竟有埋伏！有数名实力强横的法相境修士早已守候在侧，不由分说便偷袭我等！徐师兄和周师兄奋力抵抗，奈何对方人多势众，又熟悉地形，两位师兄……两位师兄不幸罹难！属下……属下拼死才侥幸逃脱，回来报信……”

她绝口不提徐峰周锐的恶行，更不提自己也曾伤人抢劫，只将一切归咎于下界修士的“凶残”与“埋伏”，将自己塑造成唯一的、忠诚的幸存者。

陆洲静静地听着，手指无意识地敲击着椅背，眼神深邃。他执掌界门司分衙多年，对徐峰、周锐的品性岂会不知？那两人仗着有点背景和修为，平日里就好惹是生非，骄横跋扈。派他们下界，本就是件不太稳妥的差事。柳媚的话，他最多信三成。

什么下界埋伏？恐怕是徐峰周锐这两个不着调的，下去之后又管不住手脚，动了不该动的东西，睡了不该睡的姑娘，惹了不该惹的人，结果踢到了铁板，被人给收拾了！柳媚嘛，估计也是从犯，能逃回来已是侥幸。

想到这里，陆洲心头一阵烦闷。他新官上任，本想借这次下界接引做出点成绩，压服那些暗地里不服的老油条，谁知派出去的人竟如此不成器，捅出这么大篓子！死了两个法相境，虽然不算顶尖战力，但也是界门司的正式编制，后续的麻烦少不了。

他盯着柳媚，语气转冷，带着训斥：“埋伏？哼！柳媚，你当本座是三岁孩童？徐峰、周锐是什么德行，你不清楚？你们三人一同下界，他们若行事荒唐，你为何不加劝阻？如今酿成大祸，折损人手，丢尽我界门司颜面！你们这样，还有一丝廉耻吗？灵界的脸都被你们丢光了！”

柳媚伏地颤抖，不敢辩驳。

陆洲越说越气：“打也打不过，跑也跑不掉？你身为同僚，眼见他们要惹祸，就不知道劝劝那两个不成器的家伙？哪怕劝不住，早些传讯回来也好！如今倒好，人死了，差事也办砸了！”

发泄了一通，陆洲也知道事已至此，迁怒无益。他冷冷道：“念在你重伤逃回，尚有报信之功，死罪可免，活罪难逃！罚你禁闭一月，好好反省！滚下去！”

柳媚如蒙大赦，忍着剧痛和屈辱，被人搀扶了下去。

空荡荡的厅堂内，陆洲揉了揉眉心，叹了口气。下界居然有能击杀徐峰周锐的法相境修士？还不止一个？这倒是出乎意料。按照以往那些怠政前任的报告，那片下界应该灵气稀薄，道统衰微，连玄脉境都少见才对。看来，那片下界似乎出了些变故，或者之前的情报有误。

他原本听说下界可能有法相境，还动过亲自走一趟、以示郑重招揽的念头。毕竟若能接引到真正的天才，也是大功一件。但现在……

看看徐峰周锐和柳媚搞出的破事！下界之人恐怕早已对灵界“界门司”恨之入骨，印象跌入谷底。这时候自己再去，岂不是送上门找不痛快？别说接引了，能不能和平对话都是问题。万一对方脾气暴烈，连自己也扣下甚至打杀，那才真是天大的笑话。

“唉……”陆洲长叹一声，满脸无奈。好好一桩可能立功的差事，硬是被这几个蠢货搞成了烫手山芋，甚至成了外交污点。

没办法，他只得硬着头皮，将此事原原本本（当然，采用了柳媚汇报的版本，强调下界修士凶悍埋伏，袭击上界使者），向上级呈报。

报告很快到了界门司更高层的手中。几位主事长老翻阅后，脸色都不太好看。折损人手是小，丢脸是大！堂堂灵界界门司使者，被下界修士给杀了？这传出去，界门司乃至皇朝的脸面往哪搁？尤其是，事情的起因很可能是己方使者行为不端……这就更难以启齿了。

经过一番商议，高层做出了决定：此事不宜声张。对外，仍沿用前任的结论，宣称下界“灵气贫瘠，道统不显，暂无资质上佳者可供接引”。至于徐峰周锐的死亡，则内部处理，找个由头遮掩过去。那片下界，暂时列入“观察区”，待风波平息，或者找到合适的契机（比如换一拨完全不知情、态度也更谨慎的使者），再行接触。

就这样，一场因恶客引发的下界风波，在灵界官方层面，被悄然掩盖。叶牧谦师徒四人即将前往的，是一个对近期发生在下界的冲突“一无所知”，却又在底层规则上默认弱肉强食的陌生世界。

\vspace{2em}

云隐别院中，四道身影在晨光中聚拢，目光投向那十几枚安静的令牌。

“东西都准备好了吗？”叶牧谦问陈千羽和沈瑶。

“衣物、灵石等物都准备好了。”沈瑶回答。

“事务交割呢？”这次是问白言冰，后者点了点头。

叶牧谦了然。他拿起一块令牌。这枚令牌的内部阵纹被他特意改了改，使预计抵达地点比原来向东偏移了三万里。他向令牌中灌入真元，那道门扉很快再次出现。

“来！”

三位弟子领命，依次踏入门扉之中，身影消失。叶牧谦也紧随其后，随即门扉彻底关闭，就像从未出现过一样。

灵界的旅途，开始了。

\vspace{2em}

踏入门扉之后，短暂的失重与空间扭曲感传来。

周遭光怪陆离，仿佛穿过了一条由无数流光构成的狭窄通道。

约莫过了十息，脚下一实，周遭景象骤然清晰。

\vspace{2em}

首先感受到的，是扑面而来的浓郁灵气。这灵气之精纯、之充沛，远超内域，更非荒域可比，呼吸间便觉周身毛孔自然舒张，若换作灵枢径修士在此，恐怕会欣喜若狂。

然而，叶牧谦的眉头却微微蹙起。

与初入内域时相似，但强烈了十倍的不适感涌上心头——那是一种源自大道感知层面的“破残”之感。此界灵气固然丰沛，但其本质似乎依旧残缺不全，带着一种难以言喻的虚假和扭曲，让他体内自行圆满运转的真元产生了一丝本能的排斥与惕然。

他立刻传音提醒三位弟子：“灵气有异，固守本心，勿要引其入体，仅作环境参照即可。”

三人凛然，各自运功，将问道真元稳固于气海经脉之中，隔绝了外界灵气的自然渗入。

举目四望，他们身处一片荒郊野岭，脚下是坚硬的褐色土地，远处山峦起伏，植被茂盛，许多树木的叶片都闪烁着淡淡的灵光，显非凡品。天空高远，呈现出一种深邃的湛蓝色，云气舒卷间，似有灵禽身影掠过。

“此地灵气浓度，确实远超下界。”白言冰感应四周，低声道，“但……总觉有些浮夸，不如我真元沉凝踏实。”

陈千羽微微点头，她对于生机气息最为敏感，此刻却能察觉到这浓郁灵气背景下，自然万物生长虽旺，却隐隐有种被催熟的痕迹，少了下界那种历经风雨、自然而然的生命韧性。

沈瑶则更关注环境细节，她指向东北方向：“那边有频繁的灵气扰动痕迹，还有隐约的喧嚣声，似有人烟聚集。”

叶牧谦神识悄然蔓延，片刻后收回：“是一处规模不小的坊市。我们过去，小心行事，先探听消息。”

四人并未驾驭遁光，以免过于招摇，只是施展身法，一步数丈，向着坊市方向掠去。沿途偶尔遇到低阶修士或凡人模样的行商，皆行色匆匆，对叶牧谦四人也只是好奇打量几眼，并未过多关注。他们的服饰在灵界看来或许有些土气，但气息收敛得极好，像是偏僻乡里前来采买物品的低级修士，倒也不显突兀。

约莫半个时辰后，一座依山傍水而建的庞大坊市出现在眼前。坊市外围有简易的栅栏和瞭望塔，入口处有修士值守，但盘查并不严格，只需缴纳一枚下品灵石即可入内。坊市内部街道纵横，店铺林立，旌旗招展，叫卖声、讨价还价声、炼器锻打声、修士交谈声混杂在一起，热闹非凡。人流如织，其中修士比例极高，从灵枢境到真符境皆可见，偶尔甚至能感受到法相境修士隐晦的气息掠过。

坊市入口处的石碑上，镌刻着四个古朴大字——“流云坊市”。旁边还有小字标注：东域·郑氏辖下。

“东域……郑氏。”叶牧谦想起徐峰等人供述的信息。

虽然他现在已经笃定这方天地就是按照他的设定集演化出来的，但具体的社会形态，在他的设定集中却没有具体描述。

于是这些东西便是自行演化出来的了——甚至连同那“灵枢径”。

“灵界五域之一，郑氏，以炼器、经商著称的超级势力。看来这流云坊市，便是郑氏门下诸多产业之一了。”

步入坊市，喧嚣更甚。街道两旁店铺售卖之物琳琅满目：各种品阶的灵石、符箓、丹药、法器、灵材、功法玉简，乃至驯化的灵兽、奇花异草，应有尽有。许多货物散发着强烈的灵气波动，品相明显优于下界所见。交易媒介主要是灵石，亦有以物易物者。

师徒四人融入人流，一边闲逛，一边耳听八方，收集着各种信息。他们进入了几家茶馆、酒肆，要了些本地特色的灵茶点心，坐在角落静静聆听。

综合各类零碎信息，徐峰等人关于灵界的基本描述得到了验证。此界修士确实将内域与荒域统称为“下界”，视为贫瘠落后的蛮荒之地。而“天界”的存在也并非秘密，许多修士交谈时提及，语气中充满向往。

“听说了吗？西疆的那个老和尚，三百年前踏入天轮境的那个，去年成功通过‘天灵走廊’，飞升天界了！”

“啧啧，天轮境啊……我等不知要苦修多少岁月。听说天界灵气如雨，大道法则显化，才是真正的修行圣地。”

“五洲皆有天灵走廊，只要修为达到天轮境，便可尝试穿越。一旦成功，便能脱离此界束缚，前往更高层次的世界追寻大道。故而灵界那些顶尖强者，到了天轮境后，大多都会选择前往天界。”

“唉，就是这灵界，对修炼灵枢径的道友压制太大了。听说下界最高只能到法相境，再往上，天地灵气浓度就不支持了。咱们灵界好点，但最高也就到天轮境。”

“下界”！

原来这称呼在这里呢！

叶牧谦不禁微笑着点头。

他们也注意到，坊市中几乎所有的修士。修炼的都是灵枢径，问道途等其他法门，在此界似乎还是闻所未闻。

\vspace{2em}

打探得差不多了，叶牧谦决定先寻个落脚点。长期住客栈不仅花费高，也过于惹眼。他们开始在坊市外围区域留意是否有房产出售的信息。

运气不错，很快便在西侧靠近流云山脚的一片相对清净的区域，找到了一位正急得团团转、欲出售自家小院的修士。那修士是个真符境，似乎急需灵石周转，见叶牧谦几人气息沉凝（虽看不出具体境界，但感觉不好惹），且有意购买，立刻热情介绍。

小院不大，但颇为雅致，有正房三间，厢房两间，还有个小小的前院和后园，布置了简单的聚灵、防护阵法。位置胜在安静，且靠近山脚，若有意外也便于退入山中。

经过一番讨价还价，叶牧谦用一笔相对于坊市均价偏低的灵石（拜徐峰等人的“遗产”所赐），顺利购得了这小院的地契和房契。那修士拿到灵石后便匆匆离去，仿佛卸下了千斤重担。

有了安身之所，四人心中稍定。接下来几日，他们便以这流云山脚小院为据点，更深入地了解灵界，尤其是修行资源方面。

这一日晚饭后，叶牧谦将三位弟子召集到院中。

月色如水，流云山影朦胧。叶牧谦缓缓开口：“我等来此灵界，非为隐居。问道途之根本，在于‘实事求是、调查研究’。闭门造车，难明大道真意，亦难知此界虚实。”

他目光扫过三人：“一直聚在一起，所见所闻有限，易成井蛙之见。我意，尔等可各自出门游历，探寻机缘，体悟此界风土人情、修行百态。范围不必局限于这一隅，可往更远处去。每月十五，回此小院相聚一次，交流见闻心得即可。”

白言冰三人闻言，先是一怔，随即眼中都流露出跃跃欲试之色。他们明白师尊的深意。灵界浩瀚，道途万千，聚在一起固然安全，却也限制了视野与各自的缘法。分头行动，方能更全面地了解这个世界，也能在独立应对中更快成长。

叶牧谦又叮嘱了一些联络方式、应急手段以及需要重点留意的信息方向，例如各域势力分布、修行资源集散地、可能存在的上古遗迹或秘境传闻、以及关于“天灵走廊”和更高境界的更多细节等。

翌日清晨，师徒四人于小院门前分别。没有太多伤感，唯有对前路的期待与互道珍重。

叶牧谦站在院门前，望着弟子们远去的方向，直至身影消失。他转身回到院中，关上院门。小院顿时安静下来，只剩下风吹树叶的沙沙声。

他缓步走到院中石凳坐下，为自己斟了一杯从坊市买来的、味道略显浮夸的灵茶。举杯轻抿，目光投向流云坊市上空那终年不散的、由无数阵法灵光与修士遁光交织而成的淡淡光晕，以及光晕之外，那更高、更远、似乎隐藏着更多秘密的灵界苍穹。

“灵界……”他低声自语，“确是一片新天地。”

只是这新，是机遇，还是更大的樊笼？

\vspace{2em}

家里总归还是得有人坐镇，于是叶牧谦便成天窝在流云坊市了。他重点考察了坊市中售卖的法器，流云坊市作为郑氏旗下产业，法器铺子众多，从低阶制式武器到高阶精品皆有陈列。

然而，仔细查看之后，叶牧谦却微微摇头。

“用料确实扎实，许多灵材是下界难寻之物。”他自言自语地点评道，“但这锻造工艺、阵法铭刻……却显得颇为粗糙，匠气太重，灵性不足。”

他甚至发现，某些标价不菲的法阶法器，内部阵法衔接处竟有细微的灵力滞涩点，长期使用恐生隐患。相比之下，当初从徐峰、周锐身上缴获的那几件灵界制式法器，工艺反而显得更精良一些。看来，徐峰等人所在的界门司，配备的可能是更上乘的货色。这流云坊市面向大众，流通的或许只是大路货，甚至不乏次品。

这一发现，让叶牧谦心中有了计较。他专门去一些炼器铺子后院，或是坊市角落的垃圾回收处，以极低的价格，收购了一批被认定为“废料”的边角料和失败品。

回到小院，他闭门不出，专门捣鼓那些废料。偶然回来的白言冰等人好奇观望，只见师尊以真元为火，以神识为锤，不断熔炼、提纯那些看似无用的碎块。最终，他从一堆灰扑扑的晶体残渣中，提炼出了十几颗米粒大小、内部仿佛有雾气流转的银色晶粒。

“空灵晶……”叶牧谦指尖捻起一粒，感受着其中蕴含的微弱但奇特的空间属性波动，“内域未有出产，倒是炼制涉及空间、隐匿类法器的好材料。”

他又取出一块在坊市购买的精钢，以特殊手法将其锤炼得薄如蝉翼却坚韧异常。随后，他将提炼出的空灵晶粉末，以独特的真元导引之法，一点点融入精钢薄片之中，并在其上刻画下繁复而契合的微型阵法。这些阵法能以自身真元为引，激发空灵晶特性，兼具锋锐、破灵、光影扭曲及短距瞬闪等多重效果。

数日之后，一柄长约一尺二寸、造型流畅简约、通体泛着淡淡银灰色流光的短刀，置于案上。刀身看似朴实，但若凝神细观，会发现其边缘光线微微扭曲，仿佛随时会融入周围环境。

叶牧谦拿起短刀，轻轻一挥，并无引人警觉的破空之声，但对面墙上悬挂的一张测试用牛皮却无声无息地裂开一道平滑的切口。他注入一丝真元，短刀银光微闪，竟瞬间从手中消失，下一刻出现在三丈外的木桩上。

“便叫‘真影刃’吧。”叶牧谦将短刀收回，评价道，“用料所限，仅能算法阶，但胜在特性独特，与瑶儿的千幻流光之术颇为契合。”

等沈瑶回来，他便将真影刃递给沈瑶。沈瑶接过，稍一感应，便觉此刀仿佛是为自己量身打造，真元注入畅通无阻，那光影扭曲与短距瞬闪之能，与她的幻术简直是天作之合。她眼中绽放出惊喜的光芒，爱不释手，连忙行礼：“多谢师尊！弟子定善加运用！”

叶牧谦微微一笑，心中却想：这流云坊市毕竟只是东域一隅、郑氏旗下的一座普通坊市，所见炼器水准或许不能代表灵界真正的高层技艺。灵界广袤，能人辈出，深处那些传承悠久的炼器世家、宗师，未必没有比自己更强的存在。

\vspace{2em}

这一日，叶牧谦信步穿行于坊市西侧一片相对冷清的巷道。

这里摊位稀疏，售卖的多是些不入流的低阶材料或是来路不明的小玩意儿，光顾者也多是些修为低下、囊中羞涩的底层修士。叶牧谦本来还是来这里采购废料的，这边的废料价格低得近乎白送。

忽而，他的脚步微微一顿，目光被角落里一个不起眼的小摊吸引了过去。

摊主是个看起来二十出头的年轻人，修为仅灵枢境，穿着一身洗得发白的粗布衣衫，面容清秀却带着几分长期劳作的风霜之色，眼神倒是明亮有神。

他的摊位不大，上面摆放的东西也与周遭格格不入——并非寻常的法器铺子那样，陈列着刀剑盾甲等杀伐护身之物，甚至大多都称不上传统意义上的“有用”。

那是一些极其精巧的小玩意儿：一只由无数细如发丝的金属簧片和微型齿轮构成的木雀，注入微弱的灵力后，竟能扑扇着翅膀在摊位上盘旋数圈，发出悦耳的啁啾声；一座微缩的亭台楼阁模型，以半透明的晶石薄片拼接，内部有灵光线路流转，模拟出日升月落、窗扉开合的景象；还有几个类似万花筒但结构复杂十倍不止的铜管，透过一端看去，能见到不断变幻、组合成各种繁复图案的灵光……

每一件都做工细腻，设计充满了天马行空的巧思，将有限的灵力通过极其精妙的结构转化为生动的动态效果，与其说是法器，不如说是凝聚了创造者心血与乐趣的灵巧机关。

叶牧谦眼中闪过一丝讶异。在灵界这个普遍追求威力、品阶、实用性的炼器氛围中，竟有人沉心于此等“奇技淫巧”，而且其设计展现出的结构思维能力与对灵力微操的想象力，远超其灵枢境的修为水平。

他驻足摊前，俯身仔细观看那只木雀，手指虚触其精密的结构，神识细细感知其内部灵力的流转路径与机械联动的精妙。

“这位前辈，可是对这小玩意儿感兴趣？”年轻摊主见叶牧谦气度不凡，虽看不出深浅，但也不敢怠慢，主动开口，声音带着些许忐忑，“都是我自己瞎琢磨做出来的，上不了台面，就是图个有趣……您要是喜欢，价格好商量。”他自称柳成。

叶牧谦拿起那个亭台模型，注入一丝真元感应，随口问道：“构思巧妙，尤其是这光影联动与机械时序的配合，颇见心思。你学炼器多久了？师承何人？”

柳成挠了挠头，有些不好意思：“回前辈，晚辈没什么正经师承，就是从小喜欢鼓捣这些机巧物件。家里原是坊市里帮炼器铺子处理边角料的，耳濡目染偷学了些皮毛，后来自己瞎琢磨……让前辈见笑了。”他话语朴实，提到自己喜爱之物时，眼睛却格外明亮。

正交谈间，巷道口传来一阵脚步声，一个衣着光鲜、腰间佩玉的年轻公子哥摇着一把折扇走了过来。这公子哥修为在真符境巅峰，在流云坊市年轻一辈中算是不错，但脸上并无寻常世家子弟的倨傲之气，反而带着几分跳脱与随意。他一眼看到柳成，眼睛一亮，快步走了过来。

“柳成！可算找到你了，上次说好的东西呢？”公子哥语气熟稔，显然与柳成相熟。

柳成见到他，脸上露出笑容，又略带歉意地对叶牧谦道：“前辈稍待，熟客来了，失陪一下。”

说完，便与那公子哥走到摊位后面用布帘简单隔出的小小后屋去了。

叶牧谦本欲离开，但神识敏锐地捕捉到两人进入后屋后，有轻微的、仿佛在翻找什么的窸窣声。他心中微动，一丝若有若无的神念悄然蔓延过去……

后屋内空间狭窄，堆满各种工具和零碎材料。只见柳成从一个锁着的小木箱里，小心翼翼地取出一个巴掌大小、外观颇为粗糙的铜制圆镜，镜框雕刻简陋，镜面也不是通常的光滑水晶或白银，而是一种略显浑浊的淡黄色晶石薄片。

“郑兄，你看看，按你要求改的，留影时间能延长到差不多一炷香了，但清晰度还是老样子，动静结合的地方有点模糊……”柳成压低声音道。

那被称作郑兄的公子哥接过铜镜，翻来覆去看了看，又注入一丝灵力激活。只见浑浊的镜面上微微亮起，浮现出一些晃动模糊的影像片段。公子哥看得仔细，时而点头，时而皱眉。叶牧谦的神念“看”到那镜中影像，内容竟是一些不堪入目的黄色画面，不知这柳成是从何处录得这些内容。

然而，更让叶牧谦注意的是这留影镜本身。它的炼制手法确实粗糙，用料也低廉，但关键在于其原理：

一种能够记录并重现动态影像的、结构相对简单、似乎具备批量复制可能性的技术！

在灵界这个信息传递仍主要依靠玉简、符信或口耳相传的世界，这种廉价的、可复制的动态影像记录技术，其潜在价值，绝不止于记录闺房之乐。无论是记录功法演示、勘探地形、传递情报，还是其他更广阔的领域，都蕴含着巨大的可能性。

这柳成，在炼器一途上，恐怕拥有着连他自己都未曾意识到的、独特的天赋：一种将复杂功能通过巧妙简化结构来实现的“工程化”思维，以及不拘一格的想象力。

此时，那郑公子已经开始与柳成讨价还价：“柳成啊，时间长了是好事，但这画质……啧，看个轮廓动静还行，细节实在糊啊。这价钱……你看能不能再让让？我知道你搞这些不容易，那些透光的水影晶太贵，你用这种廉价的石头也是没办法，但效果打了折扣嘛……”

柳成苦着脸：“郑兄，真不能再低了。为了延长留影，我改了三版阵纹，报废的材料都快抵上卖价了。这昏黄石虽然成像差，但感光蓄灵的特性正好合用，换别的成本更高……”

两人你来我往，讨价还价，语气熟稔自然。那郑公子哥全然没有高高在上的姿态，反而颇有市井小民之间“你懂我需求、我懂你难处”的实在味道，一边挑着毛病，一边又透露出理解；最终，以一个双方都勉强接受的价格成交。

郑公子揣起留影镜，又跟柳成嘀咕了几句“下次搞点新花样”“清晰度再想想办法”之类的话，便心满意足地离开了。

叶牧谦收回神念，面色平静无波，心中却已暗暗记下了这两个人——拥有奇特炼器天赋却困于底层的柳成，以及这位看似不拘小节、甚至有些“荒唐”需求，但待人接物却颇为实在、没有世家子弟常见倨傲之气的郑公子。

\vspace{2em}

接下来几日，叶牧谦成了柳成摊位的常客。他不再只是远远观察，而是真的会驻足挑选一两个有趣的小机关，付钱买下，并与柳成闲聊几句。他问的问题往往切中柳成那些作品构思的关键或难点，偶尔轻描淡写的一两句提点，却常让苦思许久的柳成茅塞顿开，越发觉得这位“叶前辈”深不可测。

一来二去，两人便熟络起来。

连带着，那位经常来找柳成定制些稀奇古怪玩意儿的郑公子，也混了个脸熟。

叶牧谦也得知了这公子哥的名字——郑天明，确实属于郑氏，但只是家族中较为偏远、不受重视的一支旁系子弟。或许正因为出身并非核心，资源有限，前途也相对暗淡，郑天明身上反而没有那些嫡系子弟常见的骄横之气。

郑天明性格憨厚直率，甚至有些过于实在，对炼器杂学感兴趣，却因天赋所限难有成就，但好在灵石充足，于是便成了柳成这些“不务正业”发明的欣赏者和赞助者。他来找柳成，有时是为了些像留影镜这样的“私人趣味”，有时则是单纯觉得某个小机关好玩，买回去解闷。他对柳成的态度，更像是对待一个有真本事的匠人朋友，而非可以随意驱使的摊贩。

观察日久，叶牧谦心中渐起波澜。他看着柳成在谈及精妙结构时眼中绽放的光彩，看着他用极其有限的材料、凭借巧思与耐心解决一个又一个难题，那种专注与热爱，与他当年在荒域初遇白言冰、在内域发现沈瑶时感受到的某种特质，隐隐相通。那是一种对“道”的纯粹追求与潜能，只是尚未找到正确的路径与引导。

叶牧谦隐隐有所感应，这位看似不起眼、挣扎在灵界底层的年轻摊主柳成，或许便是他在冥冥之中所感应到的，第四个，也是最后一个弟子缘法所在。

那种独特的、化繁为简的炼器天赋，那份不被世俗“有用无用”标准束缚的创造力，若能引入问道途，以内求真元驱动，以“实事求是”之心探究万物机理，未来在器之一道上的成就，恐怕不可限量。

然而，收徒之事，讲究缘法，更需水到渠成。叶牧谦早已不是当初那个为了“反哺”而硬收徒弟的迷茫之人。他深知，强求而来的师徒名分，远不如心灵契合、自愿追随来得稳固与有益。

柳成如今虽困顿，但心志未堕，且有郑天明这样一位不算坏的朋友，在灵界也勉强能立足。自己若贸然提出收徒，反而可能惊扰了他，或让他因处境所迫而做出并非本心的选择。

不急。玉不琢，不成器。璞玉已现，待时而动便可。

\vspace{2em}

这一日，坊市中的气氛却隐隐与往日有些不同。

一种沉在底下的、压抑着的躁动，与浮在表面的、克制的期待交织在一起，仿佛平静湖面下暗涌的潜流。空气中弥漫的不仅是灵材、丹药和食物的混合气味，更添了几分紧绷与算计。

许多相熟的修士凑在一起交谈时，会不自觉地压低声音，身体前倾，眼神却不受控制地闪烁着攫取与兴奋的光，如同嗅到了血腥气的鬣狗。

茶馆酒肆里，跑堂的小二端着茶盘穿梭的频率都似乎快了几分，而那些原本议论着东家长西家短、或是某处秘境传闻的闲谈，话题陡然集中，关于“小玄界”的议论如野火般蔓延开来。

“听真切了没？新的小玄界要开启了！入口就在咱们东域，据说那空间裂隙的落点，离咱们郑氏领地西侧的枯骨荒原不远！”一个尖嘴猴腮的修士呷了口粗茶，声音压得极低，却掩不住那份恨不得全坊市都来打听的炫耀意味。

他对面坐着的疤脸汉子眼皮一掀，瓮声瓮气道：“当真？按灵界惯例，这种无主的、新诞生的小玄界开启，那可是五域震动的大事！只要消息灵通些的，哪方势力不想来分一杯羹？看看能不能捞点上古遗泽，或是撞破天运，寻到点什么失传的传承碎片。”

“可不是嘛！”尖嘴修士一拍大腿，唾沫星子差点溅到茶碗里，“咱们郑氏这次可是近水楼台！以家族那炼器、经商起家的精明性子，还有对自家地盘的控制力，这进入小玄界的资格，还不得标上天价？我估摸着，散修和小门小派，怕是连门都摸不着，也就是那些中型势力，或者背景深厚的独行客，才有资格掏空家底去争上一争。”

叶牧谦独自坐在茶馆靠窗的角落，面前一盏清茶早已凉透。他静静听着这些毫不掩饰的议论，面上无波无澜，心中却已了然。

看来，这灵界对此类事件的反应模式，与他在内域天枢城所经历的，并无本质不同。无非是规模更大，牵扯的利益更广，参与博弈的棋子也更加强大罢了。

他想起了内域那次小玄界开启，天师府主持下，双方约法三章，仅限年轻弟子进入，生死各凭机缘。而在这里，在郑氏掌控的地盘上，规则显然更加赤裸，更加以资本和实力为门槛。

\vspace{2em}

果然，没过两日，郑氏官方便发布了措辞严谨、盖有家族法印的正式公告，确认了小玄界入口位于其势力范围边缘一处荒谷的地点的消息；并同时宣布，将公开出售有限数量的“探索资格令”。

公告一出，流云坊市便沸腾起来。那资格令的价格虽未明示，但已有小道消息传出，其数目之巨，足以让一个中产修士伤筋动骨，令绝大多数散修望而却步，只能沦为看客。然而，这骇人的价码并未吓退野心家，反而像最上等的饵料，引得东域乃至从中州、北地闻风而来的修士、商队、各方代表蠢蠢欲动，纷纷汇聚而来。

流云坊市作为郑氏在东域重要的商业节点之一，顿时人满为患。

客栈房价飞涨，酒楼座无虚席，各处情报贩子的生意火爆异常。许多外来修士涌入，目光灼灼地打探着关于荒谷地形、可能出现的宝物类型、以及郑氏核心队伍动向等一切细节，同时疯狂筹措灵石，或是变卖身上的法宝材料。相应地，坊市中疗伤丹药、防护符箓、逃遁法器、以及各类一次性攻击宝物等探索必备物资的价格，也开始一路水涨船高，弥漫着一股大战将至般的紧张与狂热。

\vspace{2em}

这一日，叶牧谦如往常般来到柳成的摊位前。柳成正全神贯注地调试着一个新做出的灵巧星象仪，巴掌大小的浑圆球体上，嵌着数十枚微缩的、散发着柔和星辉的灵晶，随着他指尖灵力的细微引导，那些“星辰”开始沿着复杂而玄妙的轨迹缓缓运转，模拟出一小片简易的星图演化。他眉头微蹙，沉浸其中，连叶牧谦走近都未曾察觉。

就在这时，一阵熟悉的、风风火火的脚步声由远及近，郑天明那魁梧的身影像一堵墙似的挤了过来，脸上带着掩饰不住、几乎要满溢出来的兴奋红光。

“柳成！柳成！机会来了！天大的机会！”郑天明一把拉住柳成拿着星象仪的手腕，力道之大让柳成一个趔趄，星图运转顿时紊乱。

郑天明这才反应过来，连忙松手，却依旧压低了声音，那激动之意却怎么都压不住，“家族内部开始分配探索资格了！嘿嘿，像我这样的旁支子弟，虽然捞不着那些给核心天才、重点培养对象的核心名额，但按惯例，也能分到一个额外的陪同名额！我想好了，这名额不给别人，就给你！咱哥俩一起进去，好好碰碰运气！说不定就能捡到啥宝贝，从此翻身！”

他虽是郑氏旁支，血脉疏远，平日不受重视，但毕竟顶着郑氏的名头，在这种带有家族福利性质、安抚旁系人心的名额分配上，倒也能勉强沾上一点边。

柳成闻言，先是一喜，眼中骤然迸发出光彩。

小玄界，那可是传说中的地方！

但这点光芒瞬息即逝，现实的冰冷迅速浇灭了他的幻想。

他脸色以肉眼可见的速度垮了下来，头摇得像狂风中的拨浪鼓，连带着手中星象仪的微光都慌乱地闪烁起来：“郑兄，使不得，万万使不得！你快醒醒！我这点修为，区区灵枢境，进去能干嘛？给里面游荡的古代妖兽当开胃小菜？给那些失了效却更危险的上古机关阵法制成的人傀当新的零件？还是给其他红了眼、杀人不眨眼的修士白白送人头？那地方是你们家族核心精英和外界各方高手博弈厮杀的战场，我进去纯属找死，还得连累你分心照顾！”

郑天明急了，他不能理解柳成这份“清醒”，只觉得朋友太过畏缩：“怕什么？不是还有我吗？我罩着你！咱们有点自知之明，不往那些标记着高危的核心区域凑，就在外围，最最边缘的地方转转！捡点前人看不上的边角料，挖几株年份浅的灵草也好啊！就当开开眼界，总比在这坊市里天天对着这些铁疙瘩强吧？”

柳成攥紧了手中的星象仪，指节有些发白，他还是坚决地摇头，语气斩钉截铁：“郑兄，你的好意，我心领了，真的。但我柳成知道自己几斤几两。这不止是为我自己，也是为你。你带我进去，就是你的累赘，遇到危险你救是不救？不救，你良心不安；救，可能就把你自己也搭进去。这名额珍贵，哪怕是个陪同名额，也是多少子弟眼巴巴望着的。你还是……另请高明吧。”

他态度坚决如铁，任凭郑天明如何晓之以情、动之以“利”，甚至急得抓耳挠腮，也无济于事。

郑天明像只被扎破的兽皮气囊，瞬间泄了气，高大身躯都佝偻了几分，沮丧地松开原本按在摊位上的手。他环顾四周，目光无意识地扫过摊位前气度沉静、正若有所思地打量着那暂时停摆的星象仪的叶牧谦，黯淡的眼睛忽然像是被火石擦亮，再次迸发出光来。

这位神秘的“叶前辈”，虽然修为深不可测，让人看不透根脚，但几次接触下来，待人接物始终平和淡然，毫无高阶修士常见的傲慢与疏离，感觉颇为可靠。而且，他能随口指点柳成炼器关窍，往往一针见血，其自身见识与修为定然都极为不凡。一个念头如同野草般在他心里疯长。

他凑到叶牧谦身边，搓着手，脸上堆起平日里求家族炼器师帮忙时那种特有的、混合着讨好与期待的憨厚笑容：“叶前辈，您看……柳成这小子，胆子比耗子还小，好说歹说就是不敢去。晚辈这个名额，空着也是浪费，家族规定又不能转让买卖……不知前辈……或者前辈门下是否有高足，近期有空，可有兴趣一同前往？彼此也好有个照应。”

他话说得小心翼翼，带着试探，眼神里却满是灼热的期待，仿佛抓住了最后一根稻草。

叶牧谦缓缓放下手中那精巧却此刻显得无比脆弱的星象仪，目光平静地看向郑天明。

又是小玄界……他正欲更深入地了解灵界的方方面面，尤其是其资源分布、势力格局以及可能存在的上古隐秘，这样一个小玄界，倒是个绝佳的观察窗口。不过，他自己亲自前往，目标太大，容易引来不必要的关注，且流云山巅那处清静小院也需有人看顾。

他沉吟片刻，心中已有人选，淡然道：“我确有一名弟子，近日游历归来，修为尚可，心性也算沉稳，或可陪你同去历练一番。”

郑天明大喜过望，几乎要跳起来，连忙躬身道谢：“太好了！多谢叶前辈成全！不知是哪位师兄？晚辈何时可以拜见？”

“他叫白言冰，不日即归。届时，我让他直接去寻你便可。”叶牧谦语气依旧平淡。

郑天明欢天喜地，仿佛那名额已经换来了绝世宝藏，又连声道谢了几次，这才脚步轻快地离去，开始火急火燎地为进入小玄界做准备——尽管他拥有的情报和资源，恐怕也有限得很。柳成看着他那仿佛捡了天大便宜的背影消失在人群里，转过头对叶牧谦露出一个混合着无奈与担忧的苦笑：“让前辈见笑了。郑兄人……是极好的，热心肠，没什么坏心眼。就是有时候，想法过于简单直率，把事情想得太过轻易。小玄界那地方，步步杀机，人心鬼蜮，岂是那么好相与的。”

叶牧谦微微颔首，目光似乎透过坊市喧嚣的穹顶，看到了更远处：“无妨。让言冰去见识一番灵界修士真正的争夺之态，于他而言，亦是一种历练。” 他语气平静，却带着一种笃定。他对白言冰有信心，那孩子心志坚毅，历经磨练，早已不是当年荒域中那个只知挥剑的少年。

\vspace{2em}

数日后，一道略带风霜却更显挺拔的身影回到了流云山小院，正是自游历归来的白言冰。他气息内敛，比离开时更为凝练扎实，眉宇间多了几分冷硬的棱角，眼底沉淀着旅途见闻带来的思索。

听闻师尊安排，他并无半分异议，沉稳抱拳领命，眼中反而掠过一丝跃跃欲试的锐芒。叶牧谦将小玄界开启的前后经过、郑氏的态度、以及郑天明此人热心却可能失于莽撞的性子大致交代一番，末了，重点叮嘱道：“此行以探查与历练为主，安全为上。玄界之内，机缘与风险并存，人心更比机关险恶。多看，多思，慎行。”

白言冰肃然道：“弟子明白，定谨记师尊教诲。”

叶牧谦又交予了白言冰一些法宝，先前捆过柳媚的那条绳子也在内，“以备不时之需”。

翌日，白言冰便按师尊给出的地址，在坊市一角寻到了正对着一堆乱七八糟符箓和矿锭发愁的郑天明。郑天明抬头，见来人身材挺拔如松，面容刚正，目光清澈而沉静，周身气息含而不露，行走间自有一股沉稳气度，一看便知绝非寻常散修或小门派弟子可比，心中顿时大定，欢喜之情溢于言表。

他连忙将一枚入手微沉、镌刻着郑氏家族独特云纹徽记的黑色金属令牌交给白言冰，口中连道“有劳白兄”，又约定好三日后清晨，在荒谷入口处集合。

\vspace{2em}

三日转瞬即过。

\vspace{2em}

往日飞沙走石、鬼哭狼嚎的荒僻山谷入口，此刻景象已然天翻地覆。人声鼎沸，灵力波动混杂冲天，足足有数百名修士聚集于此。他们依附着不同的旗帜、穿着制式各异的服饰，或三五成群，或数十人一队，泾渭分明地占据着谷口外的各处有利位置，彼此间眼神交汇时充满警惕与审视，无形的暗流在人群中涌动，压抑的敌意几乎凝成实质。

空中偶尔有刺耳的破空声或绚丽的流光划过，拖出长长的灵尾，那是某些势力的大人物驾驭的飞行法宝或遁光，它们根本不屑在此停留，直接无视下方人群，堂而皇之地落入被淡淡迷雾笼罩的谷内深处——那是提前数日便已进入的郑氏核心精英团队，以及少数与郑氏交情深厚、或是付出了外界难以想象的巨大代价，从而换取到优先进入权与更详尽情报的其他大域强势代表。

郑天明带着白言冰赶到这纷乱如集市般的谷口时，即便早有心理准备，也被这凝重而充满竞争火药味的气氛压得缩了缩脖子。

他下意识地靠近白言冰半步，低声快速说道：“白兄，看到了吧？咱们啊，还有这边绝大多数人，说白了就是来陪跑的。家族里头那些真正的天才，被宝贝一样供起来的重点苗子，还有依附于家族的强力客卿，前几天就已经跟着长老们进去了，这会儿估计都在探索核心区域，争夺最大的那份机缘了。像咱们这种旁支子弟，分到的不过是些边角料名额，还有这些外面来的、花了天价买路钱的修士，就是最后一批被放进去的。能捡到点前面人马看不上的残羹冷炙，或者运气逆天在不毛之地发现点漏网之鱼，就算烧高香了。”

他语气里有一种早已认命的平静，并无多少愤懑不平，反而透着一种在等级森严大家族中长大的、对规则天然的“本该如此”的坦然。

白言冰默默点头，凌厉的目光早已扫过周围的人群。修士修为参差不齐，以真符境、法相境居多，偶尔能感受到一两道更为隐晦强大的气息，应是法相境中的佼佼者或是压制了修为的厉害角色。不少人身上煞气隐隐，眼中带着常年刀头舔血的冷漠与贪婪，显然都绝非善类。

他暗自将警惕提到最高，丹田内真元默默流转，腰间那柄看似寻常的、叶牧谦亲手为他打造的佩剑，似乎也感受到主人的心绪，传来一丝微不可察的清鸣。

\vspace{2em}

约定的时辰一到，谷口处那由数十名身着郑氏制式灵甲、面色冷峻修士把守的临时关卡缓缓“开放”——实则只是撤去了一层警示性的灵光屏障。持有资格令的修士排成长队，依次上前，由郑氏修士验明令牌真伪并登记神识印记后，便被放入那迷雾缭绕、不知通向何方的谷中。过程沉默而迅速，带着一种不容置疑的秩序感。

踏入谷口迷雾的瞬间，仿佛穿过了一层冰凉的水膜，外界的喧嚣嘈杂骤然减弱、变形，取而代之的是一种奇异的寂静和空间上的轻微扭曲感。

郑天明深吸一口气，彻底收敛了平日跳脱的神色，变得严肃起来，甚至显得有些紧张。他迅速从怀中掏出一份边缘已经有些磨损的简陋皮质地图，摊开给白言冰看：“白兄，跟紧我！这鬼地方，家族发放的地图也语焉不详，只标了几个大概的危险区域和可能的方向。里面不光有地图上都没记载的天然陷阱、破碎的空间裂缝，还有那些上古遗留的、可能已经错乱的机关阵法，更要命的……”他顿了顿，声音压得更低，“是那些同样进来、却可能在任何时候翻脸的其他家伙！咱们不求有功，但求无过，小心驶得万年船！”

地图上线条粗犷，只有寥寥几处标记，大片区域都是空白，或者打着问号。

这确是他能从家族得到的、最粗浅的共享信息了。

白言冰目光在地图上快速掠过，将其记在心中，然后望向迷雾深处那隐约显露出的、与外界荒谷截然不同的、布满嶙峋暗色巨石与低矮扭曲植物的怪异地貌，点了点头，握紧了剑柄。

踏过那道如水波般荡漾、却散发着不稳定空间波动的入口，一股混杂着铁锈、腐朽与淡淡腥甜的沉闷气息便扑面而来，取代了外界的清新灵气。白言冰与郑天明对视一眼，都在对方眸中看到了一抹凝重。

白言冰在内域时曾进入过青鸾秘境。在那里，那份源自上古神鸟的生机道韵始终萦绕，草木葱茏，灵机盎然，连破败的殿宇都因那道韵而保持着一种神圣的衰颓感。

而这里，天空是一种永恒的、令人压抑的铅灰色，仿佛被厚厚的尘霾与尚未散尽的硝烟永久浸染，不见日月，只有朦胧胧的光晕勉强勾勒出事物的轮廓。大地是焦黑的，布满了龟裂的纹路，如同干涸了万年的河床，踩上去发出“咔嚓”的脆响，那是无数细小碎石与某种硬化的灰烬。

举目望去，没有一丝绿色。扭曲的、早已石化或炭化的树木残骸以狰狞的姿态指向灰空，像是无数死者最后的控诉。远处，影影绰绰有着巨大的轮廓，像是倒塌的城楼、折断的巨柱、或是某种庞大法器或妖兽的残骸，沉默地矗立在灰雾中，如同巨兽的骸骨。

空气中弥漫的，除了那股陈年的腐朽与血腥混合气，还有一种极为淡薄、却无孔不入的煞气，那是一种纯粹的、属于死亡与混乱的冰冷意念，微微刺激着修士的神魂，让人下意识地绷紧神经，心绪难宁。

“果然……是一处古战场。”郑天明压低声音，他虽不受待见，但家族的底蕴让他见识远超常人，此刻脸色也有些发白，“看这煞气凝结的程度，以及天地规则的隐隐扭曲……当年在此陨落的，恐怕不止是一两个道种境修士，甚至有天轮乃至圣魄！”

他们的道韵、怨念、残留的杀意，与破碎的空间、崩坏的地脉混杂在一起，经年累月，才形成了如此凶恶的小玄界。

白言冰默默点头，手不自觉地按在了腰间佩剑的剑柄上。剑身传来熟悉的冰凉触感，让他心中稍定。他回想起青鸾秘境中，众人虽也有争斗，但至少核心处尚有传承与机缘，有一线光明。而此处，目之所及，只有破败、死亡与无尽的凶险。生机盎然与死寂破败，形成了何其鲜明的对比。

两人不敢大意，将神识小心翼翼地探出，却又不敢延伸太远，唯恐触发什么未知的存在，或引起暗处其他修士的注意。他们选择了远离几处看似曾有大规模战斗痕迹、残存能量波动也最混乱的区域，向着一片相对平坦、只有零星低矮石丘的方向探索。

然而，这份“相对平坦”本身就是最大的假象与陷阱。

前行不过百丈，郑天明脚下踩中一块颜色略深的石板，石板看似与周围焦土融为一体，却在承重的瞬间无声地翻转！他反应极快，轻身提气，身体向后急仰，同时灵气在身下凝聚，狠狠轰向身旁一块凸起的岩石借力。

就在他身形将起未起之际，下方黑洞洞的坑口中，骤然暴射出数十道幽蓝的寒光——那是密密麻麻、倒插在坑底、淬着未知剧毒的金属尖刺，尖端闪烁着不祥的光泽，距离他的脚底不足半尺！

白言冰几乎在同一时间出手，剑未出鞘，一道凝实的剑气已然挥出，并非攻击，而是斜击在郑天明身侧地面，借反震之力助他彻底脱离陷坑范围。

郑天明落地，额角已见微汗，回头看向那瞬间恢复原状、只是石板颜色似乎更新了一些的陷坑，心有余悸。

“是古战场遗留的防护陷阱？还是后来者布置的？”白言冰问道，目光锐利地扫视四周。

“此地煞气侵蚀，很多阵法、机关的灵力脉络可能早已变异，或处于半激活的脆弱状态，稍加触动就会爆发。而后来者，更可能利用这些现成的危险。”

郑天明的话音刚落，侧前方一片看似无害的、贴着地面蔓延的灰褐色“枯藤”突然毫无征兆地暴起！它们快如毒蛇，柔韧似钢索，表面还生有密密麻麻的倒刺，如同标枪般直刺两人！

“小心！”白言冰低喝，长剑终于出鞘，一道清亮的剑光划破灰暗，精准地斩在最前方的几根枯藤上。然而，剑锋切入，像是砍在了浸油的牛皮上，带着滞涩感，且被斩断的枯藤切口处，竟然喷溅出几滴腥臭的暗绿色汁液，落在焦土上，发出“滋滋”的腐蚀声。

郑天明忍不住痛骂一句，身法灵动，灵气在身前凝成一片光幕，格开刺向自己的枯藤，身形急退。那些枯藤一击不中，竟似有生命般，蜿蜒扭动着缩回原地，重新伪装成毫无生气的死物，只是那片区域的灰褐颜色，似乎更深了一些。

他们更加谨慎，几乎每一步都要用神识或剑气先行试探。

然而，古战场的凶险远不止于此。

当白言冰的剑气扫过一片岩壁时，岩壁上几处不起眼的孔洞骤然喷出大股墨绿色的酸液，腐蚀性极强，连坚硬的焦土都能蚀出深坑；另一处，他们绕过半截插入地下的巨型兵刃残骸，残骸上模糊的符文却因外界的灵力扰动而微微一亮，随即从刃身上震射出数十支锈迹斑斑却淬着暗芒的弩箭，力道惊人，足以射穿真符境修士的护体灵光；还有一次，他们踏足一片相对空旷的区域，脚下却突然亮起复杂而残破的阵纹，炽热的火焰从阵纹中冲天而起，覆盖了方圆数丈，若非两人见机得快，提前感知到脚下微弱的灵力躁动而飞身后撤，恐怕已被火焰吞没。

这些机关，有些是古战场遗迹本身所有，因岁月侵蚀而变得极不稳定；有些，则明显带有“人工”痕迹——符文较新，触发机制刻意，显然是先他们一步进入的修士，或为防范他人，或为坑杀后来者所设。这种阴毒的心思，比环境本身的危险更让人心寒。

\vspace{2em}

就在他们疲于应付层出不穷的环境危机时，更大的威胁，来自同类的恶意，开始悄然显现。

探索第三日，两人在一处倾斜的巨岩下，发现了一小片罕见的阴魂菇。这种菌类只生长在极阴煞且有大量魂魄消散之地，是炼制一些特殊丹药、或修炼某些偏门魂道法术的珍贵材料。郑天明眼中刚露出一丝喜色，白言冰却猛地拉了他一把，向后急退。

几乎在他们退开的瞬间，三道凌厉的攻势从三个不同的方向袭来！

一道赤红的火蛇符箓轰在刚才郑天明站立之处，一道乌黑的梭形法器带着破空声钉入岩石，还有数枚细如牛毛的毒针无声无息地笼罩了他们原先所在的区域。

“嘿嘿，反应倒快。”一个阴恻恻的声音响起。从巨岩侧后方、一片风化严重的石笋林、以及头顶一块悬石上，各自跃出两三名修士，呈三角之势将白、郑二人隐隐围住。这些人衣着各异，显然不是来自同一势力，但眼中都闪烁着贪婪与凶光。看起来，他们是为了那丛阴魂菇，临时勾结在了一起。

“两位朋友，见者有份。留下储物袋，或者……留下命来。”为首一名脸上带着刀疤的汉子狞笑道，手中提着一把鬼头大刀，气息赫然是真符境。

郑天明脸色微沉，握紧了短刃。白言冰则面无表情，只是缓缓抬起了手中长剑，剑尖斜指地面，一股凛冽而沉静的气势缓缓升起。他并未说话，但那股经过无数次厮杀磨砺出的剑意，让围上来的几人脸色微微一变。

“点子扎手，一起上！”刀疤脸眼中凶光一闪，不再废话，率先扑上。其他几人也各持兵刃法器，围攻而来。

战斗在狭窄而危机四伏的古战场一隅爆发。白言冰的剑法沉稳精准，守得滴水不漏，偶尔一剑反击，便逼得对手手忙脚乱。郑天明身法诡谲，法术专攻要害，也牵制住了两人。然而对方毕竟人多，且实战经验丰富，配合虽然不算默契，但狠辣有余。

就在战斗胶着之际，那名躲在后面释放符箓的修士，眼中狡色一闪，并未攻击白言冰或郑天明，反而将一张爆裂符射向了那丛阴魂菇旁的一块看似松动的岩石！

“不好！”郑天明识破其企图，他是想制造混乱，甚至引发可能的连锁陷阱！

轰隆！岩石炸开，不仅阴魂菇被毁了大半，爆炸的冲击波果然触动了岩石下方某种早已脆弱的地脉节点。地面剧烈一震，数道炽热的熔岩火舌从炸开的裂缝中喷涌而出，逼得交战双方不得不仓皇后退。

趁此混乱，那名释放符箓的修士和另一人，竟毫不犹豫地转身就逃，扑向残存的几朵阴魂菇，捞起就跑，全然不顾还在与白、郑二人缠斗的同伴。

“妈的！混蛋！”刀疤脸气得大骂，心神微分。

白言冰岂会放过这等机会，剑光如电，瞬间在他肩头留下一道深可见骨的血痕。刀疤脸痛呼一声，再不敢恋战，虚晃一刀，与剩下那名同伴也狼狈逃窜，消失在嶙峋的石林之中。

郑天明惊出一身冷汗，看向白言冰的目光充满了震撼。方才白言冰那看似简单的一剑，无论是速度、精准还是其中蕴含的凝练力量，都远超他的理解。这位叶前辈的弟子，果然不凡！

战斗突兀开始，又仓促结束。现场只留下打斗的痕迹、喷射的熔岩、被毁的灵草，以及空气中淡淡的血腥味。郑天明喘着气，看着狼藉的现场和逃遁的敌人，苦笑道：“看见了吧，白兄。在这里，一株灵草，就足以让人拔刀相向，甚至不惜引发更大的危险拖所有人下水。刚才那家伙，若真引发大规模地火喷发，他自己也未必能逃掉。”

白言冰收剑入鞘，看着敌人消失的方向，眼神微冷。“他们只是更贪，且笃定自己能在危机关头逃脱，或拉别人垫背。” 他顿了顿，“我们尽量避开人群密集的方向，选择偏僻路径，但看来……偏僻之处，也未必安全。既是古战场，何处没有遗宝的诱惑？何处又没有伺机而动的猎人？”

两人简单处理了一下现场，确认没有其他隐藏陷阱，就把剩余的蘑菇泥（虽然不如蘑菇本身，但还算有点用）收集起来。然后，服下丹药略作调息，便继续向更深处探索。

\vspace{2em}

少时，白言冰、郑天明两人行至一片沼泽地。

这里的泥土呈现出一种令人不安的紫黑色，踩上去绵软湿滑，偶尔有气泡从深处汩汩冒出，破裂时散发出淡淡的硫磺与腐殖质混合的臭味。一些地方积着浅浅的死水，水色浑浊，泛着油光，水面漂浮着不知名的灰白色絮状物。更多地带则是看似坚实、实则布满孔洞的奇特地表，像是一块巨大的、正在缓慢腐烂的海绵。

暗沼区域广袤，无法绕行，只能更加小心地选择落足点。

白言冰走在前方，每一步都经过仔细感知。他的真元更加凝练纯粹，对环境中能量和物质结构的细微异动感知也更为敏锐。他能“听”到脚下泥土深处细微的流水声，能感觉到某些区域地气流转的滞涩或狂暴，往往能在郑天明察觉之前，就用剑鞘或石子试探出隐藏的流沙坑或气穴。

然而，古战场的诡谲远超想象。就在他们以为已经穿越了最危险的区域，前方出现一片地势稍高、布满了风化岩柱的干燥地带时：

那是一根看似普通、半截埋入地下的灰白色岩柱。郑天明经过时为了保持平衡，手掌下意识地在那岩柱粗糙的表面撑了一下。

触感冰凉坚硬，与周围岩石无异。

但就在他手掌离开的瞬间，岩柱表面那些看似天然形成的、扭曲的纹路，突然像活过来一般，亮起了幽暗的紫红色光芒！光芒并不耀眼，却带着一种直透骨髓的邪异。

“不好！” 白言冰厉喝，一把抓住郑天明的后襟向后急扯。

“嗡——！”

低沉而令人牙酸的震颤声从岩柱内部传来，紧接着，以岩柱为中心，方圆十丈内的地面猛地向下一沉！并非塌陷，而是仿佛变成了粘稠无比的泥潭，更准确说，是活化了的、充满恶意的沼泽！

紫黑色的泥浆翻滚上涌，瞬间没过了脚踝，并且产生一股强大无比的吸力，牢牢锁住了两人的双腿。这吸力并非单纯物理上的拉扯，更伴随着一种阴寒的能量，如同无数细小的冰锥，顺着腿部经络急速向上蔓延，疯狂吞噬着他们体内的灵力与热量！

与此同时，岩柱上那些发光的纹路如同血管般鼓胀起来，“噗”地一声，从几个孔洞中喷涌出大股浓稠的墨绿色雾气。雾气带着刺鼻的腥甜气息，迅速弥漫开来，所过之处，那些风化的岩柱表面竟发出“滋滋”的声响，被腐蚀出坑坑洼洼的痕迹。

郑天明面色惨变，他感觉自己的腿部正在迅速失去知觉，灵力如决堤般外泄，更要命的是，那毒雾已然扑面而来！

白言冰眼神一凝，手中长剑清鸣出鞘。他将精纯的真元猛地灌注剑身，手腕一抖，数道凝练如实质的淡金色剑气呈扇形向下斩出！

噗噗噗！

剑气切入泥沼，随即爆开！蕴含在剑气中的、高度压缩的问道真元瞬间释放，形成小范围的能量冲击，短暂地扰乱了泥沼底部禁制的能量流动。两人腿部的吸力为之一松。

就在这时，泥沼中“嗖嗖”射出七八条黑影，快如闪电，直缠两人腰身！那竟是方才在沼泽边缘见过的、长满倒刺的诡异藤蔓，只是此刻的藤蔓颜色更深，近乎黑色，顶端甚至裂开细小的口器，滴落着粘液。

“滚开！” 白言冰冷哼一声，剑光回旋，如同精准的手术刀，剑锋贴着郑天明和自己身侧掠过，一连串令人牙酸的断裂声响起，那些诡异藤蔓被齐根斩断。断口处喷出腥臭的汁液，落在泥沼里冒起白烟。

趁此间隙，白言冰左手在腰间一抹，一根流光溢彩的长索已然在手，正是捆过柳媚的那条法宝。

在临行前，叶牧谦把这玩意给他时，白言冰还在思考这东西对他有什么用。

没想到在救人的时候派上了用场！

他手腕一抖，长索如灵蛇般缠住郑天明的腰际，随即沉腰坐马，体内真元轰然爆发，手臂肌肉贲起，吐气开声：“起！”

“噗嗤！”

郑天明只觉一股沛然巨力从腰间传来，整个人被硬生生从粘稠如胶的泥沼中拔起，甩向后方干燥的地面。他在空中勉强调整身形，落地时踉跄几步，单膝跪倒，随即“哇”地吐出一口带着黑气的淤血，脸色由青转白，气息骤然萎靡下去。虽被拉出险境，但侵入体内的毒雾和禁制冲击已然爆发，内腑受创不轻，灵力运转几乎停滞。

白言冰紧随其后，足尖在几块凸起的、尚未被泥沼完全吞没的石块上连点，身形如轻烟般飘然而出，落在了郑天明身旁。

他顾不上查看自己是否沾染毒瘴，立刻扶住郑天明，从怀中取出一个玉瓶，倒出一粒龙眼大小、散发着沁人心脾清香的碧绿丹药，正是陈千羽以青鸾传承结合自身感悟炼制的护心丹，最擅解毒固元、平衡生机。

他迅速将丹药塞入郑天明口中，同时单掌抵住其后心，精纯温和的问道真元缓缓渡入，助其化开药力，导引药力护住心脉、驱散侵入经络的阴寒毒力。

丹药入腹，化作一股温润清流，迅速扩散。郑天明只觉得火烧火燎的五脏六腑仿佛被清泉洗涤，那股阴寒麻痹的感觉被徐徐驱散，紊乱的气息也渐渐被抚平。他剧烈咳嗽几声，又吐出几口带着腥味的黑血，脸色这才稍微恢复了一点人色，但依旧苍白如纸，虚弱不堪。

“白兄……咳咳……又救我一命。”郑天明喘息着，看向白言冰的眼神充满了感激与愧疚。若非自己一时不察，何至于此。

“少说话，凝神调息，引导药力。”白言冰言简意赅，声音依旧平稳，但他收掌之后，目光锐利地扫视着周围，特别是那片重新恢复“平静”、只是泥浆依旧缓缓翻涌的诡异沼泽，以及那根重新黯淡下去、却更显狰狞的岩柱。

他持剑而立，为其护法，心中对这片古战场、乃至对整个灵界修士生存环境的认知，变得更加冷峻切实。

弱肉强食，在此地被演绎到了极致：不仅是修士之间的互相倾轧，这环境本身，就像一头沉睡的、布满致命机关的远古凶兽，任何一丝疏忽，任何一点贪念，都可能唤醒它，将闯入者吞噬得骨头都不剩。

比起内域至少还有基本秩序、灵界小玄界外面还有世家与宗门维持，这里纯粹是力量与生存的赤裸丛林，危机四伏的程度，犹有过之。

接下来的路途，因郑天明受伤，行动能力大减，探索的主力变成了白言冰一人。郑天明更多时候是在白言冰的掩护下，艰难挪移，抓紧每一刻调息恢复。

白言冰将自身的感知能力提升到了极限。他的真元流转周身，如同最精密的雷达，捕捉着空气中灵气的每一丝异常波动，脚下土壤密度的细微差别，乃至空气中残留的、其他修士经过时留下的极淡气息。

他往往能在肉眼看见之前，就通过气流、温度、能量场的微妙变化，判断出前方是否存在隐形的空间裂缝、地火气脉的薄弱点、或是某种沉睡妖物的潜伏范围。

选择路径时，他不再仅仅避开明显的危险或人群，而是综合考量地形、能量流动、视线遮蔽以及可能的退路，近乎本能地规划出最优路线。遇到无法绕行、必须穿越的危险地带，白言冰便展现出其超凡的应对能力。那份对时机、角度、力道的精确把握，在郑天明看来已然超越了寻常术法的范畴，更近乎一种融入战斗本能的“艺”。

当然，避不开的战斗依然存在。古战场深处，亡命之徒更多，也更加狡猾和强悍。

\vspace{2em}

跋涉的日子在紧张与喘息间交替，足足过了半月有余。古战场的天空依旧是永恒的铅灰色，但周围的景象开始随着跋涉发生变化。

焦黑的大地逐渐被更多巨大的人造物残骸覆盖——崩碎的、铭刻着模糊符文的城墙断口；斜插入大地、长达数十丈、锈蚀斑驳的巨弩残骸；只剩下基座和几根巍峨巨柱的广场；还有更多倒塌的、样式古朴的殿宇楼阁。

煞气变得更加浓郁，空气中仿佛弥漫着无形的重量，让人呼吸都不由自主地放缓。但与此同时，某种苍茫、古老、宏大的气息，也开始从废墟深处隐隐透出。

终于，在穿越一片由无数碎裂法器、铠甲碎片堆积而成的、如同金属坟场般的斜坡后，眼前的视野豁然开朗。

一片难以想象的巨大废墟，展现在两人面前。

这里显然曾是一个规模极其宏大的上古宗门遗址。目之所及，尽是断壁残垣，但即便是这些残骸，也透露出昔日的辉煌与磅礴。高达百丈、即使断裂也依旧耸入灰雾的巨型门楼石柱；绵延数里、铺满了精美浮雕（大多已残缺）的广场地基；以及那些即使坍塌，其基座范围也大得惊人的殿宇群落……无一不在诉说着往昔的鼎盛。

废墟整体呈现一种悲凉的灰白色，覆盖着厚厚的岁月尘埃，许多地方生长着颜色暗沉、形态扭曲的耐煞植物，更添荒寂。而在废墟的最中央，众星捧月般的位置，一座巍峨的殿宇相对完整地矗立着。

它的墙壁上同样布满了裂痕与战斗留下的创口，檐角多有残缺，但主体结构依然稳固，透露着难以言喻的坚固与苍古。殿宇的样式古朴厚重，与周边废墟同源，却又明显高出数个等级，通体用一种非金非玉、泛着暗沉青灰色的奇异石材构筑，历经万古煞气侵蚀，依旧屹立不倒。

最引人注目的是，整座殿宇被一层淡淡的、犹如水波般流动的乳白色光晕所笼罩，那光晕柔和却坚韧，将周遭浓郁的煞气和破败之气隔绝在外，让那座殿宇在这无边废墟中，犹如风暴眼中的一点静谧之地，散发着神秘而诱人的气息。

殿门紧闭，门前有九级高大的台阶，台阶材质与殿宇相同，光洁如镜，纤尘不染，与周遭格格不入。

无需多言，两人心中同时明悟——历经艰险，他们终于抵达了这处小玄界最核心的区域。那座被神秘光晕守护的巍峨古殿，便是这一切凶险与机缘的终点，上古宗门留存的，最核心的传承所在。

站在废墟边缘，望着远处那孤高而神秘的殿宇，无论是重伤初愈、心有余悸的郑天明，还是一路披荆斩棘、冷静如冰的白言冰，心中都难以抑制地升起一股复杂的情绪。有抵达终点的如释重负，有对可能获得机缘的期待，但更多的，是一种面对未知与古老的深深敬畏。最后的考验，或许就在那紧闭的殿门之后。

然而，当两人赶到时，大殿周围已是一片狼藉，显然经历过激烈的争夺。殿门已然洞开，光晕黯淡。郑天明通过家族内部特殊的联络方式，得到了一个让他沮丧的消息：核心传承，早已被提前进入的郑氏核心队伍取走。

“果然……汤都没剩几口热的。”郑天明苦笑着对白言冰道，“白兄，看来咱们这趟，真的就是来走过场了。进去看看吧，捡点边角料，也算没白来。”

两人进入大殿。殿内空间广阔，但许多地方已被翻找得乱七八糟，有价值的物品早已被搜刮一空。他们在一些偏僻的侧殿、损坏的傀儡残骸旁，找到了一些被前人遗漏或看不上的东西：几枚记录着上古功法的残破玉简，数件灵光黯淡、有明显破损痕迹的法阶、道阶法器，以及一些零散的、品相还算不错的炼器材料和灵草。

最让白言冰眼前一亮的，是在一个倒塌的石柜下方，发现了一个隐藏的夹层，里面整整齐齐码放着数百块晶莹剔透、灵气充沛的灵石——而且是上品灵石！粗略一算，算成更常用的下品灵石，竟有数万之多！这对于目前主要依靠下界带来灵石积蓄、在灵界唯一稳定收入只有陈千羽行医所得的师徒几人来说，简直是雪中送炭，足以支撑很长一段时间的生活和基本修炼资源采购。

“太好了！白兄，咱们平分！”郑天明也高兴起来。

白言冰却摇了摇头：“郑兄，这一路多亏你带我进来，又屡次涉险。这些灵石，你多拿些。”他深知这些灵石对己方的重要性，但也不愿占朋友太多便宜。

郑天明却大手一挥，爽快道：“多分什么多分！白兄，要不是你，我这条命早丢在路上了，哪还有命拿灵石？再说了，我虽然在郑氏不受重视，但基本的月例、再做些小生意，日常花销还是不愁的，不缺这点。我看得出，你们……似乎手头不宽裕？这些灵石正好解燃眉之急。”

白言冰闻言，心中感动。郑天明看似憨直，实则心细体贴。他也不再矫情，郑重抱拳：“既如此，白某愧领了。这些灵石，确实解我师徒眼下之困。郑兄高义，白某铭记。”

他将自己面前那堆灵石悉数收起。想了想，又将那些找到的功法玉简和破损法器仔细筛选，将其中内容相对完整、法器结构尚有研究价值的，挑出大半，强行塞给郑天明：“这些你务必收下。功法虽残，或许能触类旁通；法器虽损，但材料与上古炼制手法仍有借鉴之处，对你郑家炼器之道或有益处。”

郑天明推辞不过，见白言冰态度坚决，便也爽快收下，笑道：“那便多谢白兄了！这下回去，也不算空手，好歹能交差。”

两人又在殿内细细搜寻一番，确认再无遗漏，便按原路退出小玄界。

出来时，入口处已然冷清了许多，大部分修士或早已离开，或永远留在了里面。

\vspace{2em}

回到流云坊市，郑天明坚持要亲自去流云山脚小院向叶牧谦道谢。白言冰便带他一同返回。

两人回到小院时，陈千羽正好结束一次短途采药归来，正在院中分拣药材。她一身素雅衣裙，身姿婀娜，专注的神情在午后阳光下显得格外宁静美好。她察觉到有人进来，抬头望去，见是白言冰带着一位陌生青年回来，微微颔首示意。

郑天明一眼看到陈千羽，只觉得眼前一亮。这女子容貌清丽绝俗，气质温婉中带着一丝不易察觉的坚韧，尤其那双眸子，清澈明净，仿佛能映照人心。他一时竟看得有些呆了，连白言冰的介绍都没听清。

直到陈千羽注意到郑天明气息不稳、面色隐带青灰，医者本能让她微微蹙眉，主动开口道：“这位公子似有内伤，可是在小玄界中受了暗毒侵蚀？”

郑天明这才回过神来，连忙道：“是，是，不小心吸入了毒沼雾气，多亏白兄相助……”

陈千羽放下手中药材，走上前来：“若不嫌弃，容我为你诊视一番。” 她声音柔和，却带着令人信服的力量。

郑天明晕乎乎地伸出手。陈千羽三指搭在他腕脉上，一丝精纯平和的真元探入，细细感知。片刻后，她收回手，道：“毒已解了大半，但肺脉与肝经仍有郁结，灵力运转滞涩。需服药调理数日，并辅以行气导引之法。” 说着，她便转身去屋内取来金针和纸笔，动作娴熟流畅。

郑天明看着她专注书写药方的侧影，只觉得心跳都漏了几拍。待陈千羽写好药方，又亲自为他施针疏通部分经脉后，那股滞涩感明显减轻，精神也为之一振。

“多谢……多谢姑娘！”郑天明连忙道谢，又忍不住偷偷瞥向正在收拾针具的陈千羽。

这时，白言冰走过来，很自然地接过陈千羽手中的药包，对郑天明介绍道：“郑兄，忘了介绍，这是舍妹，陈千羽。”

郑天明脸上的笑容瞬间僵住，眼睛瞪大，看看白言冰，又看看陈千羽。

不对啊，他姓白，她姓陈，怎么还能称呼“舍妹”的？

总不能是更进一步的某些关系吧？

怎么好姑娘全都名花有主了？

算了，保险起见，还是不要试图追求她了。

郑天明顿时脸拉得跟苦瓜一样长，拿了药，告辞离去。

走出小院不远，他回头望了一眼那掩映在绿树中的朴素院落，想起白言冰高深莫测的剑术，陈千羽出神入化的医术，还有那位始终气度沉静、仿佛一切尽在掌握的叶前辈，不由得低声嘟囔了一句：

“我就知道……同一门下的，都差不了……”

语气里，羡慕有之，感慨有之，或许还有一丝连他自己都未察觉的向往。

\vspace{2em}

白言冰将郑天明送至院外，目送其身影消失在巷口，这才转身回到正屋。叶牧谦正坐在窗边的竹椅上，手中把玩着那枚从小玄界带出的、记录着残缺上古功法的玉简，目光却似乎投向了更远的虚空。陈千羽已收拾好针具药箱，正安静地坐在一旁，素手烹茶，氤氲的水汽模糊了她沉静的眉眼。

“师尊。”白言冰上前一步，躬身行礼。

叶牧谦收回目光，看向他，微微颔首：“此行辛苦了。收获如何？可有受伤？” 语气平和，却透着关切。

白言冰直起身，开始详细汇报小玄界中的经历。从入口处的纷乱局势，到一路上的机关陷阱与修士争斗，再到石林迷阵中的惊险、毒沼禁制下的救援，最后是藏宝大殿内的所见所得。他语速平稳，条理清晰，将上品灵石、残破功法玉简、损坏法器以及郑天明坚持将灵石全部相让、自己则以相对完整的功法法器回赠等细节一一道来。

“那郑天明，初识时只觉其憨直跳脱，甚至有些不务正业。” 白言冰汇报完毕，略一沉吟，脸上露出些许感慨之色，“但经此一行，方知其人性情。危难时不弃友，得利时不独占，心思也算细腻，察觉我等境况，便以灵石相赠，又不愿令我觉得亏欠，言语颇为体贴。我本以为，似郑氏这等超级世家出来的公子哥，纵然是旁支，也多半是那种眼高于顶、吊儿郎当、不学无术的纨绔子弟，未曾想还有郑天明这般重情重义、忠厚老实之人。”

他这番话发自肺腑。在荒域，他是皇子，见过勋贵子弟的跋扈；在内域，他接触过各大势力的精英，深知其中不乏傲慢之辈。郑天明的确刷新了他对“世家子弟”的一部分刻板印象。

叶牧谦听着，手指轻轻摩挲着温润的玉简，脸上露出一丝淡淡的、意味深长的笑容。他抬眼看向白言冰，又瞥了一眼旁边静坐聆听的陈千羽，缓缓开口道：“言冰，你此言，却是将人看得太过简单，也看得太割裂了。”

白言冰微微一怔：“师尊的意思是？”

“你说世家公子哥多是纨绔，那为师问你，”叶牧谦语气平和，却带着引导的意味，“你白言冰，出身荒域大周皇室，是名副其实的太子，算不算‘世家公子哥’？千羽虽非皇室血脉，但自幼生长于宫廷，锦衣玉食，见识教养皆非凡俗，算不算‘世家小姐’？便是瑶儿，虽然十年出门游历，但天师府府主之女，内域顶尖势力的千金，这身份又当如何？”

白言冰与陈千羽闻言，俱是一愣，随即陷入思索。

的确，若按出身论，他们三人都可归入所谓的“世家”范畴。

叶牧谦继续道：“郑天明是郑氏旁支，资源有限，前途不明，故而性格中多了几分市井的实在与憨厚，少了几分核心子弟的骄矜。这是环境使然，却并非本质。而你，言冰，你所见所历，让你养成了沉稳刚正、心怀苍生的性子；千羽醉心医道，仁心仁术，外柔内韧；瑶儿嫉恶如仇，执拗坚定……这些，难道不也是你们出身与经历塑造的一部分吗？”

他顿了顿，声音温和却带着一种洞彻的力量：“看人看事，万不可如此非黑即白，截然割裂。世间之人，复杂多面。便是同一个人，在不同境遇下，也可能表现出截然不同的面貌。有世家子骄横，也有世家子仁厚；有寒门子勤奋，也有寒门子奸猾。重要的不是贴上一个标签，而是去理解其行为背后的缘由、处境，看到其矛盾统一之处。”

“‘一阴一阳之谓道。’ 阴阳并非对立，而是相辅相成，互为其根。人性之中，善恶、刚柔、智慧愚钝，往往并存。郑天明有其憨厚实在的一面，或许也有其局限与缺点；那些看似纨绔的子弟，或许在特定领域也有其专注或长处。取其精华，观其全貌，方能不偏不倚，接近真实。这便是‘对立统一’地、辩证地看问题。”

白言冰听着师尊的教诲，心中似有明光划过。他想起自己当初对青云门李轩等人的判断，想起在内域与各势力打交道时的观察，再结合此次与郑天明的接触，确实感到自己以往的看法有些简单化了。他肃然躬身：“弟子受教。是弟子思虑不周，以偏概全了。”

叶牧谦摆摆手：“无妨，经历便是最好的老师。日后留心观察，自会体悟更深。这些灵石来得及时，可解我们当下之困。那些功法法器虽残，亦有其参考价值。你做得不错，既未贪心，也全了朋友之义。”

他又询问了一些关于小玄界内灵气环境、残留禁制特点的细节，白言冰一一作答。师徒三人又聊了片刻，便各自忙于自己的事情了。白言冰去清点规整带回的物品，陈千羽继续研究她的医药，叶牧谦则拿着那枚玉简，若有所思。

\vspace{2em}

郑氏家族的宅邸区域。

郑天明回到自己那处不算宽敞、但比起流云山脚小院却要精致许多的独栋小楼。他是旁支，住所自然无法与核心区域的琼楼玉宇相比，但还算不错。

他心情颇好，虽然小玄界中险死还生，但结识了白言冰这等人物，又得了一些上古之物，自觉收获大于惊吓。他哼着小调，将白言冰分给他的几枚玉简和两件相对完好的破损法器小心收好，又摸了摸怀中那份陈千羽开的药方，打算明日就去抓药。

就在这时，院门被不轻不重地叩响。声音规律而带着一种不容忽视的力度。

郑天明心头一跳，这敲门风格……

他赶紧整理了一下衣衫，快步走到门前，拉开一看，门外站着的，正是他的堂兄，郑氏这一代的核心翘楚之一——郑天河。

郑天河看上去约莫二十七八岁年纪，面容俊朗，身材挺拔，穿着一身暗绣云纹的锦袍，腰间悬着一枚莹润的玉佩，气息渊深绵长，赫然已是法相境通灵期的修为。他也刚从小玄界回来不久，收获显然是颇为丰厚的。

郑天河眉宇间带着一丝惯常的沉稳与隐隐的威严，目光平静地落在郑天明身上，却让郑天明感到一股无形的压力。

“天……天河哥？你怎么来了？快请进！”郑天明连忙侧身让开，脸上堆起笑容，心里却有些打鼓。这位堂兄天赋卓绝，深受家族长老们器重，一直被作为下代家主的有力竞争者培养，平日事务繁忙，极少来他这小院。

郑天河迈步而入，目光随意地扫过小院，最后落在郑天明还有些许苍白的脸上，淡淡道：“听说你从小玄界回来了，还受了伤？看来此行不甚顺利。”

郑天明挠挠头：“是遇到些麻烦，不过都过去了，还得多亏一位朋友相助。”

“朋友？”郑天河眉头微不可察地蹙了一下，“我听闻，你此次是与一个陌生剑修同行？并非家族安排之人。”

郑天明心里咯噔一下，果然是为了这事。他不敢隐瞒，将如何结识叶牧谦、如何邀请白言冰同行、小玄界中白言冰数次相救、最后分配收获等事情，简明扼要地说了一遍，重点强调了白言冰的可靠与实力，以及对方并无恶意，只是寻常散修结交。

郑天河静静听着，手指无意识地在身旁的石桌上轻轻敲击。他了解自己这个堂弟，憨厚有余，机变不足，但本性不坏，看人的眼光也还算直白。

听他描述，那剑修白言冰倒像是个可交之人，实力心性都不差。而且对方并未觊觎郑氏什么，反而在危险中护住了郑天明。

心中的那点疑虑和淡淡的不悦渐渐消散。

郑天河身为核心，对家族利益和内部团结看得极重。起初听闻旁支子弟与来历不明的外人走得近，他本能地有些不快，觉得这是“不分内外”。但听完解释，又觉得情有可原，毕竟郑天明这次能活着回来，多亏了那人。

只要不是心存叵测之辈，结交些有本事的散修，对郑天明而言也未必是坏事。

“罢了，既然是你的救命恩人，日后当知恩图报。但也要记住，你是郑氏子弟，言行须有分寸，莫要让人轻易利用了去。”郑天河语气缓和了些，告诫道。

郑天明连忙点头：“是，天明明白。白兄他们不是那样的人。”

\vspace{2em}

正事说完，气氛稍缓。

郑天河此次过来，除了询问小玄界之事，也存了几分关心堂弟伤势的意思。他目光随意扫过屋内陈设，忽然，被桌角一件半掩在杂物下的东西吸引了视线。

那是一个巴掌大小、造型粗糙的铜镜，镜框边缘还有未打磨干净的毛刺，正是柳成制作、郑天明之前“定制”的那枚留影镜。大约是郑天明回来后随手放在那里，还没来得及收好。

“那是什么？”郑天河随口问道，走上前，伸手拿起了铜镜。入手沉甸甸，工艺实在不敢恭维，完全是坊市底层炼器学徒的水平。

郑天明脸色“腾”地一下红了，支吾道：“没……没什么，一个小玩意儿，朋友做的……”

郑天河瞥了他一眼，见他神色窘迫，心中起疑，下意识地注入一丝灵力。

浑浊的镜面亮起，一阵模糊晃动后，开始浮现出不堪入目的影像片段，虽不清晰，但足以辨认内容。

“胡闹！”郑天河脸色一沉，低声斥道，“你从哪弄来这种有伤风化的东西？简直不成体统！我郑氏子弟，岂可沉迷此等秽物？还不快拿去丢了！”

他心中涌起一股怒意，觉得这堂弟实在是不争气，修为不精进，却弄这些歪门邪道的东西。

郑天明被训得面红耳赤，喏喏不敢言，伸手就要接过铜镜。

然而，就在郑天河准备将铜镜丢还给他的瞬间，动作却猛地顿住了！他脸上的怒容迅速褪去，取而代之的是一种极度的专注与惊疑。他的目光死死盯住那依旧在播放模糊影像的镜面，瞳孔微微收缩。

不对！

这是动态的影像！

虽然粗糙模糊，但确实是连贯的动作记录！

而且，这铜镜的结构……他神识扫过，虽然工艺低劣，但核心的几处阵纹组合，却实现了一种他从未见过的、简易却有效的动态影像记录与重现功能！所用的材料，也都是最廉价低阶的昏黄石和普通赤铜！

廉价、简易、可批量复制……的动态留影技术？！

郑天河的心脏不受控制地剧烈跳动起来。

作为一个被当作下代家主培养、目光早已超越个人修为、更多投向家族产业与未来战略的核心子弟，他几乎是瞬间就洞悉了这项看似“不堪”的技术背后，所蕴含的何等恐怖的潜力与价值！

战场记录与情报传递？无需文字描述，直接呈现动态影像，敌我态势、地形地貌、功法招式，一目了然！

民用娱乐与信息传播？远非现在单调的留影石可比，可以记录戏剧、讲述故事、传播知识，甚至满足某些特殊需求（比如郑天明这种），市场空间不可估量！

技艺传承与教学？炼丹的火候、炼器的锤法、复杂法术的施展过程，都可以动态记录，供后人反复观摩学习！

无论是直接将成品作为独家商品售卖，还是将这项技术作为筹码与某些势力合作，甚至仅仅是在家族内部应用以提升效率，都能带来难以想象的巨大利益和战略优势！

这哪里是什么“有伤风化的秽物”？这分明是一座尚未有人发现的、巨大的金矿！不，是比金矿更珍贵的技术基石！

郑天河深吸一口气，强行压下心中的激动狂澜。他脸上的神情已经恢复了平静，甚至比刚才更加深沉。他缓缓收回了准备递出铜镜的手，转而将其紧紧握在掌心。

“天明，”他的声音变得异常平稳，甚至带上了一丝罕见的温和，“这东西……制作此物的朋友，是谁？”

郑天明被堂兄这瞬间的转变弄得有些懵，下意识答道：“是……是坊市西边摆摊的一个朋友，叫柳成，他喜欢鼓捣这些稀奇古怪的……”

“柳成……”郑天河默默记下这个名字，目光灼灼地看着郑天明，“这枚铜镜，暂时放在我这里。”

“啊？”郑天明一愣，“天河哥，你不是说……”

“我改主意了。”郑天河打断他，语气带着不容置疑的决断，“此物虽内容不堪，但其制作思路颇有些意思。我需要研究一下。此事，你不必再对任何人提起，包括柳成。明白吗？”

郑天明看着堂兄那深邃而充满压迫感的眼神，哪里还敢说半个不字，连忙点头：“明……明白。”

郑天河将留影镜收入自己的储物戒指，又深深看了郑天明一眼，拍了拍他的肩膀：“好好养伤。今日之事，你做得……不算全错。”

说罢，不再停留，转身大步离去，身影很快消失在院门外。

郑天明站在原地，摸了摸还有些发烫的脸，又看了看堂兄消失的方向，一头雾水。

天河哥怎么突然对那“不正经”的镜子感兴趣了？还要保密？

他摇了摇头，想不明白，索性不再去想。反正天河哥做事，总有他的道理。

\vspace{2em}

而离开小院的郑天河，步伐沉稳，心中却已翻腾起惊涛骇浪。

他漫步在郑氏宅邸区整洁的石板路上，看似从容，脑海却在飞速运转。

柳成……一个在坊市摆摊、鼓捣奇巧的底层修士。郑天明的朋友。这意味着技术源头目前掌握在一个微不足道的小人物手里，而且与郑天明有关联。这是优势，也是风险。优势在于，可以通过郑天明这条线，相对容易地控制住源头。风险在于，消息可能已经小范围泄露（至少叶牧谦师徒可能知道郑天明有这东西），必须尽快行动。

这项技术，绝不能现在就上报家族！

郑天河眼神锐利。家族长老会势力盘根错节，一旦上报，这巨大的功劳和随之而来的利益分配，立刻会成为各方争夺的焦点。他虽然是下代家主的有力竞争者，但并非没有对手。届时，这项技术很可能被某个长老派系攫取，自己最多分润一点功劳，无法将其完全转化为自己积累威望、夯实地位的基石。

他需要这个项目，一个独属于他郑天河，能彰显其眼光、能力和手腕，并能带来实实在在巨大收益和影响力的项目！ 这留影镜技术，简直是天赐良机！

“必须私下处理！”郑天河心中定计，“先控制住柳成，确保技术独家掌握。然后，以我个人的名义，秘密组建一个小型的研发和试产团队，改进工艺，提升清晰度和稳定性，开发出不同用途的型号。初期，可以借用在家族中我能完全掌控的、不那么起眼的炼器作坊和销售渠道。”

他想到了流云坊市，虽然整体上属于家族，但实际有少数门市是由他母亲那一系的旁亲在打理，忠诚度相对可靠。这可以作为起点。

“资金……我自己的积蓄，启动应该够了。关键是保密。”郑天河思忖着，“柳成那边，可以用重利诱之，也可以用郑天明的安危稍作敲打，务必让他签下最严格的保密和独家授权契约。至于郑天明，他老实，好糊弄，只需告诉他这是家族机密任务，让他配合即可，不必知道全貌。”

一个清晰而大胆的计划，在郑天河脑海中逐渐成形。他停下脚步，站在一处回廊下，望着远处流云坊市上空永不熄灭的灵光，嘴角勾起一抹志在必得的弧度。

\vspace{2em}

郑天河的动作极快。在确定那枚粗糙留影镜背后所蕴含的巨大价值，并决定将其作为自己积累资本、争夺家主之位的秘密项目后，他立刻动用了自己能够完全掌控的、隐藏在家族庶务体系中的几条暗线。

不过两日功夫，一份关于柳成的详细调查报告，便悄然呈递到了郑天河的书案上。报告内容简单得让郑天河都有些意外：柳成，年二十二，灵枢境中期修为，出身流云坊市底层，父母早亡，原在坊市一炼器铺做学徒兼处理边角料，后因“心思不专、喜好奇巧”被辞退。如今在西侧陋巷独自摆摊，售卖些自制的、无甚大用的精巧机关玩物。人际关系简单，除了与旁支子弟郑天明交好外，几乎不与任何修士深交。无师承，无背景，无异状。

“一个小小的灵枢境……背后竟真无任何靠山？”郑天河指尖敲击着报告，眼中精光闪烁。

这比他预想的最好情况还要好。

一个无根无底的底层修士，意味着掌控成本极低，技术泄露的风险也小。郑天明这条线，简直是天赐的桥梁。

他不再犹豫，唤来一名心腹护卫。此人名为李辉，真符境显化期修为，面容冷峻，是郑天河母亲一系培养出来的死士，忠诚可靠，办事利落。

“去西巷，找一个叫柳成的摊贩，请他来见我。”郑天河将一枚代表自己身份的玉牌递给李辉，语气平淡却不容置疑，“注意方式，莫要过于惊动旁人。若他不愿……你知道该怎么做。”

“遵命。”李辉接过玉牌，躬身一礼，转身离去，身影如鬼魅般融入阴影。

\vspace{2em}

流云坊市，西侧陋巷，柳成摊位。

柳成刚刚完成一件新作品——一个会模仿简单琴音的灵巧八音盒，正小心翼翼地调试着最后几个簧片。他心情不错，前几日叶前辈买走了他做的星象仪，还随口提点了两句关于共振频率与灵力波动契合的问题，让他茅塞顿开。摊位上零零散散摆着几件旧作，虽无人问津，但他乐在其中。

忽然，他感到一股无形的压力笼罩过来，周围的喧嚣似乎都安静了一瞬。他抬起头，只见一个身着郑氏护卫服饰、面容冷硬的中年男子，正站在摊位前，目光如鹰隼般盯视着他。那眼神中的审视与漠然，让柳成心头猛地一紧，手中精巧的镊子“啪嗒”一声掉在摊位上。

“你叫柳成？”李辉开口，声音干涩，不带丝毫感情。

“是……是我。这位前辈，有何指教？”柳成连忙站起身，手足无措地行礼。真符境修士，对于他这灵枢境来说，已是需要仰望的大人物，更何况对方身着郑氏护卫服饰，代表的可是流云坊市的主宰者。

李辉亮出那枚刻有“天河”二字的玉牌：“我家公子有请，随我走一趟。”

公子？天河？柳成脑子里“嗡”的一声。

郑天河！郑氏核心翘楚，下代家主的有力竞争者！

这样的大人物，怎么会突然找上自己这个微不足道的小摊贩？他瞬间联想到郑天明，以及郑天明不久前从他这里“定制”的那枚留影镜……难道是那东西惹祸了？

一股寒意从脚底直冲头顶。柳成虽然憨直，但并不傻，相反，底层挣扎求生的经历让他对危险有种本能的敏锐。

大人物相召，往往祸福难料，尤其是他这种毫无背景的小人物，被叫去后是生是死，全在对方一念之间。那留影镜里的内容……若是被郑天河那样的正经人物看到，自己恐怕吃不了兜着走！

“前……前辈，不知贵公子召见小人，所为何事？小人……小人摊位离不开人，能否……”柳成脸色发白，试图推脱，声音都有些发抖。

李辉眉头微皱，眼中闪过一丝不耐。

公子的命令是“请”，但若请不动，他也不介意用些手段。

他上前一步，压迫感更强，皮笑肉不笑地说：“郑氏主脉相请，谁敢不从？莫要敬酒不吃吃罚酒。今天你去也得去，不去也得去！”

话音落下，他一只手已缓缓抬起，指尖有灵光隐现，显然准备若柳成再敢推拒，便要强行擒拿。

柳成吓得连连后退，撞翻了身后的工具架，叮当作响。周围几个摊贩和路过的低阶修士见状，纷纷躲远，生怕被牵连。

在这流云坊市，郑氏就是天，谁敢管郑氏的闲事？

就在柳成绝望，李辉的手即将触及他肩膀的刹那——

“且慢。”

一个平和的声音忽然响起，并不高亢，却奇异地穿透了巷道的嘈杂，清晰地传入两人耳中。

李辉动作一顿，霍然转头。只见巷口不知何时多了一人，青衫磊落，面容普通却气度沉凝，正是信步走来的叶牧谦。他原本只是如往常般来坊市闲逛，顺道看看柳成是否有新作，却不想撞见了这一幕。

叶牧谦缓步走近，目光先扫过脸色惨白、惊魂未定的柳成，然后落在李辉身上，淡淡道：“这年轻人本不想去。你强行逮人，这是待客之道？”

李辉心中一凛，目光在叶牧谦身上仔细打量。气息晦涩，深浅难测！以他真符境显化期的修为，竟完全看不透对方虚实。

此人衣着朴素，不像世家大派子弟，但那份从容不迫的气度，绝非寻常散修能有。

他立刻想起修真界一些关于隐世老怪或其门人弟子喜欢隐藏修为、游戏风尘、扮猪吃虎的传闻。

眼前这人，八成就是此类！

李辉跟随郑天河日久，见识不浅，深知这类人物最是麻烦，往往实力强横且脾气古怪。自己若当街与其冲突，无论胜负，都可能给公子惹来不必要的麻烦。 公子吩咐的是“请”人，不是结仇。

电光石火间，李辉已权衡利弊。他缓缓收回手，脸上挤出一丝生硬的笑容，对叶牧谦抱拳道：“这位道友，在下奉命行事，并无恶意。既然这位小友不愿，那便罢了。”他又冷冷瞥了柳成一眼，撂下一句：“你且等着。”

说罢，竟不再纠缠，转身快步离去，显然是回去搬救兵了。

\vspace{2em}

劫后余生的柳成腿一软，差点瘫坐在地。他看向叶牧谦，眼中充满了感激与后怕：“叶……叶前辈！多谢前辈救命之恩！”

叶牧谦摆摆手，眉头微蹙。郑天河的人为何突然来找柳成？而且态度如此强硬？联想到郑天明与柳成的交集，以及郑天明可能从柳成那里得到过留影镜……叶牧谦心中隐隐有了猜测。

郑天河，恐怕是看出了那简陋留影镜背后技术的价值！

“此地不宜久留。”叶牧谦对柳成道，“郑氏的人很快会再来，而且下次来的，恐怕就不是这么好打发的了。你可有去处？”

柳成脸色惨白地摇头。他孤身一人，哪有什么靠山去处？铺子？这铺子哪里还敢要！

“若信得过我，先随我离开。”叶牧谦果断道。

柳成此刻已将叶牧谦视为唯一的救命稻草，哪有不从之理？他胡乱将摊位上几件最珍视的作品和工具塞进一个旧包袱，连摊位都顾不上了，抱起包袱就跟在叶牧谦身后，两人迅速离开陋巷，向着坊市外流云山方向而去。

一路上，柳成心有余悸，又满心疑惑与恐惧。叶牧谦并未多问，只是带着他穿街过巷，专挑人少僻静处走。直到离开坊市范围，进入流云山脚相对荒僻的小道，柳成才稍微松了口气。

“叶前辈……那郑天河，为何突然要抓我？我……我从未得罪过他啊！”柳成终于忍不住，带着哭腔问道。

叶牧谦放缓脚步，看了他一眼：“或许，与你做给郑天明的那枚‘留影镜’有关。”

柳成如遭雷击，瞬间明白了。是了，一定是那东西！郑天明拿回去，可能被郑天河看到了……那里面可是……

“前辈！我……我不是故意的！那只是郑兄他……他非要我做的……”柳成急得语无伦次，又羞又怕。

“内容无关紧要。”叶牧谦打断他，目光深邃，“重要的是那镜子本身。郑天河不是寻常纨绔，他能看出那廉价留影技术背后的价值。匹夫无罪，怀璧其罪。你身怀此技而无自保之力，便是祸端。”

柳成呆呆地看着叶牧谦，他从未从这个角度想过。自己鼓捣出来的、被人视为“奇技淫巧”甚至“不堪入目”的小玩意儿，竟然蕴含着能引来杀身之祸的价值？

恐惧过后，一股难以言喻的复杂情绪涌上心头。有被认可的微微悸动，但更多的是对未来的茫然与恐惧。郑氏如庞然大物，他一个灵枢境小修士，能逃到哪里去？叶前辈能护他一时，能护他一世吗？

或许是看出了柳成眼中的绝望与挣扎，或许是觉得时机已至，叶牧谦忽然问道：“你的炼器手艺，构思精巧，尤其善于将复杂功能简化实现，这种思路，并非完全源于自学与偷师吧？可还有其他渊源？”

柳成浑身一震，抬头看向叶牧谦。叶前辈的目光仿佛能洞彻人心。在经历了刚才的生死危机，又感受到叶牧谦毫无保留的庇护之后，柳成心中那点最后的防备也消散了。他咬了咬牙，决定说出那个埋藏心底多年的秘密。

“前辈明鉴……晚辈的炼器手段，确实……并非完全来自自身。”柳成压低声音，仿佛怕被山风听去，“约莫十年前，晚辈那时还是个半大孩子，在流云山外数十里一处叫‘落星谷’的干涸河谷里捡柴火时，偶然捡到了一枚半埋在砂石里的玉简。”

“玉简？”叶牧谦眼神微凝。

“是。”柳成点头，回忆道，“那玉简材质特殊，非金非玉，入手温润。我当时好奇，发现里面记载了一些残缺的、关于炼器构型、灵路简化、材料物性利用方面的知识与图录。虽然很多地方深奥难懂，但也给了我极大的启发。我后来喜欢鼓捣这些精巧机关，很多思路都是从那玉简记载的内容中化用、演变而来。就连那留影镜……其实核心的‘昏黄石感光蓄灵、串联重现’的思路，也是受了玉简中一段关于‘光影留迹’的残缺记载启发，我自己摸索着试出来的。”

落星谷！玉简传承！

叶牧谦心中一动。这或许不仅仅是一个简单的机缘，更可能指向某处上古遗迹或隐修洞府！若真如此，其价值恐怕远超那留影镜技术。而且，若那里足够隐蔽，或许……能成为一个暂时的避风港，甚至未来道场的备选？

“你可知那玉简具体发现的位置？可还能找到？”叶牧谦问道，语气中带着一丝罕见的急切。

柳成想了想，肯定道：“记得！那地方很偏僻，河谷里怪石很多，我当时是在一块像卧牛的大石头下面的缝隙里发现的。虽然过去十年，但地貌应该变化不大，我应该能找到。”

“好！”叶牧谦当机立断，“我们先回流云山小院，带上言冰和瑶儿，立刻前往落星谷查看。”

柳成愕然：“现在？可是郑氏的人……”

“正因如此，才要立刻动身。郑天河的人很快会查到我的住处，那里也不安全了。落星谷地处偏僻，或可暂避，顺便探查那玉简来历。”叶牧谦解释道，同时已加快脚步。

两人迅速回到流云山脚小院。白言冰刚要动身离去，沈瑶刚刚回来，只有陈千羽外出采药未归。叶牧谦简短说明了情况，白言冰与沈瑶皆是神色一凛，意识到事态严重。没有任何犹豫，四人稍作收拾，带上必要物品，便由柳成带路，悄然离开小院，向着流云山外数十里的落星谷疾行而去。

\vspace{2em}

落星谷。

正如柳成描述，这是一条早已干涸不知多少岁月的古老河床。谷地宽阔，地面布满被水流冲刷得光滑圆润的卵石和嶙峋的怪石，两侧是陡峭的、植被稀疏的山崖。夕阳西下，将整个山谷染成一片昏黄，更添荒凉寂寥之感。

在柳成的指引下，四人很快找到了那块形似卧牛的巨石，巨石底部与河床地面之间有狭窄的缝隙。柳成指着缝隙某处：“当年玉简就是卡在这里。”

叶牧谦点点头，并未急于探查缝隙。他闭上双眼，通明境界大成的心念如无形的水银般铺开，细致地感知着周围每一寸土地、每一块岩石、每一缕空气的流动与灵气残留。

片刻后，他睁开眼，眼中闪过一丝了然。他走到卧牛石侧后方约三丈处，那里看起来与周围别无二致，都是乱石堆积。但叶牧谦却伸出手指，凌空虚划，一道道精纯的真元丝线随着他的指尖流淌而出，没入虚空之中，仿佛在勾勒着某种无形的脉络。

白言冰与沈瑶凝神观看，他们也能隐约感觉到那片区域有极其微弱、近乎自然的灵气韵律，与周围环境略有不同，但极其隐晦，若非师尊点明并以特殊手法激发，他们根本无从察觉。

随着叶牧谦的真元丝线勾勒完毕，眼前的空间微微扭曲了一下，一片约莫丈许方圆的区域，景象发生了微妙变化。

乱石依旧，但在乱石之中，隐约可见一道半掩埋的、布满风蚀痕迹的石碑，碑文已模糊难辨。石碑旁，则浮现出淡淡、几乎透明的阵法灵光，结构繁复古奥，与现今灵界常见的阵法风格迥异。

“好高明的隐匿阵法！几乎与自然环境融为一体，借用地脉残余灵气维持，若非岁月久远有所衰减，又恰逢日落时分阴阳交替、灵气波动最微弱之时，恐怕更难发现。”沈瑶轻声赞叹，她在幻术禁制上造诣颇深，更能看出此阵的不凡。

“我来试试。”白言冰上前，并指如剑，一道凝练的剑气试探性地刺向阵法灵光的一处节点。剑气触及灵光，如泥牛入海，悄无声息地被吸纳化解，阵法纹丝不动。

沈瑶也出手，指尖绽放出迷离的流光，试图以幻术迷惑阵法灵机，寻找破绽。流光渗入阵法，却仿佛撞入一团混沌，反馈回来的信息杂乱无章，无从解析。她摇了摇头，退后一步：“此阵原来是一把高明的锁头，自成一体，内外隔绝，幻术难侵。”

叶牧谦示意两人退开。他再次闭上双眼，心神彻底沉入通明之境。在他的心念感知中，那隐匿阵法不再是一团灵光，而是化作了无数细密交织的线，这些线的走向、强弱、彼此的联结与共振，都映照于心。

他追索着这些线的源头与规律，寻找那最关键、最细微的线头和结点。

\vspace{2em}

时间一点点流逝，夕阳完全沉入山后，山谷被暮色笼罩。白言冰与沈瑶静静护持在侧，柳成则紧张地攥着衣角。

足足过了半日，月上中天之时，叶牧谦忽然动了。他睁开眼，眸中似有星河流转。他伸出手指，凌空连续点出七下，每一下都精准地点在虚空某处，灌注的真元强度、属性、节奏都截然不同，仿佛在弹奏一首无声的乐章。

嗡——

那沉寂的古阵发出低沉的鸣响，灵光骤然明亮了一瞬，随即如同冰雪消融般，从叶牧谦点中的七个点位开始，迅速向内坍缩、消散。几个呼吸间，阵法彻底瓦解，露出了被其保护的真实景象。

石碑依旧，但在石碑后方，原本是坚硬岩壁的地方，赫然出现了一座高约一丈、宽五尺的厚重石门！石门表面粗糙，与山岩颜色一致，紧闭着，门缝间落满了灰尘和枯叶，显然已不知多少年未曾开启。

“开了！”柳成低声惊呼。

叶牧谦上前，与白言冰合力，缓缓推开沉重的石门。石门发出沉闷的“嘎吱”声，一股混合着陈旧尘土与淡淡矿物气息的空气涌出。

门后，是一条向下延伸数丈的简短石阶，尽头是一间颇为宽敞的石室。月光从门口斜斜照入，勉强照亮内部轮廓。地上积了厚厚一层灰，空气凝滞，果然是很久很久无人来过了。

四人鱼贯而入，白言冰取出几枚照明用的萤石，柔和的光芒散开，照亮了石室全貌。

石室约有三四间寻常屋子大小，布局分明。靠近入口一侧，摆放着石桌、石凳，还有几个空置的、已经大半腐朽的木架，墙上挂着一些磨损严重的工具，依稀可辨是锤、钳、凿等物，这里似乎是工作区。另一侧则有石床、倒塌的石柜、甚至还有一个简易的石灶台痕迹，应是生活区。

最引人注目的是墙壁。许多壁面上，用某种矿物颜料绘制着一幅幅壁画，虽因年代久远而色彩斑驳，但依然能看出大致内容：有人正在烈焰熊熊的炉前锻打胚料，有人在对着一件件法器刻画阵纹，有人则在测试各种精巧机关的性能……画中人物衣着古朴，样式与现今灵界流行的宽袍大袖或劲装短打皆不相同，更显简洁利落。壁画笔法生动，将炼器时的专注、成功时的喜悦、失败时的沉思都描绘得栩栩如生。

在工作区角落，堆放着一些矿石，在萤光下闪烁着内敛的光泽。柳成拾起一块暗青色、入手沉重的矿石，仔细辨认，动容道：“这些可都是炼制高阶法器的珍稀材料，品相极佳！”

石桌上，散落着几件炼器工具，如刻刀、模具、灵笔等，看似普通，但材质特殊，历经岁月而不朽，显然非是凡品。更引人注目的是旁边一个石架上，整齐摆放着三件形态各异的法器：一柄无鞘长剑，一口三足小鼎，一面八角铜镜。这些法器灵光内蕴，虽蒙尘埃，却依旧能感受到其不凡的波动。

叶牧谦拿起那把剑，稍一感应，心中微惊。剑身内蕴含的阵法结构与灵力脉络精妙异常，远超流云坊市所见，甚至比他之前拆解的徐峰等人的法器还要高明一个层次！按照灵界的品阶划分，这些法器，恐怕达到了天阶！只是因为久置无人温养，有些许灵性沉寂，但稍加修缮，绝对能恢复威力。

石桌中央，则静静地躺着三枚玉简，与柳成当年捡到的那枚材质相似，但保存更为完好。

叶牧谦拿起一枚，神识探入。大量信息涌入脑海，是一套系统、深刻、充满独特见解的炼器理论总纲！其核心思想，强调“器为心用，物性通灵”，主张炼器者需深刻理解材料本性，以自身心意引导灵力，达到“人器合一”，而非简单粗暴地堆砌材料与阵纹。

其中对于如何简化结构、提升效能、赋予器物独特灵性有着独到的阐述。这套理论，名为——《焚心器道》。末尾的署名，是三个清秀却力透“纸”背的古字：风灵希。

柳成也凑过来，叶牧谦将玉简递给他感应。柳成只看了开头一小部分，便激动得浑身颤抖：“是它！是它！就是这个感觉！我当年捡到的玉简内容零碎，但很多思路和这里一脉相承！能串联起来！原来……原来我这些年琢磨的东西，源头在这里！”

找到了传承源头，柳成又是欣喜又是感慨。几人将另外两枚玉简也查看一番，一枚记载了些《焚心器道》的具体应用实例和独门炼器手法，另一枚则像是一些随笔心得和未完成的构思草图，同样价值不菲。

收获巨大，但探索尚未结束。工作区已查看完毕，众人的目光投向了生活区，尤其是那石床和倒塌的石柜之后，似乎还有更深的角落。

他们绕过石床，后方空间稍小，陈设更显生活化。有半面残破的梳妆铜镜倒在尘埃里，旁边还有一个打开的空首饰盒。最令人心头发紧的是，在房间中央，安静地停放着一具石棺。

棺木朴素，无任何雕饰，与石室浑然一体。看到棺木的瞬间，四人都不由自主地停下了脚步，原本因发现宝藏而有些兴奋的心情，迅速沉淀下来，转为肃穆。

叶牧谦轻叹一声：“这里原先是一位炼器大师的隐居之所。看来，当她感到大限将至时，便自己踏入这预先备好的棺木，了结尘世一生。此地，也就从一间隐居的石室，变成了她的长眠之所——墓室。”

此言一出，白言冰、沈瑶、柳成皆感恍然，随即心生敬意与一丝打扰逝者的不安。他们并非为盗墓而来，此番闯入，终究是叨扰了前辈清静。

“既入宝山，得承遗泽，虽是无心，亦当心存感恩与敬意。”叶牧谦正色道。他率先走到石棺前三尺处，整理衣冠，躬身深深一礼。白言冰、沈瑶、柳成见状，也连忙跟随，郑重行礼。

行礼既毕，心中那份不安稍减。

此地既是墓室，自然不能再作为可能的道场考虑了。叶牧谦环视一周，将工作区的那些天阶法器、炼器用物品、珍稀矿石、以及三枚记载《焚心器道》的玉简小心收起。生活区的私人物品，包括那梳妆镜和首饰盒，他们丝毫未动，保持原样。

“风灵希前辈，我等无意冒犯，今日取走您生前心血所聚之道统与作品，必不让其蒙尘。打扰之处，还请见谅。”叶牧谦对着石棺再次轻声说了一句，然后示意众人退出。

四人带着沉重又收获满满的心情，退出石室，重新来到落星谷的夜空下。叶牧谦挥手，再次激发那隐匿阵法的残存灵机，令石门缓缓合拢，阵法灵光重新浮现、黯淡，最终恢复成那片看似寻常的乱石景象，将一代炼器大师的安眠之地重新掩藏。

\vspace{2em}

他们刚刚走出不到百丈，尚未离开落星谷干涸的河床，前方和侧翼的乱石阴影中，便窸窸窣窣地闪出五六道身影，迅速呈半包围之势逼近。为首之人，正是白天去“请”柳成未果的护卫李辉！他身旁，多了两名气息同样在真符境显化期的同伴，还有三名注灵期大成的修士。

显然，李辉回去搬了救兵，并且不知用何种方法，竟然跟到了这偏僻的落星谷！

李辉看到叶牧谦四人，尤其是被护在中间的柳成，眼中闪过一丝狠色与贪婪。他白天被叶牧谦气势所慑，回去禀报后，郑天河公子大为不悦，命令他必须将柳成带回去，且要查清叶牧谦等人的底细和目的。此刻见对方从这荒谷中走出，身上似乎还带着探寻遗迹后的风尘与隐隐的收获气息，他立刻疑心大起。

“站住！”李辉厉声喝道，目光扫过叶牧谦四人，最后定格在叶牧谦身上，“尔等鬼鬼祟祟在此荒谷作甚？可是发现了什么古传承遗迹？识相的，将所得之物与这柳成一并交出，可饶尔等不死！”

他试图以势压人，同时给叶牧谦扣上“私探古迹”的帽子。

叶牧谦停下脚步，面色平静无波，甚至没有去看那几名呈包围之势的护卫，只是目光淡然地落在李辉脸上。

被这般无视，李辉心中恼怒更甚，正待再次威胁——

嗡！

一股无形无质、却又浩瀚如星空、沉凝如大地的奇异领域，毫无征兆地以叶牧谦为中心，蓦然扩散开来！瞬间笼罩了方圆数十丈的范围，将李辉等六名护卫全然覆盖其中！

这不是灵力的威压，也不是神识的冲击，而是一种更为本质的、直指心灵与认知层面的“场”。李辉等人只觉得仿佛瞬间被投入了深不见底的寒潭，五感变得迟滞，思绪运转僵硬，体内灵力的流转都仿佛生锈的齿轮，变得艰涩无比！

他们张着嘴，想要呼喝，想要运转功法抵抗，却发现连舌头都不太听使唤，喉咙里只能发出“嗬嗬”的怪响，更别提做出有效的攻击或防御动作了！

通明境界强化过的见独领域！直指本心，映照灵台，扰乱认知！对付这些心境修为远不及他的修士，效果堪称碾压。

叶牧谦这才缓缓开口，声音不大，却清晰无比地传入每一个护卫混乱的脑海中：“回去告诉郑天河，柳成，现在是我的人了。” 他顿了一下，语气依旧平淡，却带着不容置疑的意味，“你们若真的要请他，还希望拿出诚意，派个能好好说话的人来。而不是像现在这样，鬼鬼祟祟，喊打喊杀。”

说完，他心念微动，收回了领域。

“噗通！”“噗通！”几名注灵期护卫直接瘫软在地，大口喘气，冷汗瞬间湿透衣背。李辉和另外两名显化期护卫也是脸色煞白，踉跄后退几步，看向叶牧谦的眼神充满了无边的恐惧。刚才那一瞬间，他们感觉自己就像狂风暴雨中的蝼蚁，生死完全不由自己掌控！

这到底是什么手段？闻所未闻！

“滚吧。”叶牧谦吐出两个字。

李辉如蒙大赦，哪里还敢有半分停留和废话？他连狠话都不敢撂，对着两名真符境同伴使了个眼色，自己率先带上一个瘫软在地的护卫，随即转身，连滚爬带遁光，仓皇无比地向谷外逃去。另外两人也连忙跟上，一人扛了一个，甚至没敢回头看一眼。

\vspace{2em}

转眼间，落星谷内只剩下叶牧谦一行四人。

叶牧谦不再理会他们，对白言冰三人道：“此地不宜久留，郑天河不会轻易罢休。我们找个地方避避锋芒。”

几人由叶牧谦带领，在流云山脉更深处寻了一处隐蔽的山洞暂作栖身。洞内经叶牧谦简单布置了几个隐匿与警戒的禁制后，倒也成了一处临时的安身之所。白言冰、沈瑶则出去寻回陈千羽，把她也带回了山洞里面。

惊魂甫定的柳成，一路上看着叶牧谦的背影，心中激荡难平。先是坊市摊前的解围，再是落星谷外的震慑强敌，最后更是带他找到了自己炼器技艺的真正源头，获得了难以想象的传承。这位叶前辈，实力深不可测，待人却无居高临下之感，更有一种引导人看见更广阔天地的智慧与气度。

当众人安顿下来，篝火燃起，驱散了山洞中的寒意与晦暗时，柳成再也按捺不住。

他走到叶牧谦面前，撩起衣袍下摆，双膝跪地，以额触地，声音因激动而微微发颤：“前辈！晚辈柳成，蒙前辈多次相救，更得前辈指引，寻回技艺根源，得窥大道门径。前辈之恩，如再造之德。晚辈自知愚钝，出身微末，但恳请前辈收我为徒！晚辈愿追随前辈左右，研习大道，精进器艺，绝无二心！”

说罢，他伏地不起，静待裁决。

白言冰、陈千羽、沈瑶在一旁静静看着，他们能理解柳成此刻的心情。

当年他们拜师时，何尝不是被师尊那看似平凡却蕴含无限可能的身影所折服，看到了截然不同的道路？

叶牧谦看着跪伏在地的柳成，眼中流露出欣慰之色。他早已隐隐有感，柳成便是他那最后的弟子缘法。此子心思纯粹，于炼器一道天赋独特，更难得的是那份在底层挣扎中仍未磨灭的创造力与对技艺本身的热爱。如今，水到渠成。

“起来吧。”叶牧谦伸手虚扶，一股柔和的力量将柳成托起，“你既诚心向道，于炼器亦有宿慧，我便收你为徒。你当为为师座下第四弟子。”

柳成闻言，大喜过望，眼圈瞬间红了，连忙再次拜下：“弟子柳成，拜见师尊！多谢师尊成全！”

“恭喜师弟！”白言冰三人也笑着上前道贺，山洞中气氛顿时温馨起来。陈千羽细心，早已备好了简单的拜师茶，柳成恭恭敬敬地奉上，叶牧谦含笑接过饮下，这师徒名分便算正式定下。

待礼毕，叶牧谦将此次落星谷所得尽数取出。莹莹宝光顿时照亮了山洞。

他首先拿起那柄天阶长剑。剑身无鞘，长三尺有余，通体呈现一种暗沉的青灰色，似金非金，似玉非玉，入手极沉。剑刃看似无锋，但凝神细观，能见其边缘有细微的空间扭曲之感，仿佛能无声切开一切阻碍。剑身内部阵纹层层嵌套，核心处一点灵光如星子般沉寂，却蕴含着令人心悸的锋锐与破法之力。

“言冰，你剑心通明，此剑锋锐内蕴，破法诛邪，正合你用。”叶牧谦将长剑递予白言冰。

白言冰双手接过，稍一感应，便觉剑中灵性与自己通明真元隐隐相合，仿佛沉寂多年，终遇明主。他肃然道：“谢师尊赐剑！弟子必善用此剑，护道卫正。”

接着是那口三足小鼎。鼎身古朴，呈暗金色，其上天然纹路似云霞缭绕，又似草木生发。鼎腹圆润，三足稳健，虽未催动，已有一股沉凝厚重、滋养万物的气息散发开来。这东西虽然偏向于温养、淬炼、调和，但它也被刻意设计成可以轻易搬起并投掷出去的形状，甚至倒扣时还可以用于护体。

“千羽，你精研医药，仁心济世。此鼎有温养灵粹、调和药性、祛除杂质之能，于你炼丹行医大有裨益。而且看这设计，危急时刻，还能用来砸人与护身。”叶牧谦将小鼎交给陈千羽。

陈千羽轻轻抚摸鼎身，感受其中蕴含的勃勃生机与宁静之力，眼中泛起喜色：“多谢师尊！此鼎确为医道良助。”

然后是那面八角铜镜。镜背雕刻着繁复的星辰与流云图案，镜面澄澈如水，却又仿佛深不见底，倒映人影时略有扭曲，似乎能照见人心底的某些层面。这是一件偏向幻术、洞察、乃至部分防护的法器。

“瑶儿，你擅幻术，心思机敏。此镜有映照虚实、迷惑灵觉、偏转攻击之效，与你千幻流光之术相得益彰。”叶牧谦将铜镜递向沈瑶。

沈瑶接过，指尖拂过冰凉的镜背，镜面微光流转，与她体内真元产生微妙共鸣，似乎能将她幻术的威力与变化提升一个层次。“谢师尊。”她简短回应，眼中却有精光闪过，显然已在思考如何运用此镜。

最后，叶牧谦指向那堆从石室工作区带回的、品质极佳的炼器矿石，以及风灵希留下的那些明显达到道阶（对应道种境）乃至天阶层次的专用刻刀、灵笔、模具等工具。

“柳成，你既入我门下，专研器道。这些矿石与工具，便都交予你。望你善用先贤遗泽，结合自身巧思，在炼器一道上走出自己的路。”叶牧谦对柳成说道，语气中充满期许。

柳成看着那堆散发着各色宝光、灵气盎然的珍稀矿石，以及那些造型古朴、气息玄妙的炼器工具，激动得手足无措。这些都是他以前做梦都不敢想的顶级材料与工具！“弟子……弟子定不负师尊厚望！必刻苦钻研，将风前辈的传承发扬光大！”

分配完毕，众人各自熟悉新得的法器。叶牧谦则与柳成一起，开始深入研究那三枚记载着《焚心器道》的玉简。

越是深入研读，叶牧谦眉头蹙得越紧。这套炼器理论高屋建瓴，直指“人器合一”“物性通灵”的境界，其核心思想与许多精妙构想，连他也觉受益匪浅。

然而，其具体的实操法门，尤其是涉及核心炼制过程的部分，却存在一个几乎无法逾越的门槛——对灵力的质与量要求极高！

按照玉简中风灵希的记述与暗示，想要真正实践《焚心器道》的高深法门，炼制出拥有独特灵性甚至成长潜力的器物，至少需要天轮境级别的浩瀚灵力作为支撑，用以完成最后也是最关键的“赋灵”与“启慧”步骤。这对于目前最高才通明境的师徒几人来说，无异于镜花水月。叶牧谦隐隐觉得要想完全复刻这一步骤，至少也得达到凝一境界甚至更高。

更让叶牧谦感到心惊的是，风灵希在天轮境之前，用以替代高阶灵力、完成部分核心锤炼与塑形的方法——引燃自身心念，化念为火，以这“心念之火”千锤百炼器胚！

玉简中的描述看似平淡：“心念微动，火自内生，不假外物，锤炼由心。”但叶牧谦通明境界大成，对心念之力理解极深，岂会不知其中凶险？

所谓“引燃心念”，几乎等同于燃烧神魂！每一次“焚心”炼器，都是对自身心神的一次残酷拷问与损耗。那过程中的痛苦，绝非文字可以描述万一，恐怕足以让意志不坚者心神崩溃，或留下永久性损伤。

风灵希能以平淡笔触记录，要么是因为那无尽的痛苦早已在她无数次实践中被硬生生“磨平”了感知，或是她心志坚毅远超常人。

这“焚心”之名，实至名归，透着一种决绝与残酷的美。

“此法……对修炼者心志要求太高，且隐患不小。”叶牧谦放下玉简，沉吟道，“尤其对你而言，柳成，你初入道途，心念修为尚浅，强用此法，无异于自毁长城。”

柳成也看得脸色发白，他虽渴望高深技艺，但也知性命与根基更重要。

叶牧谦闭目沉思。火光在他平静的脸上跳跃。《焚心器道》的核心在于以心御器，强调炼器者自身心意、精神与器物的深度沟通与融合。

而“问道途”，恰恰是最重修心炼性的一途！

从最初的淬体体悟自身，到养气调和阴阳，再到见独明心见性、通明映照万物，无不是对内心世界与精神力量的深入挖掘与掌控。

“或许……并非没有转圜余地。”叶牧谦睁开眼，眸中闪过一丝明悟，“《焚心器道》需以心念为火，是因缺乏更高阶的能量驱动，只得行此险着，透支本源。而我问道途，修的是自身元气，养的是通明心念。心念之力，未必非要引燃才能作用于外物……”

他想到自己的通明领域，可以影响他人心神，感知万物细微。若将这种对心念的精微操控，与炼器过程中对材料物性、阵法灵机的感知引导相结合，是否能够以更温和、更持续的方式，达到类似“心念之火”锤炼、沟通、甚至赋灵的效果？

无需燃烧，而是引导、浸润、共鸣？

“问道途与《焚心器道》，或有融合互补之妙。”叶牧谦看向柳成，目光灼灼，“不过，这需要对你所修功法进行调整，至少让你在养气境时，就能初步驾驭这种心念浸润之法，而非焚心自残。此事需从长计议，我先尝试改良其中关窍。当务之急，是你需先打下问道途的根基。”

柳成重重点头：“弟子明白！全凭师尊安排！”

转修之事刻不容缓。在沈瑶的护法、陈千羽的丹药帮助下，柳成于山洞深处盘膝而坐，开始按照叶牧谦传授的入门法诀，尝试破开原有的灵枢径根基，凝聚问道途的气海。

过程比他想象的更为痛苦。灵枢径的灵力早已与肉身有了较深的结合，强行逆转散功，如同将生长在血肉中的根系一点点剥离，剧痛钻心。

柳成额上青筋暴起，冷汗瞬间湿透衣衫，但他咬牙硬撑，想起自己卑微的过去，想起师尊的期许，想起那精妙绝伦的《焚心器道》，心中有一股火在烧——那不是自焚的心火，而是不甘平庸、渴望攀登的斗志之火。

陈千羽在一旁，随时准备以银针和金针度穴之法帮他疏导紊乱的气息，缓解痛苦。沈瑶则警惕地守护在外围，防止任何意外干扰。

\vspace{2em}

与此同时，郑氏祖地。

郑天河于自己的书房内听取李辉等人狼狈逃回、添油加醋的汇报后，并未如手下预料的那般暴怒。

“叶牧谦……落星谷……”他低声咀嚼着这两个名字。护卫描述的那种无形领域，能无声无息压制真符境修士，让其心中警兆狂鸣却无从反抗，这绝非寻常手段，至少也得是法相境的修为。

是某种强大的灵魂秘术？

还是借助了了不得的法宝？

对方实力深浅不明，在郑氏地盘上却似毫无根脚背景；偏偏又能从落星谷那等荒僻的地方中，找到一处古代石室，得了某种尘封已久的古传承。

硬碰硬，不明智。对方已经有了警觉，自己若再派更强硬的力量前去，且不说能否成功，万一踢到铁板，或者动静闹大，被家族里那些盯着自己的老家伙们抓住把柄，反而不美。

但就此放弃？那可能蕴含巨大利益的留影镜技术源头，以及其背后可能更珍贵的、能令一个名不见经传的小子短时间内炼制出新颖法器的古传承？绝无可能。

“风浪越大，鱼越贵……”郑天河嘴角勾起一抹冰冷的弧度，那笑容里没有半分暖意，只有商贾评估风险与收益时的精明，以及猛兽蛰伏等待时机时的残忍。

\vspace{2em}

不过数日功夫，郑天河精心炮制的流言，便如同滴入清水中的墨汁，以流云坊市为中心，迅速在东域修士间晕染扩散开来，尤其精准地渗透进了炼器师和那些喜好探寻古迹、撞仙缘的散修圈子。

流云坊市的茶楼。

这是消息灵通的散修们最爱的聚集地之一。二楼临窗的座位，几名衣着各异、气息驳杂的修士正压低声音交谈，脸上都带着某种发现秘密的兴奋与不平。

“听说了吗？流云山脉那边的落星谷，前些日子可是出了大事！”一个瘦削的汉子抿了口灵茶，神秘兮兮地开口。

旁边一位满脸络腮胡的大汉立刻接话：“何止听说！据说是被几个不知从哪冒出来的外来修士给独吞了！一点汤都没给别人留！” 他声音不自觉地提高了几分，引来旁桌几道关注的目光。

“当真？”一个书生模样的青年修士皱眉，“古传承乃无主之物，见者有份，他们凭什么独吞？”

瘦削汉子嗤笑一声：“凭什么？凭拳头硬呗！听说那几人实力不弱，古怪得很，连郑氏的人前去交涉都吃了瘪，灰头土脸地回来了。”

“郑氏都吃了瘪？”书生修士倒吸一口凉气，“那这几人来头不小啊……不过再厉害，吃独食也不对吧？咱们东域的规矩，历来是宝物出世，有缘者得之，但也讲究个分润机缘。何况还是炼器传承，”他压低了声音，眼中闪过热切，“说不定里面真有能改变现今炼器格局的秘法呢！”

类似的对话，在坊市的炼器材料铺子前、在散修摆摊的集市角落、甚至在几个小宗门弟子聚集的酒肆里，同时或先后上演着。

如同生了翅膀，在口耳相传中被不断增添细节，越发活灵活现。虽然始终没有指名道姓说是“叶牧谦师徒”，但时间、地点、人物特征都对得上，很快便有不少消息灵通或自以为聪明的修士，将饱含各种意味的目光投向了流云山脉方向。

这刻意引导的舆论风潮，其影响很快便超出了郑氏领地的范围。东域其他几个与郑氏规模相若的势力、一些以炼器闻名的小宗门，乃至一些游历四方、消息灵通的其他地域修士，也渐渐听闻了“落星谷古炼器传承被神秘外来者独占”的消息。这引来了一个地位超然、能量特殊的势力的注意——天工阁。

天工阁，并非传统的宗门或世家，而是一个由众多痴迷于炼器、阵法、机关傀儡等技艺的匠师、大师自发组成的联盟式组织。他们或许个体战斗力不算修真界的顶尖，但汇聚的财富、交织的人脉网络、以及掌握的各类稀缺资源与核心技术，使天工阁在灵界威望极高；更关键的是，天界也有一座天工阁，与灵界天工阁本是一门分设两地，所以不少大型势力也不敢小觑。对于可能改变炼器格局的古传承，天工阁自然抱有极高的兴趣和敏感度。

很快，阁中一位级别极高、素以眼光毒辣和处事老练著称的阁老——吴毓，便接到了上层示意，开始暗中调查这则突然甚嚣尘上的流言真伪，并密切留意那几个“神秘外来者”的动向。

\vspace{2em}

与此同时，流云山脉深处，隐蔽的山洞内。

篝火跳跃，将洞壁上的影子拉得忽长忽短。柳成盘坐在最内侧，周身气息极其不稳定，时而剧烈波动，时而微弱如缕，脸色在金红火光映照下显得愈发苍白，冷汗早已浸透数次更换的衣衫。陈千羽盘坐在他侧后方，指尖捻着数根细如牛毛的银针，全神贯注地感受着柳成体内灵力流转的每一丝异常，随时准备下针疏导。白言冰抱着剑，倚靠在洞口附近的岩壁，闭目养神，但全身肌肉微微绷紧，灵觉如同蛛网般细细铺开，笼罩着山洞周边数十丈的范围。

沈瑶的身影如同轻烟般从洞外飘入，带进一丝夜晚山林特有的清冽气息。她先是对白言冰微微颔首，示意外围无恙，然后快步走到篝火旁，对正在凝视火焰、仿佛在思考着什么的叶牧谦躬身一礼。

“师尊，弟子回来了。”

叶牧谦抬眼，火光在他深邃的眸中跃动：“外面情况如何？”

沈瑶眉头微蹙，语气带着凝重：“正如师尊所料，不过数日功夫，外面已经传得沸沸扬扬。流言焦点正是落星谷古传承被我们‘独吞’，描述绘声绘色，虽未直指我们，但特征指向再明显不过。如今已有不少修士在打听流云山脉这边的动静，心怀叵测者恐怕不在少数。”她顿了顿，补充道，“还有一个值得注意的消息，那股流言似乎也引起了一个叫‘天工阁’的势力的注意，他们可能已经派人暗中调查了。”

叶牧谦用一根细长的树枝轻轻拨弄了一下篝火，让几颗火星爆开，升腾而起，又迅速湮灭在黑暗中。

“意料之中。”他淡淡道，“看来，郑天河很懂得算计啊。此举一石多鸟。散布谣言，首先能给我们制造巨大的麻烦，吸引无数贪婪的目光聚焦于此，让我们难以安心消化所得，拖延我们的时间；其次，可以借这些被流言鼓动起来的有心人，来试探我们的虚实底牌，无论谁胜谁负，他都能躲在幕后观察，收集信息；最后，局势若真乱起来，他郑氏作为地主，或许还能以调停者、仲裁者的身份介入，趁机攫取利益。”

他顿了顿，目光扫过几名弟子，语气平和却带着一种看透世情的睿智：“当舆论不利于己时，急赤白脸地辩解、反驳，与每一个传播流言者争辩，往往只会适得其反，越描越黑。”

“那我们该如何应对？总不能任由他们污蔑，坐以待毙吧？”柳成虚弱的声音从内侧传来。

叶牧谦看向柳成，眼神微缓，点了点头道：“能有此心，不错。但应对之道，未必在口舌之争，亦未必在拳脚之快。”他的目光转向沈瑶，眼中流露出信任与托付，“瑶儿，你精于幻化伪装，心思缜密多变，最擅长与人周旋，洞察人心。为师有一事，需你冒险为之。”

沈瑶立刻挺直脊背，躬身道：“师尊请吩咐。弟子万死不辞。”

“郑天河能散布谣言、搅动外部风云，我们亦可暗中布局，从其内部着手。”叶牧谦的目光变得幽深，仿佛能穿透岩壁，看到远方那座繁华又复杂的流云坊市。

“郑氏家族枝繁叶茂，内部绝非铁板一块。利益纠葛，派系倾轧，失意之人，怀才不遇者，所在多有。我需要你在郑氏内部，设法埋下一颗钉子，一个能在关键时刻为我们提供信息、或制造些许便利的暗桩。此人不必身居核心，但需有一定地位和消息渠道，更重要的是，要有能被打动的弱点或诉求。此事需极度谨慎，人选、接触方式、分寸拿捏，皆需你自行斟酌把握。”

沈瑶眼中先是闪过一丝惊讶，随即化为强烈的了然与隐隐的兴奋。

这种潜入敌营、揣摩人心、策反布局的暗战任务，正完美契合她的天赋与心性，她感觉在内域沉寂七年的血液似乎都微微发热起来。

“弟子明白！定当小心行事，不负师尊所托。”

\vspace{2em}

接下任务后，沈瑶先是花费了一天时间，仔细回忆和梳理之前数次进出流云坊市时有意无意观察到的信息，又结合叶牧谦对郑氏内部权力结构的简要分析，在心中勾勒了几个可能的方向。

次日，她便彻底改头换面。幻术微光流转，她高挑的身形微微佝偻了几分，面容变得平凡而略显沧桑，眼角添了细纹，穿着一身半新不旧、沾着些许风尘的褐色短打，背上一个鼓鼓囊囊的褡裢，腰间挂满各种稀奇古怪的小囊、皮袋，摇身一变，成了一位游走四方、见识颇广的行脚商人沈三。

这类行脚货郎在修士界底层很常见，混迹于市井，消息灵通，又不引人注目。

她刻意在流云坊市及周边几处由郑氏直接掌控或影响力较大的产业区域出没，有时在茶楼角落独坐饮茶，耳朵却捕捉着每一段对话；有时在集市摆摊，与往来客人讨价还价，实则观察着哪些郑氏中人会对自己摊位上那些带着古朴气息的破烂多看一眼；有时则“恰好”与一些郑氏中下层管事、护卫同桌吃饭，几杯酒下肚，便能听到不少牢骚和零碎八卦。

很快，一个目标清晰地进入了她的视线——郑氏长老，郑季。

郑季，法相境通灵期修为，在郑氏众多长老中，资历不算浅，但地位却颇为尴尬，甚至有些边缘。

据沈瑶从各方听来的零碎信息拼凑：约二百年前，郑季曾是家族中风光无限的精英后裔，天赋、心性、能力皆属上乘，一度与现任家主郑伯——郑天河之父——在继承人之位上争得不可开交。然而，在最后关头，他不知什么原因，最终棋差一着，与家主大位失之交臂。自那以后，胜利者郑伯一系自然稳坐核心，而郑季则被逐渐排挤出家族的核心决策圈，被打发来经营流云坊市这类虽然油水丰厚、但更多属于“边缘”的产业，美其名曰“发挥其经营之长”，实则是流放，远离了真正的权力中心。

二百年来，郑季心中那份郁结与不甘，可想而知。他对家族现任核心层，尤其是家主郑伯一系，不可能没有怨怼。

更重要的是，沈瑶打听到，郑季本人亦是一位炼器师，造诣不算顶尖，但颇有钻研精神，只因家族资源向其倾斜有限，加之心中块垒难消，始终未能达到炼器宗师之境，然而其对炼器之道的那份热忱与追求，却似乎从未完全熄灭。

这简直是天赐的目标。

\vspace{2em}

沈瑶开始有意识地在郑季可能出现的场合制造偶遇。郑季有时会去坊市内几家信誉较好的炼器材料铺子转转，有时也会去茶楼独坐片刻，沈瑶也便开始在这些地方活动。

一次，在材料铺外，郑季正背着手，看着伙计搬动一块品相不错的原矿。沈瑶“恰好”背着褡裢经过，褡裢口没系紧，“不小心”掉出一件东西，咕噜噜滚到郑季脚边。

那是一个巴掌大小、形制古拙的铜环，表面布满铜绿，但依稀可见一些难以辨认的云纹。这只是沈瑶依照《焚心器道》中某段关于“基础灵纹回路共鸣”的粗浅描述，随手用边角料做的试验品，除了注入灵力后会微微发热、让附近的低阶金属产生极其微弱的共鸣震颤外，并无大用。但那股古朴的气息和略显奇特的波动，却立刻引起了郑季的注意。

“咦？”郑季俯身捡起铜环，指尖灵力微微探入，感受到那细微却独特的震颤反馈，眼中闪过一丝讶色。

沈瑶连忙小跑过来，一脸惶恐和市侩的笑容：“哎哟，这位老爷，实在对不住，对不住！小老儿没留心，脏了您的眼。”说着就要去接铜环。

郑季却避开她的手，仔细打量着铜环，问道：“此物……你从何处得来？”

“这个啊，”沈瑶搓着手，扮出回忆的样子，“是前些年小老儿在西边捡漏得来的。看着像个古物，但也没啥大用，就是偶尔摸着有点暖乎气，怪哉。老爷您要是感兴趣，便宜卖了，三个……不，两个下品灵石就行！”她故意将价格说得低廉，符合一个不识货行脚商的身份。

郑季没有还价，反而问：“你还见过类似这种纹路，或者感觉奇特的老物件吗？”

沈瑶心中暗喜，面上却露出为难之色：“这个……倒是在别的破烂里也见过一两次类似的纹路，都残缺不全了。小老儿走南闯北，就爱收点稀奇古怪的玩意儿，有时候也听那些挖坟……咳咳，探索古迹的道友们说些奇闻。老爷您是行家，对这些古旧东西有兴趣？”

郑季不置可否，却递过去几枚灵石：“这东西我买了。你若再有类似的发现，或者听到什么关于上古炼器流派、失传技艺的传闻，可以到坊市东头的茶楼留个信，就说找郑掌柜。”

“好嘞！多谢老爷！多谢老爷！”沈瑶点头哈腰，接过灵石，目送郑季离去，眼中精光一闪而过。

鱼儿，开始试探诱饵了。

\vspace{2em}

此后数日，沈瑶又“偶然”在茗香苑茶楼与郑季“巧遇”两次，不着痕迹地谈论一些从《焚心器道》外围理论中剥离出来的、关于材料物性相生相克、灵力引导须顺应材质本源的模糊见解。

这些见解与当今主流炼器手法强调的强力淬炼、精密控制颇有不同，更注重沟通与引导。虽不涉及核心，却让郑季听得心神震动，如同在封闭的房间里推开了一扇缝隙，窥见一丝截然不同的风景。

终于，在第二次茶楼“偶遇”时，郑季主动开口，邀请这位见识似乎颇为不俗的“行脚商人”到楼上他常包的一处僻静雅间细聊。

\vspace{2em}

雅间内，清茶袅袅，隔绝了楼下的喧嚣。沈瑶坐下后，灵力微不可察地流转，卸去了部分伪装，虽然面容依旧平凡，但那份刻意伪装的市侩与怯懦气质却褪去不少，多了几分沉静。

郑季亲自斟茶，目光审视着沈瑶：“阁下……并非寻常行商吧？前番所言，虽只鳞片爪，却深得炼器三昧，许多见解闻所未闻。不知阁下师承？”

沈瑶端起茶杯，轻嗅茶香，不急不缓道：“郑长老慧眼。在下沈三，确非专事行商。早年有些机缘，偶然得到部分古老的炼器师手札残篇，奈何学识浅薄，加之残篇缺损严重，许多精妙处如雾里看花，始终难以参详透彻。久闻郑长老不仅是郑氏栋梁，更是炼器一道的行家里手，浸淫此道数百载。在下冒昧接近，实是想以这些残篇愚见，向长老请教，或能互相印证，弥补缺憾。”

郑季果然被“古老炼器师手札残篇”这个说法牢牢吸引。他追问道：“不知阁下所得，是哪位先贤的手札？又可否让老夫一观残篇？”

沈瑶叹息一声，摇头道：“残篇以古篆书写，破损不堪，且似乎并非出自单一作者，更像是一些上古炼器理念的杂录合集。实体早已在一处险地中彻底损毁，幸而其中部分内容被我强行记下。”她看着郑季眼中闪过的急切与渴望，知道火候渐至。

于是，她将《焚心器道》中一些非核心、但确实精妙的基础理论，以缓慢而清晰的语言，讲述了几个要点；却又在关键处语焉不详，留下大量想象和探究的空间。

郑季听得如痴如醉，这些理论与他毕生所学碰撞，激发出无数灵感火花。他仿佛一个在沙漠中跋涉已久的旅人，骤然看到远处模糊的绿洲轮廓，那种震撼与渴望，几乎要满溢出来。

“妙！太妙了！竟有如此思路！以心念为桥，以灵润代火淬……这，这真是……”郑季激动得手指微微发抖，他紧紧盯着沈瑶，“阁下……沈先生，您所记下的，定然只是这传承的冰山一角！这绝非普通古法，其立意之高，恐怕直指炼器大道本源！那手札……不，那传承，究竟……”他呼吸急促，眼中闪烁着近乎狂热的求知光芒。

沈瑶见火候已到，知道是摊牌引线的时刻了。

她脸上的平和神色稍稍收敛，取而代之的是一种深沉的叹息，仿佛承载着无奈与考量。她放下茶杯，目光平静地迎向郑季：“不瞒郑长老，在下亦知这些残篇价值非凡，奈何自身能力有限，福缘浅薄，难以尽窥堂奥。此次前来流云坊市，除了向长老请教，其实……亦另有一事相察。”

郑季从激动的情绪中稍稍冷静，敏锐地察觉到对方话中有话：“何事？”

沈瑶缓缓道：“在下亦听闻，郑长老在族中，德才兼备，却似乎……有些抱负难展？空有炼器之志，却囿于方寸之地。”

郑季脸色骤然一变，眼中的热切迅速被警惕和一丝被触及痛处的冷意取代。

他放下茶杯，身体微微后靠，属于法相境修士的淡淡威压无声弥漫开来，笼罩住小小的雅间。他沉声道：“阁下此言何意？莫非打听我郑氏族内之事，才是你的真正目的？”

房间内的空气仿佛凝滞了，茶香中混杂了一丝紧绷的危险气息。

沈瑶却恍若未觉那无形的压力，神色坦然依旧，甚至带着几分坦诚的锐利：“郑长老不必动怒，更无须怀疑在下的用心。明人不说暗话。”她微微前倾身体，声音压低，却字字清晰，“我知郑长老因二百年前旧事，心有不平之气。更知长老因家族资源所限，纵有凌云之志，在炼器大道上却难有根本寸进，瓶颈困锁已久。长老难道甘心，此生便如此蹉跎于这坊市之间，任凭胸中所学与心中抱负，随岁月一同蒙尘？”

郑季脸色变幻，袖中的手悄然握紧。这些话，句句戳中他心底最隐秘的痛处与不甘。

沈瑶继续道，语气带着一种蛊惑性的力量：“若现在，有一条路，或许能让长老在炼器之道上窥见更高远的风景，触及那梦寐以求的宗师乃至更高境界……甚至，或许能借这份力，在家族内重新获得某些话语权，不再被人轻易遗忘于边缘。长老您，可愿……一试？”

“你……”郑季霍然站起，周身灵力波动了一下，又强行压下。他脸色苍白中透着红潮，眼神剧烈挣扎，死死盯着沈瑶，“你……你究竟是谁？代表何方势力？”

他终于意识到，眼前之人绝非偶然，这是一场精心设计的接触！

“我是谁，代表谁，眼下并不重要。郑长老可称呼我为沈姑娘。”沈瑶也站起身来，恢复了原本清越的女声，虽然面容未改，但气质已然迥异。

“重要的是，我手中确实掌握着更完整的、源自上古某位炼器大能的传承线索。而我的‘主人’，对郑氏并无恶意，更无意与郑长老为敌。只是眼下，我主被人以阴损谣言中伤，引得群狼环伺，不得安宁。我们需要的，只是在贵家族中，寻一位可靠的、有分量的‘朋友’，在必要时，能互通有无，行个方便，仅此而已。”

当说到“主人”两字时，沈瑶的声音有些颤，但掩饰得很好，正狐疑间的郑季并没有听出来。

说着，她取出一枚早已准备好的玉简，轻轻推到郑季面前的桌上。玉简质地普通，但上面流转着隔绝探查的微光。“此乃定金，亦是诚意。其中记载了方才所谈理论的一部分，以及一种上古‘灵纹初解’的方法，虽非核心，却足以让长老验证其价值。若长老有意合作，日后自有更多深入交流。若长老觉得风险太大，无意卷入，那么，”沈瑶做了个请的手势，“便请长老当今日从未见过沈三，这枚玉简，也可当场毁去。”

郑季的目光死死黏在那枚玉简上。

理智在疯狂呐喊危险，警告他这是与虎谋皮，一旦踏入，便是万丈深渊。但内心深处，那沉寂了二百年的对炼器更高技术的渴望，那份被排挤打压的不甘与怨愤，以及玉简中可能记载的、能让他脱胎换骨的知识……如同无数只蚂蚁，啃噬着他的心防。

他颤抖着伸出手，拿起玉简，神识小心翼翼地探入。

刹那间，一股远比方才听闻更加系统、更加深邃、更加玄奥的炼器理念洪流，涌入他的识海。那是对“器”的本质的全新诠释，是对“灵”与“物”关系的颠覆性认知，虽然依旧残缺，没有具体的高深法门，但仅仅这基础的框架和方向，就已让他如遭雷击，浑身颤栗，仿佛看到了一个全新的、浩瀚无垠的炼器世界在眼前展开。

是真的！这绝非杜撰！是真正上古失传的宝物！价值无可估量！

巨大的诱惑如同最甜美的毒酒，而现实处境的不公与冷漠，则成了推动他饮下这杯毒酒的最后一把力。风险？与可能触及大道、一雪前耻的机会相比，风险又算得了什么？家族？家族何曾真正公平地待过他郑季？

他脸色苍白如纸，又泛起激动的潮红，额角渗出细密的汗珠，握着玉简的手青筋暴起。沉默了足足一炷香的时间，雅间内只闻两人轻微的呼吸声，以及郑季内心激烈交战的无声轰鸣。

最终，他像是被抽干了所有力气，又像是下定了毕生最大的决心，嗓音干涩沙哑地开口，那声音里充满了复杂的情绪：“因为家族……因为某些原因，我此生……大概率已无缘道种境了。所求无他，只是……只是希望能够在炼器之道上，真正有所建树，炼出一两件……能传世、能让后人记得我郑季名字的好法器，不枉此生修行，不被后世完全遗忘罢了。”

这话语中的落寞、不甘、以及一丝绝望中的希冀，无比真切。

沈瑶静静听着，知道对方的心防已经出现了决定性的松动，但还差最后一把火，需要他自己做出选择，强扭的瓜不甜。

“郑长老的心思，我明白了。”沈瑶将玉简收起，语气平和，“此事关乎重大，长老确实需要时间思量。”她走到窗边，看了看天色，“明日，还是此时，此地。我静候长老佳音。”

说完，她身形微微一晃，如同水中的倒影被石子打破，泛起一阵轻微的涟漪，整个人便融入了雅间角落的阴影之中，气息瞬间消失得无影无踪，仿佛从未出现过。

雅间内，只剩下郑季一人，呆呆地站在原地，望着沈瑶消失的方向，又低头看了看自己空空如也、却仿佛残留着玉简温热触感的手掌。脸色在烛火下变幻不定，时而有获得至宝线索的狂喜，时而有踏入未知险境的恐惧，时而有对家族压抑已久的愤懑，时而有对未来一丝渺茫期待的微光。

他知道，从这一刻起，自己站在了一个可能彻底改变命运，也可能万劫不复的岔路口。

\vspace{2em}

郑季恍恍惚惚地回到自己在流云坊市的宅邸。他挥手屏退了所有上前问候的侍从与管事，独自一人穿行在回廊与庭院间，脚步声在空旷的青石板上回荡，显得有些寂寥。他径直走进了自己的书房，厚重的门扉在身后无声合拢，隔绝了外界的一切。

书房内，陈设雅致，博古架上摆放着不少与炼器相关的古籍和矿石样品。郑季没有点灯，任由暮色如潮水般涌入室内，将他渐渐吞没。他在偌大的书房中来回踱步，双手负在身后，指节无意识地反复曲张。袖中那枚记载着惊世传承碎片的玉简，明明轻若无物，此刻却仿佛重若千钧，灼烧着他的肌肤，更灼烧着他的神魂。

“与那神秘莫测、正被郑天河视为眼中钉的势力合作……”

“成为暗桩，泄露家族内部动向……”

“交换那直指大道的上古炼器传承……”

这些念头如同沸水中的气泡，在他脑海中翻滚、炸裂。

风险、恐惧、对家族律法的敬畏，紧紧缠绕着他的心脏。

这份利益的诱惑有多大，背后隐藏的深渊就有多深。他几乎能预见，一旦事情败露，自己将被刻在家族的耻辱柱上，甚至可能永世不得翻身。

然而，另一股力量同样在激烈地冲撞着他的理智——那份窥见全新炼器天地的震撼与狂喜，那份困守边缘二百年的憋屈与不甘，那份渴望被重新看见、渴望在真正的大道上留下印记的炽热愿望。

退？退向何方？家族早已视他为弃子。进？前方是刀山火海，却也可能有毕生所求的风景。

他在黑暗中不知踱了多少圈，窗外的天色彻底黑透，星辰零星亮起。终于，他脚步一顿，深吸一口气，仿佛下定了某种决心，转身推门而出，向后院一处更为清静的花厅走去。

每当心中乱麻难理时，他总会去那里寻他的道侣，沈玉。

花厅建在一片精心打理的小灵植园中，夜风送来草木清气与淡淡花香。厅内只点了几盏柔和的暖光石灯，映照着一个窈窕的身影。

沈玉正站在一盆君子兰前，手持银剪，神情专注而温柔地修剪着稍显杂乱的兰叶。她身着素雅的月白长裙，长发松松挽起，几缕青丝垂在颊边，侧脸在灯光下显得温婉宁静，法相境巅峰的修为让她周身萦绕着淡淡的清灵之气，岁月似乎并未在她身上留下太多痕迹。

听到脚步声，沈玉并未回头，只是轻轻放下银剪，拿起一旁的细布擦了擦手，温声道：“回来了？看你心神不宁的，可是坊市中遇到了难事？”

郑季走到她身边，看着那盆在灵力滋养下微微发光的幽兰，沉默了片刻。他对这位夫人，不仅是爱侣，更是最信任的知己与智囊。二百五十年前，这位来自遥远下界、漂泊至东域的奇女子，以她的聪慧、见识和那份独特的温柔坚定，走进了他当时正值意气风发却也暗藏危机的人生。在他与郑伯争夺家主之位最激烈的时刻，是她为他剖析局势，出谋划策，是他最坚实的后盾；在他落败失意，被打发到这流云坊市后，也是她不离不弃，以她的方式和人脉，默默帮他经营着这份“边缘”产业，更以她的温柔与智慧，替他排解了无数烦忧。

在他心中，沈玉的分量，甚至超过了许多血脉相连的族人。

“玉儿，”郑季的声音有些干涩，“今日……我遇到一个人。不，更准确地说，是她找到了我。”

沈玉转过身，清澈的眼眸望向他，带着询问。

郑季再无隐瞒，将今日在茗香苑雅间发生的一切，原原本本、毫无保留地告诉了沈玉。从“行脚商人沈三”如何引起他注意，到雅间内对方卸去部分伪装，自称沈姑娘，展示传承碎片，提出合作（或者说，充当暗桩）的条件，再到那枚作为“定金”的玉简中蕴含的、令他神魂震撼的炼器理念，最后到对方给出的三日考虑期限以及悄然离去的神秘。

他只隐去了那沈姑娘最后提及的“我的主人正被谣言中伤”这一细节，因这几乎直接指向了近日沸沸扬扬的“落星谷事件”主角。

沈玉静静地听着，脸上温柔的神色渐渐变得沉静而专注。她走到一旁的藤椅边坐下，示意郑季也坐下，自己则微微垂眸，陷入了沉思。修剪整齐的指甲无意识地轻轻叩击着藤椅的扶手，发出极细微的“嗒、嗒”声。

良久，沈玉抬起眼眸，目光清亮而锐利，仿佛能穿透迷雾。“阿季，”她缓缓开口，声音依旧柔和，却字字清晰，“此事风险极大，毋庸置疑，可谓与虎谋皮。但……正如你方才所言，你既已答应与她面谈，更亲自查看了那传承碎片，无论你最终是否答应她的具体条件，在你与她接触的那一刻起，在某些时刻关注着落星谷风波、嗅觉灵敏的人眼中，你郑季，就已经与那‘外来的神秘传承’产生了说不清道不明的瓜葛。”

郑季闻言，背脊蓦地升起一股凉意，悚然一惊。

他之前只顾犹豫合作的风险，却未曾从旁观者、尤其是家族监控者的角度想过这一层！夫人所言，真是一针见血！

沈玉继续道，逻辑严密，层层递进：“若此事最终因其他缘由暴露，而那时，你既未能向家族提供关于这‘传承’或‘外来者’的任何有价值信息，又未能站在家族所认为的‘正确’一方，甚至被发现有暗中接触的痕迹……那你很可能会受到比现在更严重十倍、百倍的怀疑和打压。届时，恐怕连如今经营流云坊市的权力，都会被彻底剥夺，甚至可能被圈禁起来，严加审讯。到那时，才是真正的万劫不复。”

郑季额角渗出冷汗，仿佛已经看到了那可怕的场景。自己确实已经半只脚踩进了浑水，沾湿了鞋袜，再想干干净净、悄无声息地抽身退走，恐怕只是痴心妄想。

“反之，”沈玉话锋一转，目光灼灼地看向郑季，“若你选择答应她的条件，固然是与来历不明者合作，风险依旧，但至少你能先得到部分实实在在、梦寐以求的好处。而且，观那人行事，步步为营，投你所好，分寸拿捏得极准；其背后之人，绝非鲁莽无智的狂徒，或许真有几分我们无法想象的能耐与底气。你若能凭借这份传承，在炼器之道上取得前所未有的突破，甚至炼制出一两件令家族高层都不得不侧目、惊叹的法器，那么，你未必不能借此重新引起家族的重视。哪怕这种重视，最初只是将你视为一个有特殊价值的匠师，这也足以让你摆脱目前这种被彻底边缘化、有志难伸的困境。这或许是你此生摆脱桎梏、实现心中抱负的唯一，也是最后的机会了。”

她的话语，平静却充满力量，如同精准的砝码，一枚一枚地放在郑季心中那架摇摆不定的天平上。最终，“唯一机会”这四个字，仿佛最后一根稻草，带着千钧之力，彻底压垮了犹豫，也如同一点星火，投入了他心底那堆被压抑了二百年的、混合着不甘与野心的干燥柴薪之中。

“呼——”郑季长长地吐出一口浊气，眼中挣扎、恐惧、犹豫的神色如潮水般退去，取而代之的是一种破釜沉舟、豁出一切的决绝。

是啊，既然退路已窄，近乎于无，前方纵是龙潭虎穴，又何妨搏上一把？为了那窥见一丝便已心神俱醉的炼器大道，也为了……心中那口憋屈了整整二百年的恶气！

他猛地伸手，紧紧握住了沈玉微凉的柔荑，感受着那熟悉的温度和支撑，声音因激动而有些发颤，却异常坚定：“玉儿，你说得对！分析得透彻！我……我想试一试！为了你，也为了我自己，我要争这一线生机！”

沈玉反手握住他宽大却微微颤抖的手掌，将自己的温暖与力量传递过去，柔声道：“无论你最终作何选择，是进是退，我都会在你身边，支持你。只是，既然选择了这条险路，往后务必要万分小心，如履薄冰。与那等人的交易，须得约定分明，留有后手。”

郑季重重点头，心中暖流涌动，更觉底气足了几分。忽然，他像是想起了什么，犹豫了一下，脸上露出一丝古怪的苦笑，补充道：“玉儿，还有一事……说来也奇，那行脚商人，虽做男装打扮，又多有伪装，但其身形轮廓，尤其是那双眼睛的形状与神韵，还有她卸去部分伪装后，偶尔流露的些许真嗓音……都与你有五六分相似。而且，她也姓沈……当真只是巧合么？”

沈玉闻言，微微一怔，瞅了郑季一眼。尽管她面上依然保持着温婉的微笑，但握着郑季的手，却不自觉地加大了些许力度。

郑季立刻意识到自己失言了——在一个善妒且对自己全心投入的女子面前，提及另一个人与她相似，实属不智。

他赶忙讪讪陪笑，解释道：“是我多心了，天下之大，同姓之人何其多，身形相似者亦不少，想必只是巧合，巧合……”

沈玉这才神色稍霁，轻轻抽回手，重新拿起银剪，语气听不出太多波澜：“嗯，许是巧合。夜深了，此事既已暂定，便不必多思，早些歇息吧。明天你还要去见她呢。”

郑季连忙称是，心中却将那点疑虑悄悄按下，只道是自己心神激荡下的错觉。

一夜无话，但两人皆辗转反侧，各有思量。

\vspace{2em}

次日，约定的时辰将至。

郑季换了一身不起眼的深灰色常服，未带任何随从，如同一个普通的坊市老者，穿行在流云坊市纵横交错的小巷中。最终，他在靠近坊市边缘、一家专供低阶修士歇脚的简陋茶馆后院柴房旁，见到了早已等候在此的沈瑶。

地点是沈瑶通过一种隐蔽的街头标记约定的，柴房旁堆着杂物的阴影恰好形成一个视觉死角，即便白天也少有人注意。沈瑶此刻依旧是那副平凡行商的打扮，但气质却与昨日茶楼中又有所不同，少了些刻意营造的市侩，多了几分沉静与疏离，仿佛与周围油腻的环境格格不入，却又完美地融于阴影之中。

郑季深吸一口气，空气里混杂着劣质茶叶和柴火灰尘的气味。他走到沈瑶面前三步处站定，看着眼前这个依旧神秘莫测、仿佛笼罩着一层迷雾的“行脚商人”，压下心头最后的波澜，沉声开口，每一个字都像是从胸腔里用力挤出：“阁下的条件……我答应了。不过，在交易开始之前，我必须知道，我需要具体做些什么，以及，我究竟能得到什么。” 这是他身为一族长老、一个合作者，必须握有的基本知情权。

沈瑶面具下的嘴角，几不可察地微微勾起一抹弧度。

她知道，经过一夜煎熬与权衡，这枚精心挑选的钉子，终于被成功楔进了郑氏这块看似铁板一块的木头里，尽管只是浅浅的初始。她没有丝毫拖泥带水，开出的条件清晰、直接，边界分明。

“‘沈姑娘’只是个便于称呼的代号，郑长老可继续如此叫我。”沈瑶的声音平静无波，“我们需要的是信息，关于郑氏主脉，尤其是郑天河及其紧密关联的势力，近期针对流云山脉方向，或任何‘外来修士’可能采取行动的预期性情报。不需要细节，只需风向，比如不寻常的修士调动、针对性的资源调配、高层会议的异常议题、乃至郑天河亲信人马的异动。至于我们的具体身份、目的、落脚点，以及其他任何与交易无关的信息，我不会透露，也请郑长老莫要深究，这对你我都好。”

她说着，从袖中取出一枚看似普通的青色玉简。

“而你能得到的，是《焚心器道》的原始版本，分期交付，直到全部。其中所载的上古炼器理论之堂奥，足以让你目前的炼器认知发生天翻地覆的变化，远超你在郑氏，乃至东域绝大多数地方所能接触的一切。”

“信息换传承，公平交易。”沈瑶将玉简轻轻推向郑季，“郑长老只需留意约定范畴内的风声，若有所察，便以我们昨日约定的方式即可。其余诸事，与你无关，切记，知道的越少，对你越安全。”

郑季伸出略微有些颤抖的手，接过了那枚青色玉简。玉简入手温润，带着一丝奇异的、仿佛能安抚心神的凉意。他不再犹豫，当即分出一缕神识，小心地探入玉简之中，随即瞬间头皮发麻、浑身战栗，那是求知者看到真理之门在眼前洞开一隙时，最本能的震撼与狂喜！

系统、连贯的理论阐述与技艺框架！

那是一种与他二百年来所学习、所实践、所信奉的炼器体系截然不同的思路！

他强压下几乎要冲口而出的惊叹与追问，知道这仅是那浩如烟海传承的冰山一角，但仅仅是这一角，其价值已无法用寻常灵石来衡量。

风险固然如影随形，但这回报……对他这个困于瓶颈数十年、壮志难酬、未来一望到底的失意长老而言，诱惑力太大了，大到足以让他甘愿铤而走险。

迅速收回神识，郑季将玉简紧紧攥在手心，仿佛握住了命运的钥匙，也握住了滚烫的炭火。他再次深吸一口气，这次，气息稳了许多，沉声道：“交易……合理。老夫既已应下，自会履行约定，在能力范围内，提供你们所需的风向。但愿……阁下与贵主，亦能信守承诺，后续的传承……”

“自然。”沈瑶微微颔首，不再多言，“第一期内容已交付。望郑长老珍重，谨慎行事。”话音落下，她身形微微一晃，如同阳光下的露珠蒸发，又如同融入水中的墨迹，悄无声息地淡化、消失在那片杂物阴影之中，没有惊起一丝尘埃，仿佛从未出现过。

柴房旁，只剩下郑季一人，独立于略带霉味的空气中。他缓缓将玉简收入怀中最贴身的暗袋，指尖隔着衣料，反复摩挲着那温润的轮廓。心中一半是获得突破希望、触摸到更高境界的炽热火焰，一半是与未知危险共舞、背叛家族规则的冰凉寒意。

郑季回到宅邸时，已近黄昏。他努力调整着自己的表情与气息，想要显得与往常无异，但眉宇间那抹无论如何也挥之不去的深沉凝重，以及眼底深处难以完全掩饰的、因接触到全新知识领域而泛起的隐隐兴奋光芒，终究未能逃过沈玉那双关切而敏锐的眼睛。

结合几日前丈夫透露的关于“神秘行脚商人”与“上古炼器传承”的信息，聪慧如沈玉，立刻意识到事情绝不可能仅仅停留在考虑或简单合作的层面，恐怕在今日，丈夫已经做出了那个危险的决定，并且与对方达成了某种具体的、不便与她言说的契约。

是夜，更深露重，寝室内只余一对夜明珠散发着柔和的光晕。沈玉屏退了所有侍女，亲自落下内室的禁制符箓，确保绝对私密。她走到坐在窗边兀自出神的郑季身旁，轻轻拉住他的手，目光不再是白日里的冷静分析，而是充满了妻子对丈夫最本真的恳切与担忧：“阿季，这里没有旁人，你告诉我，今日……你是不是已经答应了那人什么？具体约定为何？”

郑季从沉思中惊醒，看着妻子眼中毫不掩饰的忧色，心中满是愧疚与温暖。

他苦笑了一下，反手拍了拍沈玉的手背，却避重就轻：“玉儿，我知你担心。但……箭在弦上，不得不发。那人给的定金，啊，不，今日所获，确实是我梦寐以求、足以改变我余生道路的东西。而且，正如你昨日所言，我既已与她有了那般深入的接触，知晓了那等秘密，便再难撇清干系了。如今木已成舟，也只能走一步看一步，步步为营，多加小心便是。”

他并未透露交易的具体内容与方式，这是对沈瑶承诺的遵守，某种程度上也是对沈玉的一种保护。

见丈夫语意坚决，且明显不愿多说细节，沈玉眼帘微垂，不再多劝。

她一向如此：郑季询问她的所有问题，无论是家族琐事还是修炼疑难，她都会竭尽所能，给出详尽的分析与解答，做他最可靠的智囊；但一切涉及到夫妻二人前途命运的最终拍板决断之事，她总是将选择权完全交予丈夫，自己则在他身后全力支持。

然而，此时此刻，沈玉心中的忧虑却如野草般疯长，远比以往任何时刻都要深重。

斗了二百多年，她太了解郑氏主脉，尤其是那位年轻气盛、手段狠辣、野心勃勃的郑天河了。与来历如此不明、偏偏又身怀重宝、已和郑天河结怨的外来者，做这种隐秘的、近乎背叛家族立场的信息交易，一旦有丝毫风声走漏，被郑天河或其爪牙察觉，那后果……她简直不敢细想。

强烈的危机感，以及对丈夫安危的极致关切，驱使着沈玉做出了一个决定。

她不能坐视丈夫独自在刀尖上行走，她必须做点什么，至少，要摸清一点对方的底细，为可能的危机早做准备。出于一种难以言喻的直觉，以及郑季那句“与你相似，也姓沈”带来的微妙异样感，她决定暗中动用自己的力量，进行调查。

接下来的几日，沈玉表面上依旧深居简出，照料灵植，处理一些无关紧要的府内庶务，扮演着一位安静的长老夫人。暗地里，她却悄然启动了自己二百五十年来，在经营流云坊市产业过程中，小心翼翼编织、维护的一张不起眼却触及三教九流的人脉网络。她通过绝对可靠的中间人，以“打听近期坊市异常动向，确保自家产业安稳”为借口，开始撒网，搜集一切关于近期出现在流云山脉附近、与郑氏有过交集、特别是与郑天河产生过冲突的外来者情报。

消息零碎而模糊，需要极高的耐心和洞察力去拼凑。

数日之后，一幅粗略的画像逐渐在沈玉心中成形：一个近期出现在流云山脉、似乎最初与旁支的郑天明有过接触、随后因“落星谷”之事与郑天河正式对上、目前行踪成谜的小团体。核心人物疑似一位姓叶的先生，深不可测，身边跟着几名弟子，其中包括一名冷峻的剑修，一名医术精湛的医者，一名精于幻术与伪装、很可能就是与郑季接触的“行脚商人”的女子，以及最近才出现的一个原本身份低微、却似乎颇有炼器天赋的年轻人。

“叶……牧谦……”沈玉轻轻咀嚼着这个完全陌生的名字，心中的不安如同投入石子的湖面，涟漪不断扩大。这几人如同凭空出现，在郑氏的地盘上搅动风云，得了古传承，竟还敢正面威慑、击退郑天河的护卫……这绝非寻常散修或小势力出身者所能为。

郑季这次，恐怕真是踏上了一条贼船，而船上的人，底细不明，目的成谜。

更让沈玉心湖难以平静，甚至泛起一丝奇异涟漪的，是几条极其隐晦、需要反复印证才能确定的情报反馈。

那“行脚商人”虽常做老朽伪装，但在极少数未曾完全警惕的时刻，惊鸿一瞥的真容或无意间流露的身形习惯，被极少数眼尖之人捕捉到——描述中，那女子容貌颇为清丽，身段窈窕，行动间有种特殊的韵律。甚至，有来自内域、听觉格外敏锐的渠道隐约提及，那女子在极放松或情绪波动时，口音中会带出一丝难以完全掩饰的、与东域截然不同的韵味……那韵味，竟与她沈玉的母语腔调，有几分说不出的相似。

沈……相似的身形容貌……内域南方的口音韵味……

无数被尘封的记忆碎片，在这一刻被悄然触动。沈玉站在花厅窗前，望着夜空中那轮与故乡并无不同的明月，手中无意识地摩挲着一枚一直贴身收藏的温润玉佩，眼神变得悠远而复杂。

一个荒谬绝伦、却又让她心跳莫名加速的念头，如同暗夜中的萤火，在她心底最深处，微弱地亮了一下，旋即又被她强行按灭。

不，不可能……天下之大，无奇不有，怎会如此巧合？

\vspace{2em}

沈玉离开那个被称为“下界”的内域，已经整整二百五十年了。

当年灵界使者降临，惊叹于她的修炼天赋与灵秀根骨，力邀她前往更高层次的灵界修行。怀着对大道更高风景的向往，也带着少年人特有的雄心与对家族安排的些许叛逆，她毅然辞别了父母双亲，踏上了跨越界域的旅程。

然而，灵界中州那看似繁华鼎盛的修真文明，骨子里却浸透着对下界来人根深蒂固的轻视与排斥。天赋并不能完全抹平出身的差距，冷眼、排挤、隐晦的掠夺……最终，她负气离去，在广阔的灵界四处漂泊、辗转。直至来到相对偏远、势力格局也更为松散的东域，遇到了当时尚且意气风发、更不在意她所谓下界出身、只欣赏她本人的郑季，两颗心才逐渐靠近，最终相知相守，这时她才终于停止了漂泊，安顿下来。

这二百五十年间，她与故乡——按灵界的称谓，应该是“下界”——几乎断绝了一切联系。初时是赌气，觉得混不出名堂无颜回去；后来是漫长的漂泊与挣扎，无暇他顾；等到安定下来，岁月已蹉跎太多，故乡人事皆非的恐惧，又让她怯于打探。

她只知道，在自己离开很久之后，父亲沈天穹与母亲又育有一女，但那对她而言，仅仅是一个遥远而模糊的血缘符号，甚至连妹妹的具体名讳都未曾知晓。

一个天师府的千金小姐，即便后来修炼有成，又怎么会如此巧合地、偏偏也来到这东域，卷入与郑氏如此凶险的纷争，甚至伪装成卑微的行脚商人，来与自己那失意困顿的丈夫做这等隐秘交易？

更何况界门司不是说今年也“没有好苗子”吗？

这听起来像是话本里才会出现的离奇桥段。

然而，理智的否定，却压不住直觉的疯狂滋长。那“行脚商人”行事风格中偶露的果决与灵动机变，处理问题时那种隐隐的、不同于寻常散修的格局与章法，都让她恍惚间，仿佛看到了记忆中天师府精英子弟的某些影子。

疑虑的种子一旦落下，便在心土中悄然生根发芽，再也难以拔除，反而在每一个相似细节的浇灌下，越发茁壮。

她必须知道答案。

\vspace{2em}

数日后，一次经过精心设计和多重掩饰的偶然机会，沈玉与那行脚商人约定在坊市外数十里、一处隐藏在人迹罕至山坳竹林深处的私人小筑见面，声称“有关于传承真伪与后续合作风险的重要事宜，需当面厘清，关乎重大，请务必独自前来”。

这座小筑是沈玉多年前发现并修缮的，除了郑季再无别人知道，是两人偶尔想独处时的避世之所。环境清幽绝伦，竹涛声声，灵泉潺潺，安全性与私密性都极高。

沈瑶收到那条通过特定渠道传来的口信时，心中微微一惊。她以为是郑季那边又出了什么新的变故，或是交易出了纰漏，需要紧急沟通。

虽然师尊叶牧谦曾嘱咐诸事需谨慎，但此事涉及埋下的暗桩是否稳固，她不敢怠慢。仔细查验了接头方式与暗记无误后，她一如往常般改换装束，凭借高超的身法和幻术，如同融入夜色的青烟，悄然避过可能存在的眼线，向着约定的山坳竹林潜行而去。

月华如练，倾泻在连绵的竹海之上，泛起一片朦胧的银绿色光晕。夜风穿过竹叶的缝隙，发出沙沙的轻响，如同无数细碎的私语。小筑就坐落在竹林环绕的一小片空地上，以灵竹和青石搭建，样式简朴古雅，与周围环境浑然一体，檐下悬挂着几盏未点燃的旧式风灯。

沈瑶灵觉铺开，确认方圆数里内除了些无害的小型灵兽，并无其他修士气息，小筑内也只有一道平稳而隐晦的法相境气息——有很浓重的郑季气味。

她稍稍安心，轻如鸿毛般落在小筑门前，略一犹豫，还是以“沈三”那副平凡行商的伪装模样，轻轻推开了虚掩的竹扉。

小筑内没有点灯，唯有清澈的月光透过雕花木窗，在地上投下斑驳静谧的光影。一股极淡的、似曾相识的冷梅幽香萦绕在空气中。一道身着月白长裙的窈窕身影，正背对着门口，静静立在窗边，仰头望着挂在墙上的一幅画。

沈瑶的目光立刻被那幅画吸引了过去。画纸已有些泛黄，但保存得极为精心。上面描绘的是一座巍峨古城的一角，飞檐斗拱，气象万千，城中有一高塔直插云霄，塔尖似乎隐没在流转的云层之中——那是内域天枢城的标志性建筑！

而画的笔法、用色、乃至边角处一个属于天师府的戳印，都让沈瑶的心猛地一跳。

那是姐姐当年从天师府带出来的、为数不多的故物之一？在临去灵界之前，她才听父亲沈天穹完整提起过姐姐的过往，听说她离家时，只带走了几件随身物品和一幅自己画的天枢城图，以慰思乡之情。

就在沈瑶心神剧震之际，那背对着她的女子，仿佛听到了她瞬间紊乱的呼吸与心跳，缓缓地、极其缓慢地转过了身。

月光恰好在这一刻，透过窗格，清晰地照亮了彼此的面容。

“嗡——”

时光的长河仿佛在这一刹那被无形的手截断、凝固。空气不再流动，竹涛声、泉流声、甚至夜虫的微鸣，都从感知中彻底消失。世间万物，只剩下眼前这咫尺相望的两人。

尽管沈瑶的脸上依旧覆着“沈三”那层平凡的伪装，但那伪装之下，那双因为震惊而微微睁大的眼睛——那眼型，那睫毛翘起的弧度，那瞳孔深处瞬间掠过的难以置信、急切探寻、以及一种血脉深处被唤起的、难以言喻的激动微光——都与沈玉每日在镜中看到的自己，产生了惊人的、几乎令她眩晕的重叠。

沈瑶同样死死地盯着眼前的女子。雍容的气度中沉淀着岁月的历练与漂泊的痕迹，温婉的眉宇间却藏着一股不轻易显露的柔韧与刚毅。容貌确与自己有五六分相似，尤其是挺秀的鼻梁和下巴的轮廓。但更重要的是，在如此近的距离内，一种源自同源血脉的、微妙的灵魂共鸣与亲近感，如同沉睡的琴弦被轻轻拨动，发出只有彼此能感知的颤音。而对方那双同样瞬间掀起惊涛骇浪、混杂着震惊、狂喜、困惑与不敢置信的眼眸，更是让沈瑶几乎可以肯定。

是她……真的是她吗？那个只在父亲满是思念与愧疚的讲述中存在的、素未谋面的……姐姐？

窒息的沉默在月光中流淌，仿佛持续了永恒，又仿佛只是一瞬。

“你……”沈玉张了张嘴，声音干涩得厉害，仿佛许久未曾说话，又仿佛有千言万语堵在喉头，最终只挤出一个颤抖的音节。

沈瑶没有立刻回答。她深深地、深深地吸了一口气，仿佛要将这混杂着冷梅香与竹叶清气的空气，连同眼前之人带来的巨大冲击一同纳入肺腑。然后，她抬起手，指尖闪烁着极细微的灵力光华，如同剥开一层轻柔的纱，缓缓拂过自己的脸庞。

粗糙的男性肌肤纹理褪去，暗沉的肤色变得白皙，平凡的五官如同被一只无形的手精心雕琢、复位，显露出其下清丽灵动的本来容颜——眉毛细长如远山含黛，眼眸明亮若秋水寒星，鼻梁挺直，唇色如樱。虽因长期历练带着风霜之色，不如沈玉那般养尊处优的温润，但那眉眼神情，那骨相轮廓，与沈玉如同一个模子刻出的、经过不同岁月雕琢后的版本。

卸去大部分伪装后，沈瑶依旧没有说话。她只是用那双与沈玉愈发相似、此刻却噙着隐隐水光的眼睛，静静地看着对方，然后，极其缓慢而郑重地，从自己怀中贴身的内袋里，取出一物。

那是一枚小小的玉佩。质地是最上等的羊脂灵玉，触手温润生津，在月光下流淌着柔和静谧的光泽。玉佩正面，以古朴苍劲的笔法，刻着一个古老的“沈”字。

翻到背面，则是天师府独有的、象征着云海巡天的复杂云纹标记，而在云纹的中心，以纤巧的笔触，刻着一个清晰的“瑶”字。

这玉佩，自她记事起便佩戴在身上，父亲沈天穹曾说，这是她出生那日，他亲手为她戴上，以灵气温养，寄托平安与期盼。

当那枚玉佩在月光下完全显现，当那熟悉的“沈”字与天师府云纹映入眼帘，尤其是那个小小的“瑶”字清晰无比时，沈玉脑中最后一丝残存的、支撑着“巧合”幻想的弦，彻底崩断了。

她眼睛瞬间睁到极大，瞳孔收缩，呼吸骤然停止，随即，滚烫的泪水如同决堤的洪流，毫无征兆地、汹涌地从眼眶中疯狂溢出，瞬间模糊了视线，滑过白皙的脸颊，在下颌处汇聚、滴落，在月白衣襟上晕开深色的湿痕。她颤抖着，同样从自己袖中一个几乎从不离身的暗袋里，取出一枚玉佩。

款式、质地、大小、正面那古朴的“沈”字、背面天师府的云纹标记……几乎一模一样。唯一的不同，只是云纹中心刻着的，是一个笔迹略显青涩却透着坚定的“玉”字。

两枚玉佩，在清冷的月光下，静静辉映，像是分隔了二百五十年的血脉印记，终于在此刻，跨越了无尽的空间与漫长的光阴，重新聚首。

“你……你真是……我妹妹？”沈玉的声音破碎不堪，带着剧烈的哽咽，每一个字都像是用尽了全身的力气。泪水模糊中，她死死盯着妹妹那张与自己年轻时有七分相似的脸，巨大的狂喜与尖锐的心痛交织着撕裂她的胸膛，“父亲他……父亲他可还好？你……你怎么会在这里？还卷进这些事情里来？”

“姐姐——！”

所有的伪装、所有的谨慎、所有的算计，在这一声饱含了二百五十年思念、担忧与委屈的呼唤中，彻底冰消瓦解。沈瑶的泪水也终于夺眶而出，如同断了线的珍珠。她一个箭步冲上前，张开双臂，用尽全身力气，紧紧地、紧紧地抱住了沈玉，仿佛害怕这只是一场醒来即碎的幻梦。

沈玉也反手死死抱住妹妹纤细却充满力量的身躯，将脸埋在她带着夜露凉意的肩头，压抑了二百五十年的思乡之情、漂泊之苦、对父母的愧疚、对自身命运的彷徨，以及此刻失而复得的巨大狂喜，化作无声的痛哭，肩膀剧烈地颤抖着。

姐妹二人相拥而泣。没有言语，只有滚烫的泪水浸湿彼此的衣衫，只有压抑不住的抽泣声在静谧的竹舍内回荡。跨越了界域屏障，跨越了漫漫光阴，在这异乡他域、危机四伏的险境之中，两条离散已久的血脉支流，竟以这样一种谁也无法预料的方式，轰然交汇。

积压的情感如同被封印了太久的火山，在这一刻轰然爆发，炽热的岩浆奔涌流淌，灼烧着理智，也温暖了孤寂太久的心。

\vspace{2em}

不知过了多久，窗外的月影已经悄悄偏移，竹涛声重新涌入耳中。两人的哭声渐渐低了下去，变为断续的抽噎，但拥抱的力道却未曾松懈半分。

她们相携着，在竹榻边坐下，手却依然紧紧握在一起，仿佛一松开，对方就会消失。借着清澈的月光，她们贪婪地打量着彼此的面容，从眉眼到唇角，从神情到风霜留下的细微痕迹，仿佛要将这缺失的二百五十年时光，一点一点地看回来，补回来。

“姐姐，先擦擦。”沈瑶从袖中取出干净的丝帕，有些笨拙却轻柔地替沈玉拭去脸上的泪痕，自己的眼圈依旧红红的。

沈玉接过帕子，也帮妹妹擦了擦脸，含着泪笑了，那笑容里满是失而复得的珍宝般的珍重。“瑶儿……让姐姐好好看看你。爹爹……爹爹他身体可还康健？天师府……如今怎么样了？”她迫不及待地想要知道故乡的一切，那是她午夜梦回却不敢深探的角落。

沈瑶握着姐姐的手，组织了一下语言，将她所知道的情况，娓娓道来。她讲述了父亲沈天穹这些年的状况，虽偶有旧伤隐痛，但总体康健；也提及了天师府这些年的变迁，一些人事浮沉，但大体安稳。她接着说起自己年少时不甘困于府内，执意离家外出历练，途中遭遇险境，幸得师尊叶牧谦出手相救并赏识，从此拜入门下，研习幻术、杂学与处世之道。她隐去了一些内域纷争的残酷细节，但简要提及了内域近些年并不太平，各方势力暗流涌动，但师尊全都给摆平了；后来灵界来的恶客也被师尊收拾了，她最终跟随师尊离开内域，来到灵界。

沈玉含泪听着，时而因父亲的安好而欣慰点头，时而因妹妹的惊险遭遇而揪心握紧她的手，时而因天师府的变迁而感慨唏嘘，时而又因那从未谋面的“叶牧谦”师尊对妹妹的恩情生出感激。

待沈瑶说完，她才轻轻拭去眼角的泪，开始述说自己的经历。

她讲述了当年离家后，初至灵界中州时的憧憬与随后遭遇的冷遇排挤，那份理想与现实碰撞的失落与愤怒。说起负气离开中州后，在广阔灵界独自漂泊的艰辛与孤独，遇到过危险，也得到过零星善意。最后，是如何来到相对平和的东域，在一次坊市纠纷中，偶然结识了当时尚在家族中颇受重视、为人正直热忱、更对她“下界”身份毫不在意的郑季。两人从相识到相知，彼此扶持，历经风雨，最终结为道侣。她也坦承了郑季在争夺家主之位失败后，两人在郑氏家族中所处的尴尬地位，看似守着流云坊市这份产业安稳度日，实则一直备受主脉，尤其是现任家主一系的隐形打压与排挤，抱负难展，郁郁不得志。

“阿季他……这次之所以会答应与你们交易，也是不得已。”沈玉说到此处，眼中流露出深切的疼惜与无奈，再次为丈夫分辩，“他在炼器之道上天赋不俗，也有大志向，可家族……从未给过他应有的支持与机会。那《焚心器道》的传承碎片，对他而言，无异于无尽黑暗中突然出现的一缕真实光芒，是他抓住突破瓶颈、实现自身价值的唯一可能。他现在肯定不是贪图富贵权势了，只是不甘心一生所学就此埋没，不甘心将来被人彻底遗忘，想在这条他挚爱的大道上，真真切切地留下一点属于自己的痕迹罢了。”

沈瑶认真地听着，用力回握姐姐的手，表示理解：“姐姐，我明白的。此事我已然回禀师尊。我们与郑长老的交易，本就是各取所需，互惠互利。郑长老提供了我们所需的信息风向，我们给予他探寻炼器大道的阶梯，这很公平。只是……”她神色转为凝重，“如今那郑天河在外散布谣言，步步紧逼，形势颇为险恶，我们师徒的行踪与意图，越少人知道越好。姐姐，你和姐夫身处郑氏内部，更需万分小心，谨言慎行，切莫让任何人察觉端倪。”

沈玉重重点头，眼中的泪光已被一种坚定如铁的光芒取代。血脉亲情的确认，激起了她骨子里的护短与决心。

“既然知道是你，是我的亲妹妹在此，那我更无置身事外的道理。”

姐妹二人低声且快速地交流起来。

至少现在，她们不必再担心这条情报路径会出问题了。

\vspace{2em}

流云山脉深处，叶牧谦等人的临时洞府。

石室中央，仅有一小块取自落星谷的星纹铁原矿悬浮空中，散发着微弱的、星屑般的光点。叶牧谦闭目盘坐其前，周身并无丝毫灵力外泄的波澜，甚至呼吸都微不可闻，仿佛整个人已化为一块温润的玉石，与这山腹的沉寂融为一体。然而，若有心念境界高深者在此，必能感知到，一股宁静、圆融却又浩瀚深邃的“意”，正以他为中心，如同无形的水波，柔和而持续地弥漫开来，充塞着整个石室。

他将全部心神，沉入那原矿的最深处。于是，在他那映照万物的心念感知中，眼前这块看似死物的金属，内部的一切奥秘纤毫毕现：天然生成的、宛如星辰轨迹的灵性纹路，错综复杂却又暗含某种韵律的物质结构网络，被视为瑕疵的细微杂质，其分布、性质清晰无比。

叶牧谦所走的，是不同于风灵希的另一条路。

他的通明心念，如同最温柔也最坚韧的春水，缓缓溢出，轻柔地包裹住那块星纹铁。他极其细微地调整着自身的心念波动，去悄然触碰那些天然的灵性节点，去安抚金属本身固有的、略显躁动的物性，去倾听材料在微观层面最细微的诉求。然后，他以这浸润万物的心念为引，引导着星纹铁内部的结构，朝着更坚韧、更纯净、更贴合其内在潜藏脉络的方向，极其缓慢地进行着自我调整与优化。

这个过程，远比用真火猛烈灼烧、用巨锤千锤百炼要缓慢十倍以上，但其效果却截然不同：对材料本源灵性的伤害降至最低，那些脆弱的天然星纹，在温和的共鸣滋养下愈发清晰、灵动。

数个时辰的光阴，在绝对的静默中流淌而过。洞府外日升月落，洞府内唯有心念的微光与星纹铁逐渐变化的光泽交相辉映。

终于，叶牧谦缓缓睁开了双眸。眼中神光湛然，清澈见底，不见丝毫疲惫，反而有一种水到渠成的圆满之意。他轻轻呼出一口气，那气息悠长平和，带着山洞中泥土与灵泉的清新。

眼前，那块星纹铁胚料，体积已然缩小了约三分之一，形态却更加规整。表面原本略显晦暗的星纹，此刻变得无比清晰、优美，如同将一片微缩的璀璨星河烙印在了沉凝的玄铁之上。质感内敛至极，光华含而不露，只在心念微微触动时，才有一抹深邃的星辉流转而过。虽然它还远未成器，但无论是材质的纯净度、灵性的活跃度，都已远非先前可比。

即使是这块没成材的半成品本身在流云坊市出卖，应该也是最上等的一批货了。

第一次试验还不熟练，但叶牧谦能够确定：这方法，可行！

“师尊，成功了？”一直屏息凝神、在旁护法兼观摩的柳成连忙上前，声音因激动而有些发颤，眼中充满了难以抑制的惊叹与渴望。

他虽然转修问道途初成，仅仅踏入养气、凝聚了最初级的气海，但得益于问道途对心神感知的天然侧重，他也能模模糊糊地感受到师尊刚才那种炼器方式的玄妙。那与他从《焚心器道》玉简中感受到的、充满惨烈决绝意味的“焚心”之法截然不同，少了那份令人心悸的自我毁灭感，多了几分顺应自然、引导共鸣的深邃与和谐，仿佛大道本应如此。

叶牧谦看着弟子眼中的光，微微颔首，随即又几不可察地摇了摇头：“算是初步踏出了一条路。”

他指尖轻点，那块蜕变后的星纹铁轻盈落入掌心，触感温润，仿佛有了生命。“此法以心念为桥梁，以共鸣为根本，摒弃了自燃心念的凶险法门。对心念的消耗更为温和、持久，痛苦大减，更契合我问道途循序渐进、修心炼性的根本宗旨。”

他顿了顿，目光投向柳成，既是解说，也是传授：“此法，可称之为《明心器道》：明心见性，以心映物。然此法门槛亦不低，至少需心念修为达到见独之境，心念初步凝实成形，方能稳定外放，进行如此精微的引导与共鸣。你如今初入养气，根基未稳，距此尚远。”

柳成眼中闪过一丝明显的失望，但很快被更坚定的斗志取代。有了明确且更安全的前进方向，这本身就是最大的希望。

“弟子明白！”他用力点头，“定当日夜勤修不辍，先夯实问道途根基，温养心神，早日达到修习的门槛！”

“甚好。”叶牧谦将星纹铁收起，“万丈高楼平地起，根基最为紧要。”

\vspace{2em}

与此同时，郑氏主脉。

郑天河精心策划的舆论风暴，经过近一个月的持续发酵与暗中推波助澜，已然在特定的修士圈子里达到了他所期望的沸点。

“落星谷上古炼器传承被神秘外来者独占”的消息，如同野火燎原，不仅引来了无数散修和中小势力或明或暗的觊觎目光，也果然成功引起了天工阁这一庞然大物的持续关注。郑天河通过自己遍布坊市的眼线汇报，确认天工阁已有执事级别的人物在流云坊市暗中活动，四处打探。

书房内，郑天河负手立于窗前，目光穿越庭院，遥望着远方云雾深处若隐若现的流云山脉轮廓，嘴角噙着一丝冰冷而满意的弧度。

“火候差不多了。”他低声自语。

造势已成，鱼儿已经循着饵料的血腥味聚集而来，但他依然保持着猎人般的极度谨慎。

亲自下场，与叶牧谦那种深浅不明、手段诡异、身边还聚拢了几个各有所长的弟子的人物硬碰硬，那非智者所为。

借刀杀人，驱虎吞狼，坐收渔利，才是上上之策。

他需要一把足够锋利、能伤敌见血，但又不能太过锋利以致反噬自身，且最好易于操控的“刀”。这把刀，需要对古传承，或者对方可能掌握的、某种能“直指心念神魂”的奇异能力有迫切的渴求；其实力需足以构成威胁，但又不能强到脱离掌控，甚至最好与郑氏本身有着千丝万缕的联系，以便在事后无论是安抚、奖赏还是切割、问罪，都能占据主动。

很快，密探呈上的一份名单中，一个名字引起了他的注意——寒玉门长老，冰魄，人称冰真人。

寒玉门，一个依附于郑氏生存的中型宗门，以修炼冰寒属性功法著称，门中高手多在郑氏领地内担任客卿或护卫，关系盘根错节。

冰真人乃门主，修为在法相境通灵期，在东域也算小有名气，实力不俗。但此人有一致命弱点，或者说，一个足以驱使他疯狂的执念——早年因修炼急功近利，急于突破境界，又在一次探索遗迹时遭遇强敌，伤了神魂根本。

多年来，这神魂旧伤时好时坏，不仅将他死死卡在法相境和道种境之间，更让他常年承受着神魂刺痛、偶发性幻听幻视之苦，性情也因此越发阴郁、偏激、易怒。为治此伤，他耗尽积蓄，遍访名医，试过无数偏方秘药，效果皆微乎其微，几乎已近绝望。

“就是他了。”郑天河眼中精光一闪。一个有实力、有动机、有弱点、且其宗门依附于郑氏的完美刀子。

他立刻命心腹暗中搜集、整理了一份关于叶牧谦师徒的、看似详尽实则经过精心裁剪的情报。情报着重突出了两点：其一，为首的叶先生极可能擅长某种直指心念、影响甚至操控他人神魂的奇异领域；其二，其门下姓陈的女弟子医术超凡，曾当场为旁支的郑天明诊治复杂的内腑伤势，手法精妙绝伦，堪称手到病除。

然后，郑天河以一种看似无心实则精妙的方式，让这份情报“意外”地流落到了与寒玉门素有往来、且一直有意巴结冰真人的一位郑氏旁系子弟手中，并让人隐晦暗示：此等关乎神魂伤势与高深医术的消息，或可作为一份厚礼，献予正为此烦恼不堪的冰真人，以结善缘。

钓饵精准地抛向了饥渴的鱼。

\vspace{2em}

果然，冰真人在得到这份“意外之喜”后，几乎是从座椅上弹了起来，枯瘦的手掌因激动而微微颤抖，那双因常年痛苦而布满血丝、显得阴鸷的眼睛里，爆发出骇人的精光！

擅长心念之道？直指神魂？那岂非正是修复自己这陈年痼疾的最大希望所在？！

门下还有医道圣手？或许能炼制出对症的灵丹妙药！

这简直是天赐的机缘，是照亮他黑暗绝望道途的曙光！

至于对方是否愿意帮忙？在冰真人那被痛苦和偏执浸透的思维里，这根本不是问题。几个无根无底、不知从哪个犄角旮旯冒出来的外来修士，只要擒下，以寒玉门和郑氏的势力为后盾，自有无数手段让他们乖乖吐出传承、献出医术！那所谓的上古炼器传承，届时自然也成了他囊中之物，或许其中就有辅助修炼、温养神魂的秘宝！

贪婪的火焰瞬间吞噬了本就所剩无几的理智与顾忌。他迫不及待，立刻通过那位郑氏旁系子弟作为中间人，急切地向郑天河表达了“愿为公子分忧，出手擒拿那几名不识抬举、胆敢冒犯郑氏威严的外来修士，并顺便‘请教’一番他们那奇特的医术与传承”的意愿。

郑天河在书房“欣然”接见了这位中间人，听罢转述，脸上露出恰到好处的“欣慰”与“凝重”。

他沉吟片刻，对中间人（实则是说给冰真人听）忧心忡忡地道：“那几人实力深浅，尤其那姓叶的，手段确实诡异莫测，李辉他们真符境的修为，在其面前竟也无丝毫反抗之力，估计至少是个法相。冰长老若决定出手，务必要计划周详，调集足够人手，以雷霆之势，务求一击必中，绝不能给他们喘息或走脱的机会，否则……后患无穷啊。”

他刻意反复强调叶牧谦的诡异与危险，以及陈千羽医术的超凡价值，既深深刺激了冰真人心中的贪欲和对他自身伤势治愈的渴望，也成功地在冰真人心头压上了一块名为忌惮的巨石，使其必然会下定决心，狮子搏兔亦用全力。

冰真人得到郑天河“关怀备至”的提醒后，非但没有退缩，反而更觉郑天河思虑周全，心中杀意与决心愈发炽烈。他眼中寒光四溢，对中间人冷然道：“请回禀天河公子，老夫省得了。要么，生擒活捉，撬开他们的嘴，拷问出所有传承与医术秘密；若事有蹊跷，难以生擒……便就地格杀，以绝后患！绝不可放虎归山！至于他们身上的古传承，想必也会随身携带，届时自然见分晓。”

双方通过中间人秘密敲定了行动的诸多细节：大致锁定了叶牧谦等人最可能藏身的流云山脉中部区域，约定了冰真人出动的人手规模（包括他自己，三名法相境得力长老，以及十余名真符境门下好手），甚至初步商议了事成之后，关于“传承共享”与“医术服务”的模糊分配方案。

然而，百密终有一疏。

或许是出于上位者记录重要事务以备查询的习惯，或许是郑天河潜意识里想留下一点制约或事后追究的凭据，他与冰真人之间往来沟通的玉简讯息，竟然没有全部销毁；其中一份记载了双方初步意向、尤其是冰真人那句“若事不可为，便就地格杀”的明确表态的玉简副本，被郑天河习惯性地归档，加密存放于他书房某个隐秘的暗格之中。

\vspace{2em}

数日后，正午时分。

流云山脉深处，阳光努力穿透茂密的林冠，投下斑驳晃动的光斑。鸟鸣啁啾，溪流潺潺，一切似乎与往日并无不同。

突然，一股凛冽刺骨的寒意，毫无征兆地自山林某一处爆发，并以惊人的速度席卷开来！仿佛无形的寒潮过境，所过之处，生机勃勃的翠绿草木以肉眼可见的速度凝结上一层厚厚的白霜，潺潺的溪流表面瞬间覆盖上不透明的坚冰，活跃的鸟兽虫鸣戛然而止，要么僵毙坠落，要么惊慌逃窜。温暖的阳光仿佛失去了温度，空气中的水汽凝结成细小的冰晶，在光线中闪烁，却只带来更深的寒意。

这片区域的中心，正是叶牧谦师徒暂居的那处隐蔽山洞。

“里面的人听着！”

一个苍老、嘶哑、仿佛摩擦着冰碴的声音，如同腊月寒风刮过空旷的雪原，带着浸透骨髓的冰冷与毫不掩饰的恶意，在山洞外的山谷中隆隆炸响，反复回荡，震得洞顶簌簌落下些许尘灰。

“老夫寒玉门冰魄！听闻尔等身怀上古传承，门下更有医道圣手。特来请诸位前往寒玉门做客，交流道法医术！若识相，便主动出来，交出传承，束手就擒，尚可免去一番皮肉之苦，留得性命！否则……休怪老夫玄冰真罡之下，不留全尸！”

话音未落，一道道周身散发着冰冷气息的身影，如同鬼魅般从挂满冰霜的林木后、岩壁阴影中浮现。他们动作迅捷，训练有素，迅速结成一种攻守兼备的阵势，隐隐蕴含着冰寒灵力的联动，将山洞出口方向以及周围可能逃遁的路径，封锁得水泄不通。

为首的，正是须发皆白、面皮如同覆盖着一层薄冰、看不出太多表情，唯有一双深陷眼窝中的眸子闪烁着阴鸷与贪婪寒光的冰真人。他负手而立，周身弥漫着肉眼可见的淡淡白色寒雾，法相境通灵期的威压毫不掩饰地释放开来，压迫着周遭一切。凛冽的杀气与彻骨的寒气混合交织，令这片原本清幽的山林，瞬间化作了致命的冰封绝地。

山洞内，篝火的光芒因洞口涌入的寒气而摇曳不定，将洞壁上的人影拉扯得晃动不安。

该来的，终究还是来了。

\vspace{2em}

流云山脉另一座视野开阔、可俯瞰大片山林的山巅之上。

三道身影悄然伫立，气息收敛得近乎完美。

为首者是一位身着暗金色锦缎长袍的老者，面容清癯，颧骨微凸，一缕银白长须修剪得一丝不苟，垂于胸前。他眼神温润，初看如寻常慈祥长者，但眼底深处却有着历经世事沉淀下的深邃与精明。此刻，他正举着一枚黄铜材质、刻满细密符文的单筒法器，遥遥观察着远方山洞方向的动静。

他便是天工阁派驻东域、常驻流云坊市一带的阁老，吴毓。

吴毓身后半步，垂手侍立着两名随从，一男一女。二人皆身着天工阁制式的玄色劲装，领口袖口以银线绣着代表工具与火焰的徽记，简洁干练。两人气息凝实沉厚，目光开阖间锐气隐现，显然是阁中悉心培养、经验丰富的好手。

手持一方罗盘状法器、正盯着远方冲突现场明灭灵光的男随从，低声汇报道：“阁老，寒玉门的人已经将那里围死了。冰魄那老家伙亲自来了，神识感应中，其灵力波动剧烈，确为法相境通灵期无疑。随行人员中，可辨明有三位法相境和十余名真符境修士的气息，玄脉及以下修士更多，已封死了那山洞周遭所有出口与气机。”

一旁的女随从立刻开口补充，声音清脆而条理分明，显然事先做足了功课：“综合我们近一个月在坊市及周边明查暗访所得，那山洞里藏身的，九成九便是近来传闻中得了落星谷上古炼器传承的那一伙外来修士。

“为首者姓叶，具体名讳、出身背景皆不详，深居简出，极少露面。其身边跟着两男两女共四位弟子。其中那身着青色衣裙、气质温婉的女子，疑似精通医道，曾为郑氏旁支的郑天明诊治过棘手内伤，手法迥异于常。另一白衣负剑男子，冷峻寡言，剑意凝练，观其步履气息，剑术造诣定然不凡。还有一女，行踪最为飘忽，极擅伪装幻化之术，坊间关于行脚商人的传言多与其有关。最新加入的那年轻人，则是原本在坊市外围摆摊、兜售低阶法器的炼器学徒柳成，天赋据说尚可，但修为低微。”

吴毓轻轻嗯了一声，手中千里镜镜片上流转的细微符文调整了焦点，将远处的景象拉得更近、更清晰。他看到了冰真人那张因常年痛苦与阴郁而显得扭曲僵硬、此刻却因贪婪与急迫而微微泛红的面容，也看到了寒玉门修士们井然有序、隐隐联动、散发着协同寒气的包围圈，以及那被冰霜彻底覆盖、死寂一片的山林。

“冰魄这老家伙，”吴毓嘴角泛起一丝早已洞悉的淡然笑意，把那枚极其昂贵的“千里镜”仔细地收纳进袋中，捋了捋银须，“那神魂旧伤折磨了他百余年，近乎成了心魔。一听到有擅长心念之道、可能直指神魂病灶，门下还有医道圣手的外来者出现，果然像是嗅到血腥味的豺狼，立刻就坐不住了。”

他顿了顿，语气中带着对另一方的评点，“郑天河那小子，年纪不大，这份借刀杀人、驱虎吞狼的心思，倒是用得娴熟。这把刀，选得也算好，够急，够贪，也够狠。”

“阁老，我们何时出手？”男随从低声问道，眼中闪过一丝热切，“难道真坐视寒玉门将那可能的上古炼器传承夺了去？冰魄老儿可不懂什么炼器精要，落在他手里怕是糟蹋了。”

“不急。”吴毓气定神闲，目光重新投向远方战场，仿佛在看一场与己无关的戏剧，“螳螂捕蝉，黄雀在后。先让他们斗上一斗，我们静观其变。我天工阁要的，是那传承本身，尤其是其中可能蕴含的、迥异于当今体系的上古炼器理念与秘法。至于那几个外来修士的性命，并非首要；与郑氏这地头蛇或其麾下附庸寒玉门正面冲突，更是下下之策，智者不为。

“且看那姓叶的如何应对。若他们只是纸老虎，不堪一击，迅速被寒玉门以力碾压、生擒活捉，那再好不过。届时我们再现身，以天工阁的名义与冰魄商议，或购买、或合作研习那份传承。谅他区区一个寒玉门长老，不敢不给我天工阁几分薄面，传承终究会落入我们手中，还省了力气。”

他话锋一转，“若他们真有几分出人意料的能耐，能与寒玉门周旋一阵，甚至拼个两败俱伤……哈，那便是更妙的局面。届时我们以调解纷争为由出手介入，既结了善缘，又能顺理成章地接触、评估那传承的价值，进退自如，岂不美哉？”

天工阁能富甲灵界、势力遍布，这种深入骨髓的、审时度势、谋定后动的智慧与耐心必不可少。为了一份可能极大提升阁内底蕴的上古炼器传承，值得他这位阁老亲自隐匿于此，如同最高明的棋手，等待最恰当的落子时机。

\vspace{2em}

山洞前，死寂般的对峙并未持续太久，气氛迅速升级至爆裂的边缘。

冰真人那饱含灵力、如同刮骨钢刀般的喊话在山谷间反复回荡、碰撞，激起阵阵冰屑。然而，回应他的，只有洞口呼啸而过的、愈发凛冽的寒风，以及山洞内那片深不见底的、令人心悸的沉默。

片刻的等待，如同无声的羞辱，彻底点燃了冰真人本就因神魂旧伤而异常暴躁的怒火。他兴师动众，志在必得，岂能容忍几个藏头露尾的鼠辈如此蔑视？！

“冥顽不灵！真当老夫不敢将这破山头彻底夷为平地吗？！”

冰真人怒喝一声，周身原本就已澎湃的寒气如同火山般轰然爆发！只见他身形依旧站在原地未动，但其身后的虚空却剧烈地扭曲、震荡起来！无尽的冰蓝色灵光从四面八方疯狂汇聚，空气被冻结、挤压，发出令人牙酸的“嘎吱”声。

眨眼间，一尊庞然大物拔地而起，耸立于冰真人身后的半空之中！

那是一座高达百丈、通体由绝对纯净、晶莹剔透的万载玄冰凝聚而成的巨人法相！法相面目模糊，唯有眼部位置燃烧着两团幽蓝色的冰冷魂火，俯瞰众生。它仅仅是矗立在那里，一股令灵魂战栗、血液凝固的恐怖威压便如同实质的海啸般席卷四方！

法相显现的刹那，方圆数里内的天地仿佛被瞬间拖入了极北冰原的深渊——温度骤降至真正的滴水成冰，甚至更低！原本只是挂霜的草木“咔嚓”一声彻底化为栩栩如生的冰雕，溪流冻结的坚冰厚度瞬间增加数尺，连那些裸露的岩石表面，也在“滋滋”声中覆盖上足有半尺厚的、闪烁着金属寒光的坚冰。阳光失去了所有温度，苍白无力地照在这片突然降临的冰封国度之上。

“搜！给老夫把这方圆十里，掘地三尺，一寸寸地翻过来，也要把这群藏头露尾的鼠辈揪出来！”冰真人须发皆张，厉声咆哮，声音通过百丈法相放大，如同冰川崩塌，震耳欲聋。

那百丈玄冰法相早已通灵，理解主人的意图之后，缓缓抬起了擎天柱般的巨大冰掌。五指张开，遮天蔽日，带着冻结虚空、碾碎山岳的无匹威势，作势便要朝着叶牧谦师徒藏身的山洞所在的那面山壁，狠狠拍下！

这是要以绝对的力量，粗暴地撕开一切伪装：要么逼对方现身，要么就连山带人，一并轰成齑粉！

\vspace{2em}

山洞内，篝火的光芒在透过禁制缝隙渗入的极致寒气中剧烈摇曳、黯淡，仿佛随时会熄灭。洞壁上凝结出一层厚厚的白霜，空气冷得吸入口鼻都带着刺痛。柳成脸色苍白如纸，呼吸急促，不仅仅是因为寒冷，更是因为那透过禁制依旧清晰传来的、如同天倾般的法相威压，让他气海初成的微薄修为几乎要停止运转，心脏狂跳，拳头攥得骨节发白。

叶牧谦盘坐如松，透过精心布置、兼具隐匿与窥视之能的防护禁制，将外界冰真人法相擎天、杀意盈野的景象尽收眼底。

“躲不得了。”叶牧谦平静地开口，声音不高，却瞬间驱散了洞内因外界压力而产生的恐慌气氛。

他眼神逐一扫过众弟子，清晰而快速地布置，“按既定计划行事。言冰，随我正面出去。瑶儿，依计暗中行事，务必小心谨慎，一击即走，不可恋战。千羽，你的任务是护好柳成，守住洞口这最后一道屏障，非到万不得已，不必出手暴露全部手段。”

“是，师尊！”三人齐声应道，声音斩钉截铁，毫无犹豫。

\vspace{2em}

山洞那块堵着门口的石头，在冰真人法相巨掌即将落下的千钧一发之际，悄无声息地、平稳地从内部搬开了。

两道人影一前一后从容步出，立足于这片被极致严寒统治的冰雪世界之中。

叶牧谦一袭简单青衫，布料普通，却纤尘不染。白言冰紧随其后，白衣胜雪，身形挺拔如孤松，右手自然垂于身侧，指尖距离腰间那柄古朴的法剑仅寸余之遥。两人就那样随意地站着，与前方那顶天立地的百丈玄冰法相、以及法相之下杀气腾腾的寒玉门众修，形成了体积与气势上极端悬殊、却又奇异平衡的对峙。

凛冽如刀的寒风卷起叶牧谦的衣角与发梢，他却恍若未觉，目光平静地越过那些张牙舞爪的寒玉门修士，直接落在法相之下、面色阴鸷的冰真人脸上。白言冰则眼神锐利如即将出鞘的剑锋，冰冷地扫视着周围每一个敌人，评估着他们的站位、气息与可能的攻击路线。

那柄天阶法剑上已有丝丝缕缕凝练到极致、锋锐无匹的无形剑意自行透出，将他周身三尺之内的空间悄然笼罩。那些试图侵入这个范围的刺骨寒气，一旦触及这无形的剑意领域，便如同撞上了看不见的利刃，悄无声息地被分割、湮灭。

“终于舍得从老鼠洞里爬出来了？”冰真人见正主现身，眼中贪婪与阴冷之色暴涨，几乎要化为实质。

他驱动百丈法相，声音如同万载寒冰相互摩擦、炸裂，带着直透神魂的寒意：“交出落星谷古传承与那医道秘法，然后自封修为，束手就擒！老夫或可大发慈悲，免去你们搜魂炼魄、神魂俱灭之苦！否则，今日此地，便是尔等葬身之所，形神俱灭！”

叶牧谦并未理会这色厉内荏的威胁，甚至没有去看那威慑力惊人的百丈法相。他只是微微侧首，用只有身旁白言冰能听到的音量，清晰而简洁地低语了一句：“按瑶儿此前探查所知，此人神魂有伤。我以心念牵制，你则放手施为，不必顾忌。”

白言冰微不可察地点了点头，按在剑柄上的手指，稳如磐石。

\vspace{2em}

下一刻。

冰真人只觉得自己的识海像是被一柄无形重锤狠狠敲击了一下，眼前猛地一黑，随即是无尽的晕眩与迟滞感汹涌而来，仿佛整个思维都被瞬间扔进了粘稠冰冷的沥青之中，运转艰难。

这还不是最可怕的，最令他魂飞魄散的是，神魂深处那处折磨他数十年、平日里需小心翼翼镇压的旧伤所在，毫无征兆地传来一阵尖锐至极、仿佛有冰锥在里面搅动的剧痛！与此同时，一种难以言喻的烦躁、心悸、甚至隐隐的恐惧情绪，不受控制地滋生、蔓延，疯狂干扰着他的判断与冷静。

仿佛有无数细密坚韧、无孔不入的心念之网缠绕上来，干扰着他与身后百丈玄冰法相那原本如臂使指的精妙联系，让法相的动作出现了显然的延迟与僵硬；遏制着他体内冰寒灵力如潮水般顺畅的奔流，使之变得涩滞；更如同最恶毒的刽子手，用看不见的手，精准而持续地撩拨、刺激着他神魂上最脆弱、最疼痛的那处伤疤！

叶牧谦也不下杀手。他能看出冰真人是被利用的，所以只是捆住冰真人，用神魂之刀使其彻底失去战斗力，并展开心念之网大大干扰剩余许多人的心神运转。

“这……这是什么邪门功法？！”冰真人心头巨震，骇然失色。

他闯荡修真界数百年，经历过无数厮杀，见识过各种诡谲法术、阴毒诅咒，却从未遇到过如此不讲道理、直接针对心神意识、干扰神魂本源的诡异手段！这完全超出了他以往的认知范畴。

他连忙收敛全部杂念，拼命固守识海，同时疯狂催动法相与体内灵力，试图以磅礴浩大的冰寒气息驱散这种无形无质的侵扰。然而，那把无形的刀锋随念而动，虽因他修为高深、心志也算坚韧，无法瞬间将其神魂击垮或操控，却让他如同身负千钧无形的枷锁，一举一动、一呼一吸都平添了三分滞涩与沉重。

十成的修为与战力，此刻竟只能勉强发挥出不到一半；更要时时刻刻分出一大半心神，去拼命压制那随时可能被引动、彻底爆发的神魂旧伤，可谓内外交困，苦不堪言。

剩下的修士也不好受，虽然比冰真人的感受要强些，但也出现了灵力运转不畅、法相动作迟延的问题。

就在冰真人被这诡异的心念领域牵制得难受万分、空有惊天法相却难以全力施展的同一瞬间，白言冰动了！

“锃——！”

清越如龙吟九霄的剑鸣陡然炸响，瞬间刺破了冰封世界的死寂！一道暗青色的剑光自他腰间迸发，并不如何璀璨耀眼，甚至显得有些古朴内敛，但其上蕴含的那股斩断一切、破灭万法的决绝锋芒，却让所有直视它的人，眼球都感到一阵刺痛！

白言冰身随剑走，人剑合一，化作一道撕裂寒风的白色惊鸿，以一种一往无前、睥睨天下的气势，直扑寒玉门修士最为密集的阵营核心！他选择的突击路线，恰好巧妙地避开了冰真人那百丈法相正面的绝对威慑区域，直指其侧翼。

“狂妄小辈！安敢如此？！拦住他！”冰真人虽被心念领域干扰得心烦意乱，但基本的判断与指挥能力尚在，见状又惊又怒，厉声大喝。

那三名早已蓄势待发的寒玉门法相境长老闻令，几乎同时暴起！

一人身形略胖，面色赤红，乃是法相外显期，身后隐约浮现一尊手持冰晶长鞭的怒目法相虚影。手腕一抖，那灵气长鞭如同毒龙出洞，带着刺耳的尖啸与冻结空间的寒流，密密麻麻的鞭影笼罩白言冰周身要害。

另一人身材高瘦，面容冷硬，是法相凝形期，他的法相相当奇特，是一面通体晶莹、宛如一整块万年寒玉雕琢而成的盾牌。盾面光华流转，寒气森森，防御力显然极其惊人，挡在了白言冰正前方，意图硬撼其锋。

第三人是个三角眼的老妪，同样外显期，她双手急速结印，十指翻飞如穿花蝴蝶，口中念念有词。顷刻间，方圆百丈内的水汽被强行抽取、冻结，化为无数根三尺长短、尖端闪耀着幽蓝寒光的锋利冰锥，如同被激怒的马蜂群，遮天蔽日，带着凄厉的破空声，朝着白言冰暴雨般攒射而下！

三位法相境长老，一远攻牵制，一正面防御，一范围覆盖，配合默契，瞬间便布下了天罗地网，誓要将这不知天高地厚的白衣剑修当场绞杀或擒拿。

然而，白言冰的剑，早已非吴下阿蒙。

通明境的修为，赋予他远超同阶灵气的凝练、纯粹、且与心神高度契合的真元。而他手中这柄由千年前一代炼器宗师风灵希亲手锻造的天阶法剑，更是将这份优势放大到了极致。剑锋过处，隐隐有风之音、冰之气相随，是剑器本身锋锐无匹、引动的天地微鸣，也是白言冰剑气外放的结果。

直刺、上撩、横斩、斜削。

但就是这简单的动作，在他手中施展出来，却有了化腐朽为神奇的威力！

“嗤啦——！”

长剑精准无比地点在漫天鞭影最核心的那一个点上。暗青剑光与冰蓝鞭影交击，没有惊天动地的爆炸，只有一声令人牙酸的金铁切割摩擦之音！那足以冻裂精钢、抽碎山石的寒髓长鞭，竟被这看似轻巧的一剑，生生斩出了一道深达寸许、灵光迅速黯淡的裂痕！持鞭的胖长老如遭雷击，闷哼一声，气血翻腾，鞭法顿时散乱。

剑势不停，顺势上撩，斩在那面厚实凝重的寒玉巨盾中央。

“铛——！！！” 一声洪钟大吕般的巨响震彻四野，盾牌表面光华狂闪，剧烈震颤！以剑锋落点为中心，无数道细密如蛛网的裂纹，伴随着令人心悸的“咔嚓”声，瞬间蔓延了小半个盾面！高瘦长老脸色一白，脚下不由得后退半步，眼中尽是骇然——他的法相本就专精防御，寻常法相境攻击根本难伤分毫！

与此同时，那漫天倾泻而下的冰锥暴雨已至头顶。白言冰甚至没有抬头去看，只是手腕微转，剑随身走，人与剑仿佛化作了一个浑然一体、急速旋转的青色光球。密密麻麻的冰锥撞入这光球三尺之内，便如同飞蛾扑火，连一丝声响都未曾发出，就被那绞碎一切的凌厉剑气彻底湮灭成最细微的冰晶粉末，纷纷扬扬飘散。

电光石火之间，白言冰以一敌三，仅凭一手返璞归真、快狠准绝的剑式，结合通明真元与天阶神剑之力，竟一人战平三名配合默契的寒玉门法相境长老！那精妙绝伦、已将剑道化简为繁又化繁为简的境界，看得远处山巅上一直观战的吴毓，都忍不住轻“咦”一声，一直平淡温和的眼中，首次流露出一丝毫不掩饰的惊讶与赞赏之色。

但吴毓的目光，很快便更多地聚焦在了白言冰手中那柄剑上，眼神变得灼热而探究：“好剑！剑意内蕴，风雷自生，锋芒含而不露，击坚则无往不克……这绝非寻常法器能有的气象，至少是天阶！”他捻着胡须，低声对身旁随从评论道，更像是在自言自语，“总不能说，是这几个来历不明的家伙背后，藏着一位早已隐世不出的天轮境炼器宗师，专门为他们炼制吧？ 还是说……这剑本就是那传承中的一部分，被他们捡到的？

不等两名随从回答，他又自言自语，像是在回答自己的问题：“可天阶法器，尤其是这等品质的杀伐之器，也是能随手捡来的机缘么？”

战场上，剩下的那十余名真符境寒玉门精锐，原本在三位长老出手时，已结成小型战阵，杀气腾腾地准备一拥而上，协助擒敌或阻断退路。然而，眼前这完全超出预料的战况，让他们瞬间僵在原地，满脸的难以置信与惊骇！

三位平日里在他们眼中高高在上、神通广大的法相境长老联手，竟然被对方一人一剑正面对抗还能斗成平手？

那笼罩全场、令他们灵觉刺痛、心烦意乱、本能感到恐惧的诡异领域，更是如同无形的沼泽，拖慢了他们的反应，瓦解着他们的斗志。

一时间，这些人竟骇得不敢上前，只敢在外围挥舞着兵器，虚张声势地呼喝，却无一人敢真正踏入那白衣剑修剑气笼罩的范围，更别说去冲击一直静立不动，只是与冰真人隔空对峙、却给人以深不可测之感的叶牧谦了。

至于那些境界更低、只是来壮声势或搞后勤的弟子们，早已被这等级别的交锋吓得魂飞魄散，远远退开，生怕被一点余波扫中，便化为齑粉。

战场，竟一时陷入了诡异的僵持。寒风呼啸，冰晶闪烁，杀意凝而不发，唯有剑气破空与冰盾哀鸣之声，偶尔打破这片冰封世界的死寂。

\vspace{2em}

然而，就在山洞前方战况激烈之时，一道几乎透明的虚影，借助山林阴影和精妙的千幻流光诀，早已悄然脱离了战场，向着流云山脉外围疾驰而去。

正是奉命暗中行事的沈瑶。

她的目标是寒玉门此次行动设在山脉外围一处隐秘山谷中的临时补给与联络营地。那里留守的人员不会太多，但一旦被捣毁，足以进一步打击众多低阶修士的士气。

凭借高超的潜行技巧，沈瑶如鬼魅般潜入寒玉门外围营地。营地中留守的不过五六名灵枢境、玄脉境修士，正在闲聊，浑然不觉有人靠近。沈瑶忽然在营地的东南角，发现了一个小型、正在微微发光的传讯法阵，旁边还有几枚使用过的传讯玉符残片，销毁得不甚彻底。

沈瑶寻思，如果能找到些线索，那自然是更好的。

她目光一凝，悄然隐匿身形靠近，捞起那些玉符残片就跑。随即再次飞往高空，取出得自落星谷的那面八角铜镜，镜面对准那传讯法阵。

注入一丝真元，铜镜背面星辰图案微亮，镜面荡起涟漪。

下一刻，那传讯法阵周围的空间微微扭曲，景象骤然变化！仿佛有无形之手将其中的灵力结构粗暴搅乱，又幻化出数道凌厉的虚影剑光劈砍而下！

“咔嚓”、“噗嗤”几声轻响，传讯法阵光芒骤熄，阵基碎裂，化为齑粉。整个过程几乎无声无息，且看起来像是遭受了某种突然的、无差别的范围性能量冲击或是内部故障。

沈瑶得手后毫不留恋，身形再度融入阴影，迅速远离；只留下那群在下面没反应过来发生了什么，拿起武器吱哇乱叫的低阶修士。

“快，快去报冰真人！”

\vspace{2em}

山洞外，冰真人百丈玄冰法相的威压如同十万丈深海下的寒流，不仅冰冷刺骨，更带着碾碎神魂的沉重。即便隔着一层隐匿防护禁制，那威压依旧丝丝缕缕地渗入石隙。

洞内，原本稳定跃动的篝火，火苗猛地一矮，色泽变得惨淡泛青，仿佛下一刻就要冻结成冰，光晕在岩壁上投下的影子也随之扭曲颤栗，映得柳成的脸半明半暗，苍白得没有一丝血色。

他耳朵里灌满了外面的毁灭之音：玄冰凝结又崩裂的脆响，如同冰山倾覆；剑气纵横撕裂空气的尖锐啸叫；还有冰真人那仿佛万载寒冰相互摩擦碾压的怒喝，每一个音节都带着冻裂魂魄的恶意。

更深处，一种超越听觉的感知在震颤——师尊叶牧谦那里，传来的并非惊天动地的轰鸣，而是一种浩瀚如星空、平静如古潭的心念波动，它无形无质，却稳稳抵住了外界滔天的寒潮；师兄白言冰的方向，则是一道凌厉睥睨、宁折不弯的剑意，如同刺破黑暗的极光，每一次闪烁激荡，都让外部的冰寒领域为之动荡。

无力感，化作无数冰冷滑腻的藤蔓，死死勒住他的心脏，每一次心跳都带来窒息般的疼痛。

他只是个刚刚转修问道途的养气境小修士。在法相境大能那引动天地之威的恐怖存在面前，他渺小得连尘埃都不如。沈瑶师姐行险潜入敌后，师尊与师兄正面抗衡强敌，就连守护洞穴的千羽师姐，修为也远超于他。

唯有他，被安置在这最安全的角落，像一个需要被严密保护的累赘。

“我必须做些什么……不能只是看着，不能只是等着……”

这个念头如同岩浆般在他心底翻滚、沸腾，灼烧着他的五脏六腑。

对自身弱小的痛恨，对师尊师兄安危的焦灼，还有传承自《焚心器道》《明心器道》两部大作的责任感，混合成一股近乎暴烈的冲动。他眼前闪过师尊改良“焚心”为“明心”时，那手指轻抚材料，万物随之律动的玄妙景象；以及，风灵希留在玉简中那句看似平淡，此刻却重若星辰的话语——“心念微动，火自内生”。

他的目光，猛地投向了角落那堆从落星谷得来的矿石。星纹铁闪烁着细碎的银蓝星光，空冥玉流淌着如梦似幻的氤氲光晕，其他一些晶体也折射着瑰丽的色彩。

然而，在这些璀璨之中，几块约莫拳头大小的青黑色玉石，却显得格格不入。它们表面粗糙，毫无光泽，像是河床底被冲刷了千万年的顽石，甚至带着一种笨拙的丑陋，静静躺在珍宝之中，仿佛在刻意隐匿自己的存在。

这是风灵希藏品里连柳成都无法辨认的少数几种矿石之一，当初带走，仅仅是因为它们坚硬得不可思议，并且指尖触碰时，会有一种极其微弱的、仿佛要将神识轻轻粘住的吸附感。

鬼使神差地，柳成伸出手，捡起了其中一块青黑色玉石。入手并非纯粹的冰凉，而是一种沉甸甸的、仿佛握着一段凝固岁月的实质感。那股微弱的吸附感再次传来，这次更清晰了些，像是有无数细小的、冰冷的针尖，在轻轻探触他指尖流转的微弱元气和逸散的精神。

一个近乎疯狂却又带着燎原之势的念头，如同闪电般劈中了他：师尊能以通明境心念“浸润”引导材料本质，我远远不及，但我有刚刚诞生的气海元气，有不顾一切的决心！

这石头……它能吃神念。我要是模仿师尊“因势利导”的意境，去“诱发”它本身的反应……会不会，哪怕只有一瞬，能创造出一点点干扰？哪怕只是让外面那尊冰疙瘩的光芒暗淡一丝，为师尊师兄争取一线微不足道的空隙也好！

“柳成？你要做什么？”洞口附近，正全神贯注感应外界、周身已有淡金色护体真元隐约流转的陈千羽，敏锐地察觉到了他气息和情绪的剧烈波动，蓦然回首，清丽的眼眸中满是惊疑与担忧，秀眉紧紧蹙起。

“千羽师姐，我……我想试试。”柳成的声音干涩沙哑，仿佛砂纸摩擦，那是极度紧张下喉骨的痉挛。然而，他的眼神却在此刻亮得骇人，像是燃尽一切的火种，带着一种近乎偏执的、豁出性命的决绝。

“我不能干等着。这块石头……它有点特别，我觉得……它能回应我。”

陈千羽的心猛地一沉。她深知这个小师弟骨子里的执拗，但更无比清醒地认识到，他想做的事情，与他目前的修为之间，横亘着何等可怕的天堑。

“太危险了！师尊的心念境界，岂是你能轻易模仿的？那是水到渠成，不是蛮干硬闯！外面有师尊和言冰，我们的任务是守住这里，不添乱就是……”

她急促劝阻的话语尚未说完，柳成已经猛地盘膝坐下，将那块青黑色玉石郑重地置于身前布满尘埃的地面，双手虚悬其上，掌心遥遥相对，闭上了眼睛。

意识沉入丹田，竭力调动起气海中那团初生不久、尚且稀薄微弱的元气漩涡，同时强迫自己摒除所有杂念——对战斗的恐惧、对自身弱小的不甘、对成败的忧虑——将全部的精神，所有的意念，如同压缩到极点的弹簧，聚焦于掌心之下的顽石。

\vspace{2em}

山洞前的战斗仍在继续。冰真人越打越是心惊，叶牧谦那无形的心念干扰让他烦躁不已，白言冰的剑更是凌厉得超乎想象，三名长老已有一人负伤，但他的剑势却越发凶猛。

久攻不下，冰真人心中焦躁更甚。

就在这时，一名留守营地的寒玉门修士连滚爬跑地从前线后方赶来，脸色惨白地禀报：“长……长老！不好了！我们设在东南山谷的临时营地……突然遭到不明袭击！传讯法阵和所有记录都被毁了！留守弟子只看到光影乱闪，没看清是什么人干的！”

“什么？！”冰真人大怒，心神又是一阵波动，玄冰法相都晃动了一下。

后院起火！难道是对方还有埋伏的人手？还是有其他势力插手了？

就在他惊疑不定、攻势为之一缓的关头，远处天际忽然传来一阵低沉的嗡鸣声。只见一辆装饰华贵、流线优美、通体散发着柔和灵光的云车，正破开云层，朝着战场方向疾驰而来。云车侧面，清晰地绘制着郑氏家族的徽记！

“郑家的人？”冰真人眉头紧锁，心中升起不祥预感。

郑天河不是说好不出面，由他全权处理吗？这来的又是谁？

云车悬停于战场边缘，舱门打开，一道身影跃下，凌空踏步而来。来人一身郑氏长老服饰，面容儒雅中带着几分久居上位的沉稳，正是郑季！

郑季目光扫过狼藉的战场，在叶牧谦和白言冰身上略作停留，最后看向冰真人，拱手道：“冰魄真人，且慢动手！”

冰真人收了收法相威势，冷声道：“郑季长老？你来此何意？莫非郑天河公子改了主意？”

郑季脸上露出恰到好处的无奈与劝解之色：“冰魄真人误会了。老夫此来，非是奉天河之命，实是不忍见同道相残，传承蒙尘。更不愿见我郑氏领地之内，酿成无法挽回的血案。”

他顿了顿，声音提高，确保双方都能听清，“叶先生，冰魄真人，今日之事，不过是一场误会。依老夫愚见，不如双方各退一步，就此罢手如何？”

他看向叶牧谦，语气诚恳：“叶先生，落星谷古传承之事，如今已传得沸沸扬扬，怀璧其罪，终非长久之计。不若将其部分精要公之于众，或与天工阁这般专研技艺的势力合作研习，既可平息纷争，亦可造福炼器同道，岂不两全？”

又转向冰真人：“冰魄真人，你为疗伤之心，老夫理解。但强取豪夺，非正道所为，亦有损寒玉门清誉。不若由老夫做个中间人，请叶先生这位高徒陈姑娘为你诊治一番，或可缓解伤痛。至于传承之事，可从长计议，何必刀兵相见？”

郑季这番话，听起来冠冕堂皇，一副和事佬的模样。但叶牧谦瞬间便明白了他的真实意图——郑季哪里看不出来这寒玉门是受人指使？但如果传承真叫寒玉门以及其背后的郑氏主脉夺了去，那算他自己任务失败，叶牧谦等人肯定不会把传承给他。而郑氏主脉，自然更是不会给他核心传承了。

他此番前来，名为劝和，实则是要搅乱郑天河的计划，保住传承不落入主脉之手，同时为自己日后接触、谋取传承创造机会和借口！他甚至主动提到了天工阁，既是为了增加说服力，或许也存了借天工阁之势平衡郑天河的心思。

冰真人闻言，脸色变幻不定。他看了看依旧深不可测的叶牧谦，又看了看剑势逼人的白言冰，再想到刚刚被毁的营地，心中已然萌生退意。继续打下去，未必能讨到好，反而可能两败俱伤。郑季的提议，至少给了他一个台阶下，还能有机会接触到那位医道弟子……

就在冰真人踌躇间，又一个声音响起了。

“郑季长老此言，甚合我意。”

只见远处山巅上，三道身影飘然而下，正是观望许久的吴毓及其随从。吴毓笑容可掬，朝着双方拱了拱手：“老夫天工阁吴毓，恰逢其会。争斗确实无益，尤其是为了可能蕴含先贤智慧的古传承。若因争斗而损毁了传承，或伤了人才，那才是真正的损失。不若暂且罢手，坐下商谈。我天工阁愿做个见证，也愿为传承的研究与保护，略尽绵薄之力。”

吴毓的出现，让冰真人彻底失去了继续战斗的底气。

天工阁！这可是连郑氏都要给几分面子的庞然大物！他们明显是冲着传承来的，自己若再坚持，岂不是同时得罪郑氏长老和天工阁？

冰真人脸色铁青，心中将郑天河骂了无数遍，但形势比人强。他狠狠瞪了叶牧谦一眼，又看了看劝架的郑季和吴毓，终于借坡下驴，冷哼一声：“既然郑长老与吴阁老都如此说……罢了！今日便给二位一个面子！”

\vspace{2em}

山洞内。

柳成一直试图用那粗糙的神识去触摸玉石内部的纹理，用微弱的元气波动去叩响那层微弱的吸附之力，然后，带着恳求与引导，将自己的元气与心神，缓缓递送过去，想象着它们如同春雨，渗入干涸的土地，去唤醒某种沉睡的本能——就像他曾惊鸿一瞥目睹的师尊的手法：因势利导，润物无声。

然而，柳成那点可怜的控制力，在叶牧谦的通明境界面前，用“萤火与皓月争辉”来形容都有些不够。他那笨拙、强烈却毫无章法的引导意念，非但没能唤醒玉石内任何温和的回应，反而像是一脚狠狠踹在了一扇无比凶险的禁忌之门上！

就在他那稀薄的元气与凝聚的心神，颤巍巍地扭转了玉石核心某处纹路的刹那，那股原本只是微微吸附的感觉，骤然坍缩、旋转，化作一个恐怖绝伦的吞噬漩涡！

这下不是他在引导玉石，而是这青黑色的石头，瞬间变成了一头苏醒的贪婪凶兽，反过来疯狂地吞噬他的一切！

气海中那点可怜的元气，连挣扎的余地都没有，如同溃堤的洪流，汹涌倾泻而出。更可怕的是，他的神念、意识、甚至生命本身，都像是被无数冰冷的黑色触手死死缠住、咬碎，然后无可抗拒地拖向那青黑色玉石张开的无底深渊！

“呃啊——！”一声短促扭曲、不似人声的惨嚎刚从柳成喉间挤出半截，便戛然而止。他只觉得眼前猛地一黑，灵魂仿佛要被从躯壳里硬生生抽离、碾碎，身体瞬间失去了所有力量，如同被抽去骨骼的皮囊，软软地向一侧歪倒。

与此同时，他手中那块青黑色玉石，骤然大放幽光——深邃、混乱、仿佛能吞噬周遭所有光线与声音的幽暗漩涡，在石头表面急速流转膨胀！

“柳成！”陈千羽脸色瞬间惨白如纸。

柳成的气息如同风中残烛般骤降、飘摇欲灭，面色死灰；周身灵气波动更是混乱狂暴到极点，如同沸腾的油锅！

出大事了！

而且是足以瞬间致人魂飞魄散的天大祸事！

她几乎瞬移般掠至柳成身边，素手疾挥，带起一片残影，指尖早已夹住的数根细如牛毛、却闪烁着淡金色玄奥纹路的金针，以肉眼难辨的速度，精准无比地刺入柳成百会、膻中等几处关乎神魂稳固的生死大穴，强行钉住他即将溃散离体的神魂！

同一时间，她的另一只手掌死死的扒开柳成的手掌，把那块玉石直接抠了出来，然后握住他的手掌，向空虚的柳成强行渡入她自己的真元。

至少先续住他的命！

或许是陈千羽的及时介入，暂时平衡了那可怕的吞噬；或许是柳成在意识彻底沉沦前，那拼死一搏、不甘就此沉寂的原始意念，与玉石深处某种机制产生了无法理解的共鸣；又或者，这本就并非一块纯粹带来毁灭的邪物。

那恐怖的吞噬过程，在几乎将柳成掏成空壳的临界点，骤然停滞了一瞬，然后发生了谁也预料不到的逆转！

它仿佛饱足了柳成献祭般的元气、神念乃至部分生命印记，内部某种亘古沉睡的机制被彻底激活。一种连陈千羽都难以理解其原因、但深入生命本源层次的链接，以玉石为中心，强行与柳成建立起来，那是一种近乎血脉共生、魂魄相连的紧密联系。

紧接着，这块被从柳成手里强行取下的玉石，竟然自己浮动起来……

\vspace{2em}

刚刚收敛了百丈玄冰法相，满心憋屈与不甘，正欲转身离去的冰真人，脚步尚未迈出，便被身后这骤然爆发、改天换地般的剧变骇得浑身寒毛倒竖，猛然扭回头！

只见他刚刚放弃攻击、认为已无大利可图的山洞所在山峰，此刻已被浩瀚如海、璀璨夺目的灵气潮汐与天地异象彻底淹没！那高悬于天的千丈灵气漩涡，那垂落涤荡的瑞气霞光，那弥漫开来、让他神魂旧伤都感到一丝舒缓的古老道韵……一切的一切，其规模之宏大，气象之纯粹，道韵之玄奥，远比他方才全力施展法相神通时引动的天地寒气，要恢弘壮丽何止十倍！

原本覆盖数里的铅灰色冰寒云气，被一股无可抵挡的伟力从中心粗暴地撕开、驱散！

一个直径超过千丈、缓缓旋转的巨型灵气漩涡，在高空赫然成型，漩涡中心垂直指向下方山洞，如同连接天地的漏斗！漩涡之中，电蛇狂舞，雷声沉闷；霞光万道，交织流淌，将整片天空渲染得如同神国降临；千条瑞气如匹练垂落，氤氲蒸腾，散发出令人心旷神怡的清香。

一股宏大、古老、苍茫，仿佛经历了无数纪元沉淀的磅礴道韵，伴随着这席卷天地的灵气潮汐，无声无息地弥漫开来，笼罩四野。在这道韵笼罩下，方才战斗遗留的肃杀、冰寒气息被迅速涤荡，甚至一些被波及而枯萎的草木，都隐隐有嫩芽挣扎欲出。

以山洞为中心，方圆数十里内的流云山脉，天地灵气彻底暴走！

肉眼可见的，赤色的火灵、青色的木灵、黄色的土灵、白色的金灵、黑色的水灵，乃至更多色泽难以定义的变异灵气光点，如同被无形的巨手从沉睡中粗暴拽出，从每一寸山岩、每一片树叶、每一道溪流、甚至虚空中本身析出、剥离！

它们起初是星星点点，旋即汇成溪流，溪流奔涌成河，最终化作一道道绚丽夺目、色彩斑斓的庞大灵气洪流，如同朝拜君王的臣民，又如同扑向火焰的飞蛾，从四面八方，以山洞口为目标，疯狂地奔涌汇聚！

洞内瞬间被浓郁到化不开的灵光充斥，呼吸间都感到灵力澎湃。

\vspace{2em}

这绝非任何法相境修士能够人为制造的场景，甚至道种境、天轮境都做不到！

\vspace{2em}

“这……这是？！”冰真人瞳孔缩成了针尖大小，心脏在胸腔里疯狂擂动，几乎要撞碎肋骨跃出咽喉。

一股灼热到让他灵魂战栗的贪婪，混合着排山倒海般的悔恨，瞬间冲垮了他残存的理智。他能清晰地感觉到，那异象核心处隐隐散发出的能量层次与法则韵味，高渺深邃得令他这法相境修士都感到自身渺小与心颤！

不远处，原本智珠在握、现身劝和的吴毓，那从容不迫的姿态早已荡然无存。他也猛得回头，脸上每一道皱纹都写满了极致的震惊与难以置信，嘴唇哆嗦着：“天地灵气自发灌注，万道霞光为其洗礼，法则道韵随之显化……这、这至少是天阶上品，不，是天阶极品法器出世都未必有的气象！甚至……这煌煌天威，这涤荡道韵，难道是圣阶？！”

以吴毓在炼器一途上的见识，天阶法器确实见了几十件，也亲自参与过其中几件的炼制或修复工作；但天阶法器也已经是灵界各大势力争抢之物，就连中州皇朝所持有的天阶法器估计也不到十件。相对的，他没见过任何圣阶法器：至少在灵界，那是传说中的传说，每一件都蕴含莫测威能，足以镇压一方超级大教气运万载不朽！

一旁的郑季，同样满脸愕然，张着嘴，半晌发不出声音。他完全没预料到，山洞内除了叶牧谦师徒几人，竟还藏着能引发如此惊天动地变故的后手！

这声势……哪里还是后手？简直是捅破了天！

冰真人死死盯着那灵气漩涡的中心，目光炽热得几乎要喷出火来，脸色在青、红、白之间急剧变换，胸口旧伤处传来一阵阵撕裂般的绞痛，呼吸都变得无比困难，仿佛溺水之人。

“早知道……早知这山洞里竟藏着能引动如此异象的宝物或逆天传承……老夫说什么也要拼到底啊！”一个疯狂的声音在他脑海中咆哮：拼着那道困扰他多年的神魂旧伤彻底爆发，拼着与郑季、吴毓乃至他们背后的势力当场翻脸血战，也一定要攻破山洞，将那引发异象之物夺到手中！

若能得此圣物，寒玉门何愁不能崛起！自己那棘手的神魂旧伤，说不定就能借此蕴含的道韵一举治愈！甚至……借此突破那道种境界的壁垒，也未必是奢望！

要是这样，他还用得着郁郁久居郑氏之下！

巨大的悔恨如同千百条毒蛇，钻心噬骨，啃噬着他的五脏六腑，让他喉头一甜，几乎真要喷出血来。

\vspace{2em}

吴毓到底是见惯大风浪的天工阁高层，深吸了好几口气，勉强压下心中的惊涛骇浪，但那眼神中的热切与震撼却丝毫无法掩饰。他率先开口，声音因激动而带着明显的变调和干涩：“叶……叶先生，敢问洞中……究竟发生了何事？这、这惊世异象……”

此时白言冰已经得了叶牧谦指令，迅速遁回山洞；沈瑶估计也已经回洞里了。

叶牧谦目光扫过空中那仍在缓缓旋转、倾泻下磅礴灵气的巨型漩涡，再感受着那弥漫天地、却与山洞内柳成气息隐隐水乳交融的奇异道韵，心中瞬间转过无数念头。再加些探查，已有了七八分猜测。

他面上却不动声色，甚至恰到好处地流露出一丝无奈与淡淡的责备，平静回答道：“让吴阁老见笑了。是我那新收的不成器小徒柳成，方才见外面情势紧张，少年心性，情急之下胡闹了一番，不想竟鲁莽引动了灵气躁动，闹出这般大动静，惊扰各位了。”

“柳成？那个……那个修为不过灵枢境的炼器小子？”冰真人如同被踩了尾巴的猫，失声尖叫起来，眼珠子暴突，几乎要从眼眶里掉出来。“他？就凭他？引动这等……这等堪比圣物出世的天地异象？胡闹！”

他想狂吼“这分明是圣器现世”，可话到嘴边，看着叶牧谦那平静无奈的表情，又觉得此事荒谬绝伦到了极点，硬生生咽了回去，憋得脸色更显灰败。

吴毓也是满脸的不可思议，仿佛听到了天方夜谭：“叶先生，您是说……引发这笼罩方圆数十里、引动天地道韵的异象的……是一件器物？而且是高徒柳成……以区区灵枢境的修为，炼制或……激发出来的？”

叶牧谦微微颔首，故意轻叹一声，语气带着师长对莽撞弟子的责备与后怕，也是对误解的一番解释：“我等所修，非灵枢径，而是‘问道途’。柳成不过新近转修，而问道途的第一境界便是‘养气’。这孩子根基都未稳，便学着我的手段，强行炼器逞能，险些榨干真气，酿成大祸。幸得小徒千羽精通医道，及时护持，才勉强保住了他的根本，未伤及本源。”

叶牧谦目光扫过空中渐息的异象，继续道：“至于这异象……叶某实际也不知。需待他稳定后，细细探查方能知晓。但叶某猜测，估计是他用了某些特殊材料。”

这番话听在吴毓、冰真人这等老江湖耳中反而比直接承认柳成炼出圣器要可信得多。

但冰真人听得此言，只觉得心口那股绞痛更剧烈了，脸色灰败如同死人。

一个灵枢境的小蚂蚁，用一块可能是圣阶的逆天材料，误打误撞，就弄出了他这法相真人都要仰望的天地异象？这运气……这机缘……凭什么就落在那么个毛头小子头上？！

郑季则是眼神急剧闪烁，心中念头飞转如电。柳成此子，果然不能以常理度之！这份“误打误撞”便能引动天威的能耐，其背后代表的隐秘气运或特殊天赋，恐怕比他明面上展现出的炼器天赋，更加惊人，也更加值得投资。叶先生门下，当真是深不可测，藏龙卧虎！

吴毓毕竟是天工阁阁老，心性沉稳，很快再次强迫自己从极致的震撼中冷静下来，分析利弊。但眼中的热切与重视，却提升到了前所未有的高度。他深深看了叶牧谦一眼：原来你叶牧谦也是个炼器师，而且水平比这柳成只强不弱，估计至少是个天阶炼器师！

随即，吴毓郑重拱手道：“叶先生门下英才辈出，气运惊世，实在令人叹为观止。不论具体缘由如何，能引动此等惊世异象，高徒柳成，乃至其所关联之物，必是非同小可。”他斟酌了一下词语，试探道：“不知吴某可否有缘……稍作瞻仰？”

他想说“仔细探查”，但话到嘴边换成了更客气的“瞻仰”，毕竟涉及他人核心隐秘，太过直接反而落了下乘。

叶牧谦自然明白他的意图。但此刻洞内情况未明，柳成生死安危尚在未定，岂容外人靠近？

他面上露出恰到好处的歉意与坚持，婉拒道：“吴阁老厚爱，叶某心领。只是小徒方才历经生死险关，此刻神魂肉身皆极度虚弱，状态极不稳定。洞内因方才变故，气息混乱狂暴，残留的灵力乱流与未散的道韵碎片，恐怕对旁人亦有风险，实在不宜外人进入探查。”

吴毓虽然心中非常想亲眼看看那引发异象之物，甚至看看柳成本人，但也知叶牧谦所言在理，且态度坚决，强求不得，反而可能恶了关系。他按下急切，点头表示理解：“理应如此，是吴某心切，唐突了。”

他心中却早已打定主意：不惜代价，也要将叶牧谦一行人——尤其是叶牧谦和柳成——牢牢绑定在天工阁的战车上。他们身上蕴含的潜在价值，甚至可能超越一件已知的天阶法器！

叶牧谦担心柳成安危，于是又客套了几句，便回去了。

冰真人见叶牧谦已然回返，郑季与吴毓明显也站在对方一边，自己再留下去不仅没有丝毫机会，反而要承受那天地异象引来更多目光后的风险。他只能将满腔几乎要化为实质的悔恨与贪婪，死死压在心底，化作一声充满憋屈与愤懑的冷哼。

再也懒得维持半点风度，对着门人弟子低吼一声：“走！”随即，袖袍一卷，带着寒玉门众人，化作数道仓促的冰寒遁光，头也不回地朝着远离流云山脉的方向狼狈遁去，那背影在尚未完全散尽的霞光映照下，显得格外灰暗颓丧。

郑季也知此时不是深谈之时，今日变故太多，信息量巨大。他亦向吴毓拱手告辞，临走前，又深深看了一眼那灵气渐平的山洞，目光复杂难明。

吴毓是最后离开的。他最后仰头望了一眼那正缓缓收缩消散、却依旧残留着震撼余威的天地异象，转身对两名随从无比郑重地说道：“今日之事，动静实在太大。流云山脉并非僻壤，此刻恐怕方圆数百里内的修士都已有所感应，不日便会传遍东域，乃至引起更遥远势力的关注。务必把此事上报至总舵，并提高至最高优先级。”

两名弟子齐声应诺，随即一行人也就此离去。

\vspace{2em}

天空中，那巨大的灵气漩涡终于缓缓消散，最后一丝霞光隐入天际，漫天瑞气也回归天地。唯有山林间，灵气浓度陡增了数倍，草木格外葱翠，一些灵花异草甚至提前绽放；无数被惊动的飞禽走兽，仍在不安地嘶鸣奔窜。

而远方天际，已能隐约感觉到一些强弱不一、带着惊疑与探究气息的神念波动，正在小心翼翼地朝着这片区域扫来，显然是被方才那惊天异象吸引而来的其他修士。

但他们显然是连口汤都喝不上了。

\vspace{2em}

洞内。

刚刚天地灵气一番灌注，那青黑色玉石却没有给这些灵气全部吃下。洞内灵气充盈得过分，甚至要在洞壁上凝出液体，呼吸间都带着清灵之感。

叶牧谦回来的时候，看到白言冰和沈瑶正研究着一块青黑色玉石。陈千羽扶着刚刚从深度昏迷中苏醒、连坐直身体都勉强、虚弱得如同大病一场的柳成，小心地将一枚专门温养受损神魂、弥补元气精血的淡青色丹药喂入他口中。丹药入口即化，一股温和的暖流散向四肢百骸，柳成灰败的脸上这才有了一丝血色。

而那块引发了一切变故的青黑色玉石，此刻正安静地悬浮在柳成身前尺许处的空中，颜色似乎比之前深邃内敛了许多，表面流转着一层温润莹莹、仿佛活物呼吸般的柔和幽光。

光芒的明暗节奏，竟隐隐与柳成微弱的心跳同步。

它散发出的气息，不再是冰冷死寂或贪婪狂暴，而是一种奇异的、安宁祥和的、仿佛与柳成生命韵律浑然一体的波动。

“师尊……”柳成看到叶牧谦返回，挣扎着想要起身，却被叶牧谦一步上前，轻轻按住肩膀。

“莫动，静心调息。”叶牧谦的声音平和而安定人心。同时他已伸出两指，轻轻搭在柳成腕脉之上，一缕精纯通明、细致入微的心念顺着接触渗入柳成体内，开始全面探查他的伤势与变化。

“感觉如何？可有哪处特别不适？”

“弟子……弟子只觉得……空……非常空……”柳成声音虚弱，断断续续，“好像身体被……被掏空了一大半，骨头都是软的，脑子里也昏沉沉的，想集中精神都难。”

他吃力地眨了眨眼，目光落在那悬浮的玉石上，眼神复杂无比，有心有余悸的后怕，也有一种莫名诞生的亲近与好奇。“但……但又好像……和这块石头，有了一种……很奇怪的联系。我说不清，好像它能感觉到我，我也能……模模糊糊感觉到它，不靠眼睛，也不靠神识。”

沈瑶已先一步对那小块玉石早已有了些探查：“奇怪……这能量波动……温和是温和了。但这强度似乎并不高啊，远达不到天阶法器应有的那种磅礴内蕴的压迫感。甚至感觉其目前的能量层次，仅仅相当于一件品质不错的灵阶法器？但方才外面那引动方圆数十里灵气、道韵天成的惊天异象……”

她百思不得其解，甚至尝试着分出一缕极其细微柔和的真元，轻轻触碰玉石表面。玉石只是微微一颤，表面的幽光略略明亮了一丝，反馈回来的力量确实平和，但也确实不强，与那撼天动地的出世景象完全不符。这与那个天阶法镜相比，何啻云泥之别，甚至连真影刃都不如。

叶牧谦此时已收回探查柳成的手，示意陈千羽继续照顾柳成调息。他伸手虚招，那块悬浮的玉石便温顺地落入他掌心。

入手不再是之前的冰凉沉实，而是带着一种与柳成体温相近的微温，那股与柳成同呼吸、共命运的奇妙联系感，在他通明心念的感知下，变得更加清晰而深刻。

他闭上双眼，全部心神沉入其中。

这一次，他看到了柳成和沈瑶都无法感知的真相。

这玉石内部的结构，根本就不是寻常矿物晶体的排列，也不是任何已知法器的阵法符文烙印。

它是活的。

至少是一种极其高明玄奥的生命雏形。

叶牧谦又迅速翻查从风灵希石室得到的其他物品，找到了一块外观类似的青黑色玉石进行对比。

那一块虽然坚硬，也有微弱吸附感，但叶牧谦能够断定它是死的。

那事实就昭然若揭了。

柳成方才那番不计后果、近乎献祭的元气与神念冲击，虽然是一把错误的钥匙，却阴差阳错地捅开了最核心的锁。然后这生命便吸收了柳成的一部分生命印记与心神烙印作为最初的火种，并与其建立了最原始、最紧密、无法分割的共生联系。

“此物……”叶牧谦睁开双眼，缓缓说道，一向平静的眼眸中也忍不住掠过一丝惊叹，“并非死物法器，至少现在不再是了。它是活的，或者说，它拥有了成长为某种活物的根基与潜质。”

他看向柳成，语气严肃中带着感慨：“它与你，柳成，已是同呼吸，共命运。一荣俱荣，一损俱损。它的品阶、未来的威能，恐怕并非固定不变，而是会随着你的成长、你与它的共同温养而逐步蜕变、进化。方才那席卷天地的异象，倒像是它从死物变为活物这一根本性蜕变时，鲸吞了大量天地灵气形成漩涡，进而引动了深层次天地法则的共鸣与庆贺，可以看作是天地为它的新生降下的洗礼。”

“‘先天’伴生之灵？成长型法器？”白言冰闻言，冷峻的脸上也露出一丝讶色。“这玩意不是话本材料吗？”

“目前看来，倾向于此。”叶牧谦将温热的玉石轻轻放回柳成手中。“但它此刻还很弱小，如同初生的婴儿，需要你以自身元气、心神，乃至对道的感悟，去慢慢温养哺育。”

“你方才的行为，凶险万分，实乃取死之道！若非千羽及时出手，你早已被它吸干一切，神魂俱灭，真灵不存！”他的语气转为严厉：“记住，炼器之道，亦是修行之道，讲究量力而行，循序渐进，厚积薄发！切不可再如此鲁莽冲动，以命换器！机缘背后，往往是等同的凶险！”

柳成双手捧住那块与自己血脉相连的玉石，感受着那仿佛另一个心跳般的微弱脉动，想起方才意识沉沦前那无边无际的冰冷与绝望，后背瞬间又被涔涔冷汗浸透。连忙挣扎着恭敬应道：“弟子……弟子明白！弟子一定谨记师尊教诲，绝不敢再如此不知死活！”

那种灵魂都要被撕碎抽离的恐怖，一次就够了，足够成为他修行路上最深刻的警钟。无论如何，活着才是根本。

“给它起个名字吧。”陈千羽见柳成知错，气氛有些沉重，便微笑着开口缓和，目光柔和地看着柳成和他手中的玉石。“它因你而生，与你一体，当有个名号。”

柳成低头，抚摸着玉石温润的表面，那幽光随着他的触摸微微流转，仿佛在回应。他思索片刻，轻声道：“它因我模仿风灵希前辈《焚心器道》中的决绝之意而现世，也差点让我真正‘焚’了心。 就叫它……‘焚心印’吧。既是纪念风灵希前辈的传承，也是时刻警醒我自己，日后炼器修道，绝不可再如此急功近利、行险侥幸，需得脚踏实地，一步一个脚印。”

“焚心印……名实相副，很好。”叶牧谦点了点头，眼中露出赞许之色。

玉石似乎也很高兴，表面的幽光又变亮了几分。

能于生死危机后有所领悟，而非一味沉溺后怕或狂喜，这份心性，亦是难得。

\vspace{2em}

见此事告一段落，沈瑶又拿出了他在寒玉门行动营地获得的那几个销毁不彻底的传信玉简：“师尊，你看看这个。”

叶牧谦接过玉简，把其中的内容简单拼凑了一下，面色顿时难看了起来：“这个郑天河，原来还想着要我们的命？”

沈瑶凝重地点了点头。

叶牧谦思索了一下，道：“目前来看，他借刀杀人的计划是失败了，但依然不可不防。不过他们敢于这么做，那确实有失公义。这些玉简残片，你先保留着，它们或可成为我们将来的利器。”

沈瑶沉声应诺。

\vspace{2em}

经此一役，叶牧谦认为短期内没人敢招惹他们了，于是几人搬回了流云山脚小院，一边给柳成的恢复调养提供一个更好的环境，一边测试这新生的“焚心印”。测试结果令众人啧啧称奇。

首先，焚心印对柳成的指令响应极为顺畅，心念一动即可催发。它能引动附近环境中的天地灵气与五行元素，形成小范围的灵气冲击或元素攻击，虽然威力目前不大，但胜在灵活且消耗极低，仿佛能自行从环境中汲取能量补充。

其次，它似乎对“杂质”有着超乎寻常的喜好。柳成尝试将一些炼废的法器碎片、陈千羽炼丹失败的药渣靠近它，焚心印便会主动散发幽光，将这些“杂质”中残存的混乱灵气或有害物质吸收、提纯，然后释放出相对精纯平和的灵气反馈给柳成或周围环境。而这种灵气，竟然是叶牧谦本人都挑不出毛病的、真正圆满的灵气。

最神奇的是，当柳成或其他人修炼时，将焚心印置于身旁，它能散发一种安宁祥和的气息，有助于平心静气，抵御心魔。它似乎特别喜欢“吃掉”修炼者心中产生的浮躁、焦虑、杂念等负面情绪，让人更容易进入深度入定状态。

这真是好东西！

然而，焚心印也有其脾气：只服从柳成的指令。叶牧谦不知为何竟然也能随意使用，甚至用得比柳成更顺畅。至于白言冰、陈千羽、沈瑶三人，无论他们如何尝试，焚心印都毫无反应，仿佛一块普通的顽石。

\vspace{2em}

柳成炼制焚心印引发的天地异象，范围极广，道韵恢弘，根本瞒不住人。很快，“流云山脉有绝世异宝出世”的消息，以比之前“古传承”流言更迅猛的速度传遍了东域，甚至向着更远的地域扩散。

而随着当日部分在场修士（尤其是寒玉门那些心有不甘的弟子）的私下传播，一个更加夸张、更加引人瞩目的说法迅速成型并如野火般蔓延——“灵枢境小修，炼出天阶法器！”

尽管绝大多数听到这说法的修士第一反应都是嗤之以鼻：“胡说八道！灵枢境和天轮境中间隔了四个大境界，提鞋都不配，还炼制天阶法器？吹牛也不打草稿！”

但是，紧接着补充的细节——“当时天地变色，灵气倒卷，道韵弥漫方圆数十里，寒玉门冰魄真人、天工阁吴毓阁老、郑氏郑季长老皆在场目睹，吴阁老亲口评价‘至少天阶，甚至圣阶’！”——就让许多人将信将疑起来。

尤其是，当日流云山脉上空的天地异象，确实有许多远处修士亲眼目睹，那等声势做不得假。于是，怀疑变成了好奇，好奇变成了探究，探究又滋生出各种各样的心思：羡慕、嫉妒、贪婪、招揽之意……

郑天河在府中听到这个完整消息时，手中的玉杯“啪”地一声捏得粉碎，脸色难看至极。他千算万算，算计着借刀杀人，算计着舆论造势，却唯独没算到，那个他原本只是视为“技术源头”、可以随意拿捏的底层小修士柳成，竟然能搞出如此惊天动地的动静！

“天阶甚至圣阶法器……灵枢境炼制？”郑天河在书房中来回踱步，心绪难平。

他比外人更清楚柳成的底细，也更能判断这消息恐怕并非空穴来风。柳成得了落星谷完整传承，又有叶牧谦这等神秘高人教导，或许真有什么不为人知的秘法或机缘？

随即是巨大的悔恨。早知道柳成有如此潜力，当初就不该用强，应该不惜代价拉拢！现在好了，人得罪死了，还让对方在绝境中爆发出如此光华，名声大噪！此事牵扯到可能的天阶乃至圣器，以及一个能“越阶炼器”的惊世奇才；再加上叶牧谦这个估摸着至少道阶甚至天阶炼器师，这已经不是他郑天河个人能私下处理、积累资本的项目了！

他不敢再有任何隐瞒和拖延，立即前往家族核心区域，求见当代家主，他的父亲——郑伯。他将事情原委（当然，略去了自己暗中指使寒玉门等不光彩细节）以及柳成引发天地异象、疑似炼制成长型高阶法器、吴毓评价等情报，一五一十详细禀报。

郑伯端坐于上首，面容威严，听完郑天河的汇报，沉默良久。他先是对郑天河此前隐瞒不报、试图私吞“留影镜”技术并擅自行动的做法，给予了严厉的批评：“天河，你急功近利，行事欠妥。如此人才与机缘，关乎家族长远，岂能因一己之私而置于险地？若非那叶牧谦确有手段，柳成此子怕是早已遭难，我郑氏不仅错失机缘，更可能平添大敌！”

郑天河冷汗涔涔，躬身认错。

郑伯训斥之后，话锋一转：“你能为家族分忧，这是好的。以后要记住，这种事情最好还是拿到台面上讨论一下。事已至此，那叶牧谦既已通过郑季之口，表达了愿意‘公开部分传承’以平息纷争的态度，倒也算识时务，给了我郑氏一个台阶。”

他沉吟片刻，下令道：“立即以家族名义发布公告：为促进炼器之道交流，弘扬先贤智慧，郑氏将于十日之后，在流云坊市中央广场，举办‘流云论器大会’。届时，将特邀近日获得落星谷传承的叶牧谦先生及其高徒，登台分享炼器心得，展示古法精妙。诚邀东域乃至天下同道，共襄盛举！”

这一招，既顺应了叶牧谦方面“公开传承”的表态，又将主动权抓回郑氏手中，将一场可能引发争夺的危机，转化为一场由郑氏主导、彰显家族气度与影响力的盛会。同时，也能借此机会，近距离观察叶牧谦师徒，尤其是那个柳成，以及那件神秘的“焚心印”。

这种事情，叶牧谦评价为“把我们放在火上烤”，但“又不得不去”。

公告一出，东域震动。无数修士、势力，尤其是炼器师和相关宗门，都将目光投向了流云坊市；一个月前因为新小玄界开启而人满为患的流云坊市，再一次人满为患起来。

天工阁吴毓也加紧了行动。他亲自再次登门，明确提出了拉拢条件：希望聘请叶牧谦为天工阁客卿长老，地位尊崇，资源供应丰厚，且只需在炼器之道上给予一些指点，或合作研究一些项目，无需承担固定职责，自由度极高。这几乎是天工阁能开出的最优惠条件之一。

然而，叶牧谦依旧没有立即接受。他深知树大招风的道理，也明白郑氏举办论器大会背后的复杂意图。他只是对吴毓温和而坚定地说：“吴阁老厚爱，叶某心领。客卿之事，关系重大，请容叶某再思量一二。一切……或许可在论器大会之后，再见分晓。”

吴毓见他态度坚决，也不强求，反而更添几分敬重，留下一些有助于柳成恢复和修炼的珍稀材料作为礼物，便告辞离去。

流云山脚的小院，暂时恢复了表面的宁静。但师徒几人都知道，这宁静之下，暗流愈发汹涌。十日后的论器大会，才是真正的大戏开场。